%============================================================
%  Bd. 17 - Ganze Zahlen
%============================================================

\documentclass[main.tex]{subfiles}

\ifSubfilesClassLoaded{
  \BandSetupStandalone{B17}
}{
}

\title{Bd. 17 - Ganze Zahlen}
\author{Martin Kunze}
\date{}
\setcounter{file}{17}

\hypersetup{
  pdftitle={Bd. 17 - Ganze Zahlen},
  pdfauthor={Martin Kunze}
}

\begin{document}

\title{Bd. 17 - Ganze Zahlen}
\author{Martin Kunze}
\date{}
\hypersetup{pageanchor=false}
\maketitle
\hypersetup{pageanchor=true}
\setcounter{tocdepth}{3}
\tableofcontents
\setcounter{file}{17}
\renewcommand{\FormulaBandID}{17}

\chapter{Einleitung}

Die ganzen Zahlen entstehen aus Paaren natürlicher Zahlen. Das Paar
\((a,b)\) wird als formale Differenz \(a-b\) gelesen. Da verschiedene Paare
dieselbe Differenz beschreiben können, werden genau diejenigen Paare
identifiziert, deren Kreuzsummen übereinstimmen. Die allgemeine Theorie der
Äquivalenzrelationen und Quotienten aus Band~11 liefert anschließend die
Menge der ganzen Zahlen.

Nach der Konstruktion werden Nachfolger, Vorgänger, Ordnung und die
zweiseitige Induktion zunächst vollständig im Quotientenmodell bewiesen.
Die dabei schrittweise gewonnenen Gesetze bilden noch keinen
Abstraktionsschnitt: Bis zum abschließenden Modellnachweis dürfen Klassen und
Repräsentanten weiterhin verwendet werden, jedoch ausschließlich zur
Herleitung der späteren Axiome. Erst im letzten Kapitel wird die Konstruktion
verlassen. Von dort an beschreibt allein die vollständige Axiomatik die
ganzen Zahlen; Addition, Multiplikation, Negation, Ordnung, Nachfolger und
Vorgänger werden dann ohne Rückgriff auf Differenzenpaare oder Klassen
verwendet.

\chapter{Konstruktion der ganzen Zahlen}

\section{Differenzenpaare}

\FormulaDefDeltaK[Träger der Differenzenpaare]{%
  D_{\mathbb Z}\coloneqq\mathbb{N}\times\mathbb{N}%
}{IntegerPairCarrier}{%
  \DeltaRow{Mengen}{D_{\mathbb Z}}%
  \DeltaRow{Neue Symbole}{D_{\mathbb Z}}%
}

Metasprachlich soll ein Paar \((a,b)\) die formale Differenz \(a-b\)
darstellen. Diese Lesart kann hier noch nicht als Definition verwendet
werden, denn eine Subtraktion mit Werten in \(\mathbb N\) steht für
\(a<b\) nicht zur Verfügung. Die anschauliche Gleichheit
\(a-b=c-d\) lässt sich jedoch ohne Subtraktion als
\(a+d=b+c\) schreiben. Die folgende Relation verwendet deshalb nur die
bereits definierte Addition natürlicher Zahlen und identifiziert genau die
Paare, die dieselbe formale Differenz beschreiben sollen.

\FormulaDefDeltaK[Äquivalenz von Differenzenpaaren]{%
  a,b,c,d\in\mathbb{N}\vdash
  (a,b)\IntPairRel(c,d)\coloneqq a+d=b+c%
}{IntPairRelDef}{%
  \DeltaRow{Mengen}{a\dsep b\dsep c\dsep d}%
  \DeltaRow{Relationsoperatoren}{\IntPairRel}%
  \DeltaRow{Trägermenge}{D_{\mathbb Z}}%
  \DeltaRow{Neue Symbole}{\IntPairRel}%
}

\FormulaThmDeltaK[Reflexivität auf Differenzenpaaren]{%
  a,b\in\mathbb{N}\vdash
  (a,b)\IntPairRel(a,b)%
}{IntRelationReflexivePairs}{%
  \DeltaRow{Mengen}{a\dsep b}%
  \DeltaRow{Relationsoperatoren}{\IntPairRel}%
}
\begin{tabproof}
  \proofstep{1}{a,b\in\mathbb{N}}{\rA}
  \proofstep{1}{a+b=b+a}{\FormulaRefAuto{PeanoAddCommutative}{1}}
  \proofstep{1}{(a,b)\IntPairRel(a,b)}{%
    \FormulaRefAuto{IntPairRelDef}[def]{1,2}}
\end{tabproof}

\FormulaThmDeltaK[Symmetrie auf Differenzenpaaren]{%
  a,b,c,d\in\mathbb{N}\dsep
  (a,b)\IntPairRel(c,d)
  \vdash (c,d)\IntPairRel(a,b)%
}{IntRelationSymmetricPairs}{%
  \DeltaRow{Mengen}{a\dsep b\dsep c\dsep d}%
  \DeltaRow{Relationsoperatoren}{\IntPairRel}%
}
\begin{tabproofwide}
  \proofstepwidestar[1]{a,b,c,d\in\mathbb{N}}{\rA}
  \proofstepwidestar[2]{(a,b)\IntPairRel(c,d)}{\rA}
  \proofstepwide[1,2]{a+d}{=}{b+c}{%
    \FormulaRefAuto{IntPairRelDef}[def]{1,2}}
  \proofstepwide[1,2]{c+b}{=}{b+c}{%
    \FormulaRefAuto{PeanoAddCommutative}{1}}
  \proofstepwide[1,2]{}{=}{a+d}{%
    \FormulaRefAuto{a=b\vdash b=a}{3}}
  \proofstepwide[1,2]{}{=}{d+a}{%
    \FormulaRefAuto{PeanoAddCommutative}{1}}
  \proofstepwide[1,2]{c+b}{=}{d+a}{%
    \rChain{4,5,6}}
  \proofstepwidestar[1,2]{(c,d)\IntPairRel(a,b)}{%
    \FormulaRefAuto{IntPairRelDef}[def]{1,7}}
\end{tabproofwide}

\FormulaThmDeltaK[Transitivität auf Differenzenpaaren]{%
  \begin{aligned}[t]
    &a,b,c,d,e,f\in\mathbb{N}\dsep
      (a,b)\IntPairRel(c,d)\dsep{}\\
    &(c,d)\IntPairRel(e,f)
      \vdash (a,b)\IntPairRel(e,f)
  \end{aligned}%
}{IntRelationTransitivePairs}{%
  \DeltaRow{Mengen}{a\dsep b\dsep c\dsep d\dsep e\dsep f}%
  \DeltaRow{Relationsoperatoren}{\IntPairRel}%
}
\begin{tabproofsplitwide}
  \proofpartwideindR[Zusammenführen der Kreuzsummen]{%
    \begin{aligned}[t]
      &a,b,c,d,e,f\in\mathbb{N}\dsep
      a+d=b+c\dsep c+f=d+e\vdash{}\\
      &(a+f)+(c+d)=(b+e)+(c+d)
    \end{aligned}
  }{%
    a,b,c,d,e,f\in\mathbb{N}\dsep
    a+d=b+c\dsep c+f=d+e
    \vdash (a+f)+(c+d)=(b+e)+(c+d)
  }[IntRelationTransitiveRearrangement]
    \proofstepwidestar[1]{a,b,c,d,e,f\in\mathbb{N}}{\rA}
    \proofstepwidestar[2]{a+d=b+c}{\rA}
    \proofstepwidestar[3]{c+f=d+e}{\rA}
    \proofstepwidestar[1,3]{f+c=e+d}{%
      \rChain{%
        \FormulaRefAuto{PeanoAddCommutative}{1},3,%
        \FormulaRefAuto{PeanoAddCommutative}{1}}}
    \proofstepwidestar[1]{d+c=c+d}{%
      \FormulaRefAuto{PeanoAddCommutative}{1}}
    \proofstepwidestar[1,2,3]{%
      (a+f)+(c+d)=(b+e)+(c+d)}{%
      \rIE{%
        \FormulaRefAuto{PeanoSumEqualitiesAddition}{1,2,4},5}}
  \closeproofpartwideindR

  \proofpartwideindR[Kürzen der gemeinsamen Summe]{%
    \begin{aligned}[t]
      &a,b,c,d,e,f\in\mathbb{N}\dsep
      (a+f)+(c+d)=(b+e)+(c+d)\\
      &\vdash a+f=b+e
    \end{aligned}
  }{%
    a,b,c,d,e,f\in\mathbb{N}\dsep
    (a+f)+(c+d)=(b+e)+(c+d)\vdash a+f=b+e
  }[IntRelationTransitiveCancellation]
    \proofstepwidestar[1]{a,b,c,d,e,f\in\mathbb{N}}{\rA}
    \proofstepwidestar[2]{(a+f)+(c+d)=(b+e)+(c+d)}{\rA}
    \proofstepwidestar[1]{a+f\in\mathbb{N}}{\FormulaRefAuto{PeanoAddClosure}{1}}
    \proofstepwidestar[1]{b+e\in\mathbb{N}}{\FormulaRefAuto{PeanoAddClosure}{1}}
    \proofstepwidestar[1]{c+d\in\mathbb{N}}{\FormulaRefAuto{PeanoAddClosure}{1}}
    \proofstepwidestar[1,2]{a+f=b+e}{%
      \FormulaRefAuto{PeanoAddRightCancellation}{3,4,5,2}}
  \closeproofpartwideindR

  \proofpartwide{Schluss}
    \proofstepwidestar[1]{a,b,c,d,e,f\in\mathbb{N}}{\rA}
    \proofstepwidestar[2]{(a,b)\IntPairRel(c,d)}{\rA}
    \proofstepwidestar[3]{(c,d)\IntPairRel(e,f)}{\rA}
    \proofstepwidestar[1,2]{a+d=b+c}{%
      \FormulaRefAuto{IntPairRelDef}[def]{1,2}}
    \proofstepwidestar[1,3]{c+f=d+e}{%
      \FormulaRefAuto{IntPairRelDef}[def]{1,3}}
    \proofstepwidestar[1,2,3]{(a+f)+(c+d)=(b+e)+(c+d)}{%
      \FormulaRefAuto{IntRelationTransitiveRearrangement}{1,4,5}}
    \proofstepwidestar[1,2,3]{a+f=b+e}{%
      \FormulaRefAuto{IntRelationTransitiveCancellation}{1,6}}
    \proofstepwidestar[1,2,3]{(a,b)\IntPairRel(e,f)}{%
      \FormulaRefAuto{IntPairRelDef}[def]{1,7}}
  \closeproofpartwide
\end{tabproofsplitwide}

\FormulaThmDeltaK[Der Äquivalenzoperator auf Differenzenpaaren ist eine Äquivalenzrelation]{%
  \EqRel{D_{\mathbb Z},\IntPairRel}%
}{IntegerPairEquivalence}{%
  \DeltaRow{Mengen}{D_{\mathbb Z}}%
  \DeltaRow{Relationsoperatoren}{\IntPairRel}%
}
\begin{tabproofsplitwide}
  \proofpartwideindR[Reflexivität]{%
    x\in D_{\mathbb Z}\vdash x\IntPairRel x%
  }{%
    x\in D_{\mathbb Z}\vdash x\IntPairRel x%
  }[IntegerPairEquivalenceReflexive]
    \proofstepwidestar[1]{x\in D_{\mathbb Z}}{\rA}
    \proofstepwidestar[1]{\exists a\in\mathbb{N}\,
      \exists b\in\mathbb{N}\,x=(a,b)}{%
      \FormulaRefAuto{x\in A\times B\eqvdash
        \exists a\in A\,\exists b\in B\,x=(a,b)}{1}}
    \proofstepwidestar[3]{a\in\mathbb{N}\land
      b\in\mathbb{N}\land x=(a,b)}{\rA}
    \proofstepwidestar[3]{(a,b)\IntPairRel(a,b)}{%
      \FormulaRefAuto{IntRelationReflexivePairs}{3}}
    \proofstepwidestar[3]{x\IntPairRel x}{\rIE{3,4}}
    \proofstepwidestar[1]{x\IntPairRel x}{\rEE{2,3,5}}
  \closeproofpartwideindR

  \proofpartwideindR[Symmetrie]{%
    x\IntPairRel y\vdash y\IntPairRel x%
  }{%
    x\IntPairRel y\vdash y\IntPairRel x%
  }[IntegerPairEquivalenceSymmetric]
    \proofstepwidestar[1]{x\IntPairRel y}{\rA}
    \proofstepwidestar[1]{y\IntPairRel x}{%
      \FormulaRefAuto{IntRelationSymmetricPairs}[thm]{1}}
  \closeproofpartwideindR

  \proofpartwideindR[Transitivität]{%
    x\IntPairRel y\dsep y\IntPairRel z
    \vdash x\IntPairRel z%
  }{%
    x\IntPairRel y\dsep y\IntPairRel z
    \vdash x\IntPairRel z%
  }[IntegerPairEquivalenceTransitive]
    \proofstepwidestar[1]{x\IntPairRel y}{\rA}
    \proofstepwidestar[2]{y\IntPairRel z}{\rA}
    \proofstepwidestar[1,2]{x\IntPairRel z}{%
      \FormulaRefAuto{IntRelationTransitivePairs}[thm]{1,2}}
  \closeproofpartwideindR

  \proofpartwide{Schluss}
    \proofstepwidestar[]{x\in D_{\mathbb Z}\rightarrow x\IntPairRel x}{%
      \FormulaRefAuto{IntegerPairEquivalenceReflexive}}
    \proofstepwidestar[]{x\IntPairRel y\rightarrow y\IntPairRel x}{%
      \FormulaRefAuto{IntegerPairEquivalenceSymmetric}}
    \proofstepwidestar[3]{x\IntPairRel y}{\rA}
    \proofstepwidestar[4]{y\IntPairRel z}{\rA}
    \proofstepwidestar[3,4]{x\IntPairRel z}{%
      \FormulaRefAuto{IntegerPairEquivalenceTransitive}[thm]{3,4}}
    \proofstepwidestar[]{\EqRel{D_{\mathbb Z},\IntPairRel}}{%
      \FormulaRefAuto{EqRelDef}{1,2,5}}
  \closeproofpartwide
\end{tabproofsplitwide}

\section{Quotient und Klassen}

\FormulaDefDeltaK[Menge der ganzen Zahlen]{%
  \mathbb{Z}\coloneqq
  \QuotientSet{D_{\mathbb Z}}{\IntPairRel}%
}{IntegersAsQuotient}{%
  \DeltaRow{Mengen}{D_{\mathbb Z}\dsep\mathbb Z}%
  \DeltaRow{Relationsoperatoren}{\IntPairRel}%
  \DeltaRow{Neue Symbole}{\mathbb Z}%
}

\FormulaDefDeltaK[Klasse eines Differenzenpaares]{%
  a,b\in\mathbb{N}\vdash
  \IntClass{a}{b}\coloneqq
  \EqClass{(a,b)}{\IntPairRel}{D_{\mathbb Z}}%
}{IntegerClassDef}{%
  \DeltaRow{Mengen}{a\dsep b\dsep D_{\mathbb Z}}%
  \DeltaRow{Relationsoperatoren}{\IntPairRel}%
  \DeltaRow{Neue Symbole}{\IntClass{a}{b}}%
}

\FormulaThmDeltaK[Klassen sind ganze Zahlen]{%
  a,b\in\mathbb{N}\vdash\IntClass{a}{b}\in\mathbb{Z}%
}{IntegerClassInZ}{%
  \DeltaRow{Mengen}{a\dsep b}%
}
\begin{tabproof}
  \proofstep{1}{a,b\in\mathbb{N}}{\rA}
  \proofstep{1}{(a,b)\in D_{\mathbb Z}}{%
    \FormulaRefAuto{(a,b)\in A\times B\eqvdash a\in A\land b\in B}{1}}
  \proofstep{}{\EqRel{D_{\mathbb Z},\IntPairRel}}{%
    \FormulaRefAuto{IntegerPairEquivalence}}
  \proofstep{1}{\EqClass{(a,b)}{\IntPairRel}{D_{\mathbb Z}}
    \in\QuotientSet{D_{\mathbb Z}}{\IntPairRel}}{%
    \FormulaRefAuto{EqClassInQuotient}{3,2}}
  \proofstep{1}{\IntClass{a}{b}=
    \EqClass{(a,b)}{\IntPairRel}{D_{\mathbb Z}}}{%
    \FormulaRefAuto{IntegerClassDef}{1}}
  \proofstep{}{\mathbb Z=
    \QuotientSet{D_{\mathbb Z}}{\IntPairRel}}{%
    \FormulaRefAuto{IntegersAsQuotient}}
  \proofstep{1}{\IntClass{a}{b}\in\mathbb Z}{\rIE{5,6,4}}
\end{tabproof}

\FormulaThmDeltaK[Gleichheit von Ganzzahlklassen]{%
  a,b,c,d\in\mathbb{N}\vdash
  \IntClass{a}{b}=\IntClass{c}{d}
  \leftrightarrow a+d=b+c%
}{IntegerClassEquality}{%
  \DeltaRow{Mengen}{a\dsep b\dsep c\dsep d}%
  \DeltaRow{Relationsoperatoren}{\IntPairRel}%
}
\begin{tabproofsplitwide}
  \proofpartwideindR[Klassenkriterium]{%
    \begin{aligned}[t]
      &a,b,c,d\in\mathbb{N}\vdash{}\\
      &\IntClass{a}{b}=\IntClass{c}{d}
      \leftrightarrow (a,b)\IntPairRel(c,d)
    \end{aligned}
  }{%
    a,b,c,d\in\mathbb{N}\vdash
    \IntClass{a}{b}=\IntClass{c}{d}
    \leftrightarrow (a,b)\IntPairRel(c,d)
  }[IntegerClassEqualityViaRelation]
    \proofstepwidestarbreakkeep[1]{a,b,c,d\in\mathbb{N}}{\rA}
    \proofstepwidestarbreakkeep[]{\EqRel{D_{\mathbb Z},\IntPairRel}}{%
      \FormulaRefAuto{IntegerPairEquivalence}}
    \proofstepwidestarbreakkeep[1]{(a,b)\in D_{\mathbb Z}}{%
      \FormulaRefAuto{(a,b)\in A\times B\eqvdash a\in A\land b\in B}{1}}
    \proofstepwidestarbreakkeep[1]{(c,d)\in D_{\mathbb Z}}{%
      \FormulaRefAuto{(a,b)\in A\times B\eqvdash a\in A\land b\in B}{1}}
    \proofstepwidestarbreakkeep[1]{%
      \EqClass{(a,b)}{\IntPairRel}{D_{\mathbb Z}}
      =\EqClass{(c,d)}{\IntPairRel}{D_{\mathbb Z}}
      \leftrightarrow (a,b)\IntPairRel(c,d)}{%
      \FormulaRefAuto{EqClassEqualIffRel}{2,3,4}}
    \proofstepwidestarbreakkeep[1]{\IntClass{a}{b}=
      \EqClass{(a,b)}{\IntPairRel}{D_{\mathbb Z}}}{%
      \FormulaRefAuto{IntegerClassDef}{1}}
    \proofstepwidestarbreakkeep[1]{\IntClass{c}{d}=
      \EqClass{(c,d)}{\IntPairRel}{D_{\mathbb Z}}}{%
      \FormulaRefAuto{IntegerClassDef}{1}}
    \proofstepwidestarbreakkeep[1]{%
      \IntClass{a}{b}=\IntClass{c}{d}
      \leftrightarrow (a,b)\IntPairRel(c,d)}{\rIE{5,6,7}}
  \closeproofpartwideindR

  \proofpartwide{Schluss}
    \proofstepwidestarbreakkeep[1]{a,b,c,d\in\mathbb{N}}{\rA}
    \proofstepwidestarbreakkeep[1]{%
      \IntClass{a}{b}=\IntClass{c}{d}
      \leftrightarrow (a,b)\IntPairRel(c,d)}{%
      \FormulaRefAuto{IntegerClassEqualityViaRelation}{1}}
    \proofstepwidestarbreakkeep[1]{%
      (a,b)\IntPairRel(c,d)\leftrightarrow a+d=b+c}{%
      \FormulaRefAuto{IntPairRelDef}[def]{1}}
    \proofstepwidestarbreakkeep[1]{%
      \IntClass{a}{b}=\IntClass{c}{d}
      \leftrightarrow a+d=b+c}{%
      \FormulaRefAuto{P\leftrightarrow Q\dsep
        Q\leftrightarrow R\vdash P\leftrightarrow R}{2,3}}
  \closeproofpartwide
\end{tabproofsplitwide}

\FormulaThmDeltaK[Jede ganze Zahl besitzt ein Differenzenpaar]{%
  z\in\mathbb Z\vdash
  \exists a,b\in\mathbb N\,
  z=\IntClass{a}{b}%
}{IntegerPairRepresentative}{%
  \DeltaRow{Mengen}{z\dsep a\dsep b\dsep D_{\mathbb Z}}%
}
\begin{tabproofsplitwide}
  \proofpartwideindR[Quotientenrepräsentant]{%
    z\in\mathbb Z\vdash
    \exists x\in D_{\mathbb Z}\,
    z=\EqClass{x}{\IntPairRel}{D_{\mathbb Z}}%
  }{%
    z\in\mathbb Z\vdash
    \exists x\in D_{\mathbb Z}\,
    z=\EqClass{x}{\IntPairRel}{D_{\mathbb Z}}%
  }[IntegerQuotientRepresentative]
    \proofstepwidestar[1]{z\in\mathbb Z}{\rA}
    \proofstepwidestar[]{\mathbb Z=
      \QuotientSet{D_{\mathbb Z}}{\IntPairRel}}{%
      \FormulaRefAuto{IntegersAsQuotient}}
    \proofstepwidestar[1]{z\in
      \QuotientSet{D_{\mathbb Z}}{\IntPairRel}}{\rIE{2,1}}
    \proofstepwidestar[]{\EqRel{D_{\mathbb Z},\IntPairRel}}{%
      \FormulaRefAuto{IntegerPairEquivalence}}
    \proofstepwidestar[1]{\exists x\in D_{\mathbb Z}\,
      z=\EqClass{x}{\IntPairRel}{D_{\mathbb Z}}}{%
      \FormulaRefAuto{QuotientRepresentative}{4,3}}
  \closeproofpartwideindR

  \proofpartwide{Schluss}
    \proofstepwidestar[1]{z\in\mathbb Z}{\rA}
    \proofstepwidestar[1]{\exists x\in D_{\mathbb Z}\,
      z=\EqClass{x}{\IntPairRel}{D_{\mathbb Z}}}{%
      \FormulaRefAuto{IntegerQuotientRepresentative}{1}}
    \proofstepwidestar[3]{x\in D_{\mathbb Z}\land
      z=\EqClass{x}{\IntPairRel}{D_{\mathbb Z}}}{\rA}
    \proofstepwidestar[3]{\exists a,b\in\mathbb N\,x=(a,b)}{%
      \FormulaRefAuto{x\in A\times B\eqvdash
        \exists a\in A\,\exists b\in B\,x=(a,b)}{\rAEa{3}}}
    \proofstepwidestar[5]{a,b\in\mathbb N\land x=(a,b)}{\rA}
    \proofstepwidestar[3,5]{z=\EqClass{(a,b)}{\IntPairRel}{D_{\mathbb Z}}}{%
      \rIE{5,\rAEb{3}}}
    \proofstepwidestar[5]{\IntClass{a}{b}=
      \EqClass{(a,b)}{\IntPairRel}{D_{\mathbb Z}}}{%
      \FormulaRefAuto{IntegerClassDef}{5}}
    \proofstepwidestar[3,5]{z=\IntClass{a}{b}}{%
      \FormulaRefAuto{a=b\dsep c=b\vdash a=c}{6,7}}
    \proofstepwidestar[3]{\exists a,b\in\mathbb N\,
      z=\IntClass{a}{b}}{\rEE{4,5,8}}
    \proofstepwidestar[1]{\exists a,b\in\mathbb N\,
      z=\IntClass{a}{b}}{\rEE{2,3,9}}
  \closeproofpartwide
\end{tabproofsplitwide}

\section{Einbettung der natürlichen Zahlen}

\FormulaDefDeltaK[Natürliche Einbettung]{%
  \begin{aligned}[t]
    &\IntEmbed\colon\mathbb N\to\mathbb Z,\\[-2pt]
    &\forall n\in\mathbb N\,
      \bigl(\IntEmbed(n)\coloneqq\IntClass{n}{0}\bigr)
  \end{aligned}%
}{NaturalIntegerEmbedding}{%
  \DeltaRow{Mengen}{n}%
  \DeltaRow{Funktionensymbol}{\IntEmbed}%
  \DeltaRow{Neue Symbole}{\IntEmbed}%
}
\begin{tabproof}
  \proofstep{1}{n\in\mathbb N}{\rA}
  \proofstep{}{0\in\mathbb N}{\FormulaRefAuto{PeanoZero}[ax]}
  \proofstep{1}{\IntClass{n}{0}\in\mathbb Z}{%
    \FormulaRefAuto{IntegerClassInZ}[thm]{1,2}}
  \proofstep{}{\forall n\in\mathbb N\,\bigl(\IntClass{n}{0}\in\mathbb Z\bigr)}{%
    \rUI{\rRI{1,3}}}
\end{tabproof}

\FormulaThmDeltaK[Typisierung der natürlichen Einbettung]{%
  \IntEmbed\colon\mathbb N\to\mathbb Z%
}{NaturalIntegerEmbeddingFunction}{%
  \DeltaRow{Funktion}{\IntEmbed}%
}
\begin{tabproof}
  \proofstep{}{\IntEmbed\colon\mathbb N\to\mathbb Z}{%
    \FormulaRefAuto{NaturalIntegerEmbedding}[def]}
\end{tabproof}

\FormulaThmDeltaK[Grundgleichung der natürlichen Einbettung]{%
  n\in\mathbb N\vdash \IntEmbed(n)=\IntClass{n}{0}%
}{NaturalIntegerEmbeddingValue}{%
  \DeltaRow{Mengen}{n}%
  \DeltaRow{Funktion}{\IntEmbed}%
}
\begin{tabproof}
  \proofstep{1}{n\in\mathbb N}{\rA}
  \proofstep{}{%
    \forall u\in\mathbb N\,
      \bigl(\IntEmbed(u)=\IntClass{u}{0}\bigr)}{%
    \FormulaRefAuto{NaturalIntegerEmbedding}[def]}
  \proofstep{1}{\IntEmbed(n)=\IntClass{n}{0}}{%
    \rRE{\rUE{2},1}}
\end{tabproof}

\FormulaThmDeltaK[Injektivität der natürlichen Einbettung]{%
  \IntEmbed\colon\mathbb N\inj\mathbb Z%
}{NaturalIntegerEmbeddingInjective}{%
  \DeltaRow{Mengen}{m\dsep n}%
  \DeltaRow{Funktion}{\IntEmbed}%
}
\begin{tabproofsplitwide}
  \proofpartwideindR[Gleichheit von eingebetteten Zahlen]{%
    m,n\in\mathbb N\dsep
    \IntEmbed(m)=\IntEmbed(n)\vdash m=n%
  }{%
    m,n\in\mathbb N\dsep
    \IntEmbed(m)=\IntEmbed(n)\vdash m=n%
  }[NaturalIntegerEmbeddingReflectsEquality]
    \proofstepwidestar[1]{m,n\in\mathbb N}{\rA}
    \proofstepwidestar[2]{\IntEmbed(m)=\IntEmbed(n)}{\rA}
    \proofstepwidestar[1]{\IntEmbed(m)=\IntClass{m}{0}}{%
      \FormulaRefAuto{NaturalIntegerEmbeddingValue}[thm]{1}}
    \proofstepwidestar[1]{\IntEmbed(n)=\IntClass{n}{0}}{%
      \FormulaRefAuto{NaturalIntegerEmbeddingValue}[thm]{1}}
    \proofstepwidestar[1,2]{\IntClass{m}{0}=\IntClass{n}{0}}{\rIE{3,4,2}}
    \proofstepwidestar[1,2]{m+0=0+n}{%
      \FormulaRefAuto{IntegerClassEquality}{1,5}}
    \proofstepwide[1]{m}{=}{m+0}{%
      \FormulaRefAuto{a=b\vdash b=a}{%
        \FormulaRefAuto{PeanoAddRightZero}{1}}}
    \proofstepwide[1,2]{}{=}{0+n}{\rIE{6}}
    \proofstepwide[1]{}{=}{n}{\FormulaRefAuto{PeanoAddLeftZero}{1}}
    \proofstepwidestar[1,2]{m=n}{\rChain{7,8,9}}
  \closeproofpartwideindR

  \proofpartwide{Schluss}
    \proofstepwidestar[]{\IntEmbed\colon\mathbb N\to\mathbb Z}{%
      \FormulaRefAuto{NaturalIntegerEmbeddingFunction}}
    \proofstepwidestar[2]{m\in\mathbb N}{\rA}
    \proofstepwidestar[3]{n\in\mathbb N}{\rA}
    \proofstepwidestar[4]{\IntEmbed(m)=\IntEmbed(n)}{\rA}
    \proofstepwidestar[2,3,4]{m=n}{%
      \FormulaRefAuto{NaturalIntegerEmbeddingReflectsEquality}{2,3,4}}
    \proofstepwidestar[]{\IntEmbed\colon\mathbb N\inj\mathbb Z}{%
      \FormulaRefAuto{Injektive Funktion}{1,\rUI{\rRI{2,3,4,5}}}}
  \closeproofpartwide
\end{tabproofsplitwide}

\subsection{Kanonische Teilmengenkopie}

Der Quotiententräger \(\mathbb Z\) enthält die natürlichen Zahlen in der
vorliegenden Konstruktion nicht wörtlich.  Das Bild der natürlichen
Einbettung liefert jedoch eine kanonische Kopie, die tatsächlich eine
Teilmenge von \(\mathbb Z\) ist.

\FormulaDefDeltaK[Kanonische natürliche Kopie in den ganzen Zahlen]{%
  \NatCopyInInt\coloneqq\IntEmbed[\mathbb N]%
}{NaturalCopyInIntegers}{%
  \DeltaRow{Mengen}{\NatCopyInInt}%
  \DeltaRow{Funktion}{\IntEmbed}%
  \DeltaRow{Neue Symbole}{\NatCopyInInt}%
}

\FormulaThmDeltaK[Die natürliche Kopie ist eine Teilmenge der ganzen Zahlen]{%
  \NatCopyInInt\subseteq\mathbb Z%
}{NaturalCopySubsetIntegers}{%
  \DeltaRow{Mengen}{\NatCopyInInt}%
  \DeltaRow{Funktion}{\IntEmbed}%
}
\begin{tabproof}
  \proofstep{}{\IntEmbed\colon\mathbb N\to\mathbb Z}{%
    \FormulaRefAuto{NaturalIntegerEmbeddingFunction}[thm]}
  \proofstep{}{\IntEmbed[\mathbb N]\subseteq\mathbb Z}{%
    \FormulaRefAuto{F\colon A\to B\vdash F[A]\subseteq B}[thm]{1}}
  \proofstep{}{\NatCopyInInt=\IntEmbed[\mathbb N]}{%
    \FormulaRefAuto{NaturalCopyInIntegers}[def]}
  \proofstep{}{\NatCopyInInt\subseteq\mathbb Z}{\rIE{3,2}}
\end{tabproof}

\FormulaThmDeltaK[Bijektive Identifikation der natürlichen Kopie]{%
  \IntEmbed\colon\mathbb N\bij\NatCopyInInt%
}{NaturalCopyIntegerBijection}{%
  \DeltaRow{Mengen}{\NatCopyInInt}%
  \DeltaRow{Funktion}{\IntEmbed}%
}
\begin{tabproof}
  \proofstep{}{\IntEmbed\colon\mathbb N\inj\mathbb Z}{%
    \FormulaRefAuto{NaturalIntegerEmbeddingInjective}[thm]}
  \proofstep{}{\IntEmbed\colon\mathbb N\bij\IntEmbed[\mathbb N]}{%
    \FormulaRefAuto{F\colon A\inj B\vdash F\colon A\bij F[A]}[thm]{1}}
  \proofstep{}{\NatCopyInInt=\IntEmbed[\mathbb N]}{%
    \FormulaRefAuto{NaturalCopyInIntegers}[def]}
  \proofstep{}{\IntEmbed\colon\mathbb N\bij\NatCopyInInt}{\rIE{3,2}}
\end{tabproof}

Im ganzzahligen Kontext verwenden wir von nun an dieselben sichtbaren
Zahlzeichen wie bei den natürlichen Zahlen.  Semantisch bezeichnet
\(\IntNumeral{n}\) für \(n\in\mathbb N\) stets das Bild \(\IntOfNat{n}\).
Die sichtbaren Zeichen allein behaupten keine Gleichheit der ursprünglichen
Träger.  Die wörtliche Inklusion wird vielmehr durch
\(\NatCopyInInt\subseteq\mathbb Z\) ausgedrückt; die vorstehende Bijektion
rechtfertigt die übliche Identifikation mit \(\mathbb N\).  Wo der umgebende
Träger nicht eindeutig ist, bleibt die Einbettung mit \(\IntOfNat{n}\)
ausdrücklich sichtbar.
Insbesondere stehen \(\IntZero,\IntOne,\IntNumeral{2},\IntNumeral{3},
\IntNumeral{4},\ldots\) im ganzzahligen Kontext ohne Trägerindex.

\FormulaDefDeltaK[Unindizierte ganzzahlige Zahlzeichen]{%
  n\in\mathbb N\vdash
  \IntNumeral{n}\coloneqq\IntOfNat{n}%
}{IntegerNumeralConvention}{%
  \DeltaRow{Mengen}{n}%
  \DeltaRow{Funktionensymbol}{\IntEmbed}%
  \DeltaRow{Neue Schreibweise}{\IntNumeral{n}}%
}

\FormulaDefDeltaK[Null der ganzen Zahlen]{%
  \IntZero\coloneqq\IntOfNat{0}%
}{IntegerZeroDef}{%
  \DeltaRow{Mengen}{\IntZero}%
  \DeltaRow{Neue Symbole}{\IntZero}%
}

\FormulaDefDeltaK[Eins der ganzen Zahlen]{%
  \IntOne\coloneqq\IntOfNat{1}%
}{IntegerOneDef}{%
  \DeltaRow{Mengen}{\IntOne}%
  \DeltaRow{Neue Symbole}{\IntOne}%
}

\FormulaThmDeltaK[Nullklasse]{%
  \IntZero=\IntClass{0}{0}%
}{IntegerZeroClass}{%
  \DeltaRow{Mengen}{\IntZero}%
}
\begin{tabproof}
  \proofstep{}{\IntZero=\IntEmbed(0)}{\FormulaRefAuto{IntegerZeroDef}}
  \proofstep{}{0\in\mathbb N}{\FormulaRefAuto{PeanoZero}[ax]}
  \proofstep{}{\IntEmbed(0)=\IntClass{0}{0}}{%
    \FormulaRefAuto{NaturalIntegerEmbeddingValue}[thm]{2}}
  \proofstep{}{\IntZero=\IntClass{0}{0}}{%
    \rChain{1,3}}
\end{tabproof}

\FormulaThmDeltaK[Die Ganzzahlnull liegt im Quotienten]{%
  \IntZero\in\mathbb Z%
}{IntZeroCarrier}{%
  \DeltaRow{Mengen}{\IntZero}%
}
\begin{tabproof}
  \proofstep{}{\IntZero=\IntClass{0}{0}}{%
    \FormulaRefAuto{IntegerZeroClass}}
  \proofstep{}{0\in\mathbb N}{\FormulaRefAuto{PeanoZero}[ax]}
  \proofstep{}{\IntClass{0}{0}\in\mathbb Z}{%
    \FormulaRefAuto{IntegerClassInZ}{2,2}}
  \proofstep{}{\IntZero\in\mathbb Z}{\rIE{1,3}}
\end{tabproof}

\FormulaThmDeltaK[Die Ganzzahleins liegt im Quotienten]{%
  \IntOne\in\mathbb Z%
}{IntOneCarrier}{%
  \DeltaRow{Mengen}{\IntOne}%
}
\begin{tabproof}
  \proofstep{}{\IntOne=\IntEmbed(1)}{%
    \FormulaRefAuto{IntegerOneDef}}
  \proofstep{}{1\in\mathbb N}{\FormulaRefAuto{PeanoOneInN}}
  \proofstep{}{\IntEmbed\colon\mathbb N\to\mathbb Z}{%
    \FormulaRefAuto{NaturalIntegerEmbeddingFunction}}
  \proofstep{}{\IntEmbed(1)\in\mathbb Z}{%
    \FormulaRefAuto{F\colon A\to B\dsep x\in A\vdash F(x)\in B}{3,2}}
  \proofstep{}{\IntOne\in\mathbb Z}{\rIE{1,4}}
\end{tabproof}

\section{Negation und Addition}

\subsection{Negation}

\FormulaThmDeltaK[Repräsentantenunabhängigkeit der Negation]{%
  \begin{aligned}[t]
    &a,b,c,d\in\mathbb N\dsep
      \IntClass{a}{b}=\IntClass{c}{d}\vdash{}\\
    &\IntClass{b}{a}=\IntClass{d}{c}
  \end{aligned}%
}{IntegerNegationCompatible}{%
  \DeltaRow{Mengen}{a\dsep b\dsep c\dsep d}%
}
\begin{tabproofsplitwide}
  \proofpartwideindR[Kreuzsummen des Ausgangspaars]{%
    a,b,c,d\in\mathbb N\dsep
    \IntClass{a}{b}=\IntClass{c}{d}
    \vdash a+d=b+c%
  }{%
    a,b,c,d\in\mathbb N\dsep
    \IntClass{a}{b}=\IntClass{c}{d}
    \vdash a+d=b+c%
  }[IntegerNegationCompatibleCrossSum]
    \proofstepwidestar[1]{a,b,c,d\in\mathbb N}{\rA}
    \proofstepwidestar[2]{\IntClass{a}{b}=\IntClass{c}{d}}{\rA}
    \proofstepwidestar[1,2]{a+d=b+c}{%
      \FormulaRefAuto{IntegerClassEquality}{1,2}}
  \closeproofpartwideindR

  \proofpartwide{Vertauschen der Koordinaten}
    \proofstepwidestar[1]{a,b,c,d\in\mathbb N}{\rA}
    \proofstepwidestar[2]{\IntClass{a}{b}=\IntClass{c}{d}}{\rA}
    \proofstepwidestar[1,2]{a+d=b+c}{%
      \FormulaRefAuto{IntegerNegationCompatibleCrossSum}{1,2}}
    \proofstepwidestar[1,2]{b+c=a+d}{%
      \FormulaRefAuto{a=b\vdash b=a}{3}}
    \proofstepwidestar[1,2]{\IntClass{b}{a}=\IntClass{d}{c}}{%
      \FormulaRefAuto{IntegerClassEquality}{1,4}}
  \closeproofpartwide
\end{tabproofsplitwide}

\FormulaThmDeltaK[Quotientenabstieg der Negation]{%
  \begin{aligned}[t]
    z\in\mathbb Z\vdash
    \exists! u\in\mathbb Z\,
    \exists p\in\mathbb N\,\exists q\in\mathbb N\,
    \bigl(z=\IntClass{p}{q}\land u=\IntClass{q}{p}\bigr)
  \end{aligned}%
}{IntegerNegationQuotientDescent}{%
  \DeltaRow{Mengen}{z\dsep u\dsep v\dsep
    a\dsep b\dsep c\dsep d\dsep p\dsep q}%
}
\begin{tabproofsplitwide}
  \proofpartwideindR[Existenz aus einem Repräsentanten]{%
    \begin{aligned}[t]
      z\in\mathbb Z\vdash
      \exists u\in\mathbb Z\,
      \exists p\in\mathbb N\,\exists q\in\mathbb N\,
      \bigl(z=\IntClass{p}{q}\land u=\IntClass{q}{p}\bigr)
    \end{aligned}%
  }{%
    z\in\mathbb Z\vdash
    \exists u\in\mathbb Z\,
    \exists p\in\mathbb N\,\exists q\in\mathbb N\,
    \bigl(z=\IntClass{p}{q}\land u=\IntClass{q}{p}\bigr)%
  }[IntegerNegationDescentExistence]
    \proofstepwidestar[1]{z\in\mathbb Z}{\rA}
    \proofstepwidestar[1]{%
      \exists a\in\mathbb N\,\exists b\in\mathbb N\,
      z=\IntClass{a}{b}}{%
      \FormulaRefAuto{IntegerPairRepresentative}{1}}
    \proofstepwidestar[3]{%
      a,b\in\mathbb N\land z=\IntClass{a}{b}}{\rA}
    \proofstepwidestar[3]{\IntClass{b}{a}\in\mathbb Z}{%
      \FormulaRefAuto{IntegerClassInZ}{3}}
    \proofstepwidestar[]{%
      \IntClass{b}{a}=\IntClass{b}{a}}{\rII}
    \proofstepwidestar[3]{%
      \exists p\in\mathbb N\,\exists q\in\mathbb N\,
      \bigl(z=\IntClass{p}{q}\land
      \IntClass{b}{a}=\IntClass{q}{p}\bigr)}{%
      \rEI{\rAI{3,5}}}
    \proofstepwidestar[3]{%
      \IntClass{b}{a}\in\mathbb Z\land
      \exists p\in\mathbb N\,\exists q\in\mathbb N\,
      \bigl(z=\IntClass{p}{q}\land
      \IntClass{b}{a}=\IntClass{q}{p}\bigr)}{%
      \rAI{4,6}}
    \proofstepwidestar[3]{%
      \exists u\in\mathbb Z\,
      \exists p\in\mathbb N\,\exists q\in\mathbb N\,
      \bigl(z=\IntClass{p}{q}\land u=\IntClass{q}{p}\bigr)}{%
      \rEI{7}}
    \proofstepwidestar[1]{%
      \exists u\in\mathbb Z\,
      \exists p\in\mathbb N\,\exists q\in\mathbb N\,
      \bigl(z=\IntClass{p}{q}\land u=\IntClass{q}{p}\bigr)}{%
      \rEE{2,3,8}}
  \closeproofpartwideindR

  \proofpartwideindR[Eindeutigkeit für feste Repräsentanten]{%
    \begin{aligned}[t]
      &a,b,c,d\in\mathbb N\dsep
      z=\IntClass{a}{b}\dsep u=\IntClass{b}{a}\dsep{}\\
      &z=\IntClass{c}{d}\dsep v=\IntClass{d}{c}
      \vdash u=v
    \end{aligned}%
  }{%
    a,b,c,d\in\mathbb N\dsep
    z=\IntClass{a}{b}\dsep u=\IntClass{b}{a}\dsep
    z=\IntClass{c}{d}\dsep v=\IntClass{d}{c}
    \vdash u=v%
  }[IntegerNegationDescentFixedUniqueness]
    \proofstepwidestar[1]{a,b,c,d\in\mathbb N}{\rA}
    \proofstepwidestar[2]{z=\IntClass{a}{b}}{\rA}
    \proofstepwidestar[3]{u=\IntClass{b}{a}}{\rA}
    \proofstepwidestar[4]{z=\IntClass{c}{d}}{\rA}
    \proofstepwidestar[5]{v=\IntClass{d}{c}}{\rA}
    \proofstepwidestar[2,4]{%
      \IntClass{a}{b}=\IntClass{c}{d}}{\rIE{2,4}}
    \proofstepwidestar[1,2,4]{%
      \IntClass{b}{a}=\IntClass{d}{c}}{%
      \FormulaRefAuto{IntegerNegationCompatible}{1,6}}
    \proofstepwidestar[1,2,3,4,5]{u=v}{\rIE{3,5,7}}
  \closeproofpartwideindR

  \proofpartwideindR[Eindeutigkeit beliebiger Ergebniszeugen]{%
    \begin{aligned}[t]
      &\bigl(u\in\mathbb Z\land
      \exists a\in\mathbb N\,\exists b\in\mathbb N\,
        (z=\IntClass{a}{b}\land u=\IntClass{b}{a})\bigr)\dsep{}\\
      &\bigl(v\in\mathbb Z\land
      \exists c\in\mathbb N\,\exists d\in\mathbb N\,
        (z=\IntClass{c}{d}\land v=\IntClass{d}{c})\bigr)
      \vdash u=v
    \end{aligned}%
  }{%
    \bigl(u\in\mathbb Z\land
    \exists a\in\mathbb N\,\exists b\in\mathbb N\,
      (z=\IntClass{a}{b}\land u=\IntClass{b}{a})\bigr)\dsep
    \bigl(v\in\mathbb Z\land
    \exists c\in\mathbb N\,\exists d\in\mathbb N\,
      (z=\IntClass{c}{d}\land v=\IntClass{d}{c})\bigr)
    \vdash u=v%
  }[IntegerNegationDescentUniqueness]
    \proofstepwidestar[1]{%
      u\in\mathbb Z\land
      \exists a\in\mathbb N\,\exists b\in\mathbb N\,
      (z=\IntClass{a}{b}\land u=\IntClass{b}{a})}{\rA}
    \proofstepwidestar[2]{%
      v\in\mathbb Z\land
      \exists c\in\mathbb N\,\exists d\in\mathbb N\,
      (z=\IntClass{c}{d}\land v=\IntClass{d}{c})}{\rA}
    \proofstepwidestar[1]{%
      \exists a\in\mathbb N\,\exists b\in\mathbb N\,
      (z=\IntClass{a}{b}\land u=\IntClass{b}{a})}{\rAEb{1}}
    \proofstepwidestar[2]{%
      \exists c\in\mathbb N\,\exists d\in\mathbb N\,
      (z=\IntClass{c}{d}\land v=\IntClass{d}{c})}{\rAEb{2}}
    \proofstepwidestar[5]{%
      a,b\in\mathbb N\land
      z=\IntClass{a}{b}\land u=\IntClass{b}{a}}{\rA}
    \proofstepwidestar[6]{%
      c,d\in\mathbb N\land
      z=\IntClass{c}{d}\land v=\IntClass{d}{c}}{\rA}
    \proofstepwidestar[5]{a,b\in\mathbb N}{\rAEa{5}}
    \proofstepwidestar[5]{z=\IntClass{a}{b}}{\rAEa{\rAEb{5}}}
    \proofstepwidestar[5]{u=\IntClass{b}{a}}{\rAEb{\rAEb{5}}}
    \proofstepwidestar[6]{c,d\in\mathbb N}{\rAEa{6}}
    \proofstepwidestar[6]{z=\IntClass{c}{d}}{\rAEa{\rAEb{6}}}
    \proofstepwidestar[6]{v=\IntClass{d}{c}}{\rAEb{\rAEb{6}}}
    \proofstepwidestar[5,6]{a,b,c,d\in\mathbb N}{\rAI{7,10}}
    \proofstepwidestar[5,6]{u=v}{%
      \FormulaRefAuto{IntegerNegationDescentFixedUniqueness}{%
        13,8,9,11,12}}
    \proofstepwidestar[5]{u=v}{\rEE{4,6,14}}
    \proofstepwidestar[1,2]{u=v}{\rEE{3,5,15}}
  \closeproofpartwideindR

  \proofpartwideindR[Abschluss des Abstiegs]{%
    \begin{aligned}[t]
      z\in\mathbb Z\vdash
      \exists! u\in\mathbb Z\,
      \exists p\in\mathbb N\,\exists q\in\mathbb N\,
      \bigl(z=\IntClass{p}{q}\land u=\IntClass{q}{p}\bigr)
    \end{aligned}%
  }{%
    z\in\mathbb Z\vdash
    \exists! u\in\mathbb Z\,
    \exists p\in\mathbb N\,\exists q\in\mathbb N\,
    \bigl(z=\IntClass{p}{q}\land u=\IntClass{q}{p}\bigr)%
  }[IntegerNegationDescentConclusion]
    \proofstepwidestar[1]{z\in\mathbb Z}{\rA}
    \proofstepwidestar[1]{%
      \exists u\in\mathbb Z\,
      \exists p\in\mathbb N\,\exists q\in\mathbb N\,
      (z=\IntClass{p}{q}\land u=\IntClass{q}{p})}{%
      \FormulaRefAuto{IntegerNegationDescentExistence}{1}}
    \proofstepwidestar[3]{%
      u\in\mathbb Z\land
      \exists p\in\mathbb N\,\exists q\in\mathbb N\,
      (z=\IntClass{p}{q}\land u=\IntClass{q}{p})}{\rA}
    \proofstepwidestar[4]{%
      v\in\mathbb Z\land
      \exists p\in\mathbb N\,\exists q\in\mathbb N\,
      (z=\IntClass{p}{q}\land v=\IntClass{q}{p})}{\rA}
    \proofstepwidestar[3,4]{u=v}{%
      \FormulaRefAuto{IntegerNegationDescentUniqueness}{3,4}}
    \proofstepwidestar[1]{%
      \exists! u\in\mathbb Z\,
      \exists p\in\mathbb N\,\exists q\in\mathbb N\,
      (z=\IntClass{p}{q}\land u=\IntClass{q}{p})}{%
      \UEI{2,3,4,5}}
  \closeproofpartwideindR
\end{tabproofsplitwide}

Der vorangehende Satz zeigt, dass die folgende Klassenregel unabhängig
vom gewählten Differenzenpaar genau ein Ergebnis in \(\mathbb Z\) festlegt.

\FormulaDefDeltaK[Negation ganzer Zahlen]{%
  a,b\in\mathbb N\vdash
  -\IntClass{a}{b}\coloneqq\IntClass{b}{a}%
}{IntegerNegationDef}{%
  \DeltaRow{Mengen}{a\dsep b}%
  \DeltaRow{Neue Symbole}{-}%
}

\FormulaThmDeltaK[Abgeschlossenheit der Negation]{%
  z\in\mathbb Z\vdash -z\in\mathbb Z%
}{IntNegClosure}{%
  \DeltaRow{Mengen}{z\dsep a\dsep b}%
}
\begin{tabproof}
  \proofstep{1}{z\in\mathbb Z}{\rA}
  \proofstep{1}{\exists a\in\mathbb N\,\exists b\in\mathbb N\,
    z=\IntClass{a}{b}}{\FormulaRefAuto{IntegerPairRepresentative}{1}}
  \proofstep{3}{a,b\in\mathbb N\land z=\IntClass{a}{b}}{\rA}
  \proofstep{3}{-\IntClass{a}{b}=\IntClass{b}{a}}{%
    \FormulaRefAuto{IntegerNegationDef}{3}}
  \proofstep{3}{\IntClass{b}{a}\in\mathbb Z}{%
    \FormulaRefAuto{IntegerClassInZ}{3}}
  \proofstep{3}{-z\in\mathbb Z}{\rIE{3,4,5}}
  \proofstep{1}{-z\in\mathbb Z}{\rEE{2,3,6}}
\end{tabproof}

\FormulaThmDeltaK[Doppelte Negation]{%
  z\in\mathbb Z\vdash -(-z)=z%
}{IntegerDoubleNegation}{%
  \DeltaRow{Mengen}{z\dsep a\dsep b}%
}
\begin{tabproof}
  \proofstep{1}{z\in\mathbb Z}{\rA}
  \proofstep{1}{\exists a,b\in\mathbb N\,
    z=\IntClass{a}{b}}{\FormulaRefAuto{IntegerPairRepresentative}{1}}
  \proofstep{3}{a,b\in\mathbb N\land z=\IntClass{a}{b}}{\rA}
  \proofstep{3}{-\IntClass{a}{b}=\IntClass{b}{a}}{%
    \FormulaRefAuto{IntegerNegationDef}{3}}
  \proofstep{3}{-\IntClass{b}{a}=\IntClass{a}{b}}{%
    \FormulaRefAuto{IntegerNegationDef}{3}}
  \proofstep{3}{-(-z)=z}{\rIE{3,4,5}}
  \proofstep{1}{-(-z)=z}{\rEE{2,3,6}}
\end{tabproof}

\subsection{Addition}

\FormulaThmDeltaK[Repräsentantenunabhängigkeit der Addition]{%
  \begin{aligned}[t]
    &a,b,a',b',c,d,c',d'\in\mathbb N\dsep
      \IntClass{a}{b}=\IntClass{a'}{b'}\dsep
      \IntClass{c}{d}=\IntClass{c'}{d'}\vdash{}\\
    &\IntClass{a+c}{b+d}=\IntClass{a'+c'}{b'+d'}
  \end{aligned}%
}{IntegerAdditionCompatible}{%
  \DeltaRow{Mengen}{a\dsep b\dsep a'\dsep b'\dsep
    c\dsep d\dsep c'\dsep d'}%
}
\begin{tabproofsplitwide}
  \proofpartwideindR[Kreuzsummen der beiden Eingaben]{%
    \begin{aligned}[t]
      &a,b,a',b',c,d,c',d'\in\mathbb N\dsep
      \IntClass{a}{b}=\IntClass{a'}{b'}\dsep
      \IntClass{c}{d}=\IntClass{c'}{d'}\vdash{}\\
      &a+b'=b+a'\quad\land\quad c+d'=d+c'
    \end{aligned}
  }{%
    a,b,a',b',c,d,c',d'\in\mathbb N\dsep
    \IntClass{a}{b}=\IntClass{a'}{b'}\dsep
    \IntClass{c}{d}=\IntClass{c'}{d'}
    \vdash a+b'=b+a'\land c+d'=d+c'
  }[IntegerAdditionInputCrossSums]
    \proofstepwidestar[1]{a,b,a',b',c,d,c',d'\in\mathbb N}{\rA}
    \proofstepwidestar[2]{\IntClass{a}{b}=\IntClass{a'}{b'}}{\rA}
    \proofstepwidestar[3]{\IntClass{c}{d}=\IntClass{c'}{d'}}{\rA}
    \proofstepwidestar[1,2]{a+b'=b+a'}{%
      \FormulaRefAuto{IntegerClassEquality}{1,2}}
    \proofstepwidestar[1,3]{c+d'=d+c'}{%
      \FormulaRefAuto{IntegerClassEquality}{1,3}}
    \proofstepwidestar[1,2,3]{a+b'=b+a'\land c+d'=d+c'}{\rAI{4,5}}
  \closeproofpartwideindR

  \proofpartwideindR[Zusammenführen und Umsortieren]{%
    \begin{aligned}[t]
      &a,b,a',b',c,d,c',d'\in\mathbb N\dsep
      a+b'=b+a'\dsep c+d'=d+c'\vdash{}\\
      &(a+c)+(b'+d')=(b+d)+(a'+c')
    \end{aligned}
  }{%
    a,b,a',b',c,d,c',d'\in\mathbb N\dsep
    a+b'=b+a'\dsep c+d'=d+c'
    \vdash (a+c)+(b'+d')=(b+d)+(a'+c')
  }[IntegerAdditionOutputCrossSum]
    \proofstepwidestar[1]{a,b,a',b',c,d,c',d'\in\mathbb N}{\rA}
    \proofstepwidestar[2]{a+b'=b+a'}{\rA}
    \proofstepwidestar[3]{c+d'=d+c'}{\rA}
    \proofstepwide[1,2,3]{(a+c)+(b'+d')}{=}{(b+d)+(a'+c')}{%
      \FormulaRefAuto{PeanoSumEqualitiesAddition}{1,2,3}}
  \closeproofpartwideindR

  \proofpartwide{Schluss}
    \proofstepwidestar[1]{a,b,a',b',c,d,c',d'\in\mathbb N}{\rA}
    \proofstepwidestar[2]{\IntClass{a}{b}=\IntClass{a'}{b'}}{\rA}
    \proofstepwidestar[3]{\IntClass{c}{d}=\IntClass{c'}{d'}}{\rA}
    \proofstepwidestar[1,2,3]{a+b'=b+a'\land c+d'=d+c'}{%
      \FormulaRefAuto{IntegerAdditionInputCrossSums}{1,2,3}}
    \proofstepwidestar[1,2,3]{(a+c)+(b'+d')=(b+d)+(a'+c')}{%
      \FormulaRefAuto{IntegerAdditionOutputCrossSum}{1,\rAEa{4},\rAEb{4}}}
    \proofstepwidestar[1,2,3]{%
      \IntClass{a+c}{b+d}=\IntClass{a'+c'}{b'+d'}}{%
      \FormulaRefAuto{IntegerClassEquality}{1,5}}
  \closeproofpartwide
\end{tabproofsplitwide}

\FormulaThmDeltaK[Quotientenabstieg der Addition]{%
  \begin{aligned}[t]
    &z,w\in\mathbb Z\vdash
      \exists!u\in\mathbb Z\,
      \exists p,q,r,s\in\mathbb N\,{}\\[-2pt]
    &\qquad
      \left\{
      \begin{aligned}
        z&=\IntClass{p}{q},\\
        w&=\IntClass{r}{s},\\
        u&=\IntClass{p+r}{q+s}.
      \end{aligned}
      \right.
  \end{aligned}%
}{IntegerAdditionQuotientDescent}{%
  \DeltaRow{Mengen}{z\dsep w\dsep u\dsep v\dsep
    a\dsep b\dsep c\dsep d\dsep
    a'\dsep b'\dsep c'\dsep d'\dsep
    p\dsep q\dsep r\dsep s}%
}
\begin{tabproofsplitwide}
  \proofpartwideindR[Existenz durch Wahl von Repräsentanten]{%
    \begin{aligned}[t]
      &z,w\in\mathbb Z\vdash
        \exists u\in\mathbb Z\,
        \exists p,q,r,s\in\mathbb N\,{}\\[-2pt]
      &\qquad
        \left\{
        \begin{aligned}
          z&=\IntClass{p}{q},\\
          w&=\IntClass{r}{s},\\
          u&=\IntClass{p+r}{q+s}.
        \end{aligned}
        \right.
    \end{aligned}%
  }{%
    z,w\in\mathbb Z\vdash
    \exists u\in\mathbb Z\,
    \exists p,q,r,s\in\mathbb N\,
    \left\{
    \begin{aligned}
      z&=\IntClass{p}{q},\\
      w&=\IntClass{r}{s},\\
      u&=\IntClass{p+r}{q+s}.
    \end{aligned}
    \right.%
  }[IntegerAdditionDescentExistence]
    \proofstepwidestar[1]{z,w\in\mathbb Z}{\rA}
    \proofstepwidestar[1]{\exists a,b\in\mathbb N\,
      z=\IntClass{a}{b}}{%
      \FormulaRefAuto{IntegerPairRepresentative}{\rAEa{1}}}
    \proofstepwidestar[1]{\exists c,d\in\mathbb N\,
      w=\IntClass{c}{d}}{%
      \FormulaRefAuto{IntegerPairRepresentative}{\rAEb{1}}}
    \proofstepwidestar[4]{a,b\in\mathbb N\land
      z=\IntClass{a}{b}}{\rA}
    \proofstepwidestar[5]{c,d\in\mathbb N\land
      w=\IntClass{c}{d}}{\rA}
    \proofstepwidestar[4,5]{a,b,c,d\in\mathbb N}{%
      \rAI{\rAEa{4},\rAEa{5}}}
    \proofstepwidestar[4,5]{a+c\in\mathbb N}{%
      \FormulaRefAuto{PeanoAddClosure}{6}}
    \proofstepwidestar[4,5]{b+d\in\mathbb N}{%
      \FormulaRefAuto{PeanoAddClosure}{6}}
    \proofstepwidestar[4,5]{\IntClass{a+c}{b+d}\in\mathbb Z}{%
      \FormulaRefAuto{IntegerClassInZ}{7,8}}
    \proofstepwidestar[]{%
      \IntClass{a+c}{b+d}=\IntClass{a+c}{b+d}}{\rII}
    \proofstepwidestarbreak[4,5]{%
      \left\{
      \begin{aligned}
        z&=\IntClass{a}{b},\\
        w&=\IntClass{c}{d},\\
        \IntClass{a+c}{b+d}&=\IntClass{a+c}{b+d}.
      \end{aligned}
      \right.}{%
      \rAI{\rAEb{4},\rAI{\rAEb{5},10}}}
    \proofstepwidestarbreak[4,5]{%
      \exists p,q,r,s\in\mathbb N\,
      \left\{
      \begin{aligned}
        z&=\IntClass{p}{q},\\
        w&=\IntClass{r}{s},\\
        \IntClass{a+c}{b+d}&=\IntClass{p+r}{q+s}.
      \end{aligned}
      \right.}{%
      \rEI{\rAI{6,11}}}
    \proofstepwidestarbreak[4,5]{%
      \IntClass{a+c}{b+d}\in\mathbb Z\land
      \exists p,q,r,s\in\mathbb N\,
      \left\{
      \begin{aligned}
        z&=\IntClass{p}{q},\\
        w&=\IntClass{r}{s},\\
        \IntClass{a+c}{b+d}&=\IntClass{p+r}{q+s}.
      \end{aligned}
      \right.}{\rAI{9,12}}
    \proofstepwidestarbreak[4,5]{%
      \exists u\in\mathbb Z\,
      \exists p,q,r,s\in\mathbb N\,
      \left\{
      \begin{aligned}
        z&=\IntClass{p}{q},\\
        w&=\IntClass{r}{s},\\
        u&=\IntClass{p+r}{q+s}.
      \end{aligned}
      \right.}{\rEI{13}}
    \proofstepwidestarbreak[1,4]{%
      \exists u\in\mathbb Z\,
      \exists p,q,r,s\in\mathbb N\,
      \left\{
      \begin{aligned}
        z&=\IntClass{p}{q},\\
        w&=\IntClass{r}{s},\\
        u&=\IntClass{p+r}{q+s}.
      \end{aligned}
      \right.}{\rEE{3,5,14}}
    \proofstepwidestarbreak[1]{%
      \exists u\in\mathbb Z\,
      \exists p,q,r,s\in\mathbb N\,
      \left\{
      \begin{aligned}
        z&=\IntClass{p}{q},\\
        w&=\IntClass{r}{s},\\
        u&=\IntClass{p+r}{q+s}.
      \end{aligned}
      \right.}{\rEE{2,4,15}}
  \closeproofpartwideindR

  \proofpartwideindR[Repräsentantenunabhängigkeit und Eindeutigkeit]{%
    \begin{aligned}[t]
      &\left(
        u\in\mathbb Z\land
        \exists a,b,c,d\in\mathbb N\,
        \left\{
        \begin{aligned}
          z&=\IntClass{a}{b},\\
          w&=\IntClass{c}{d},\\
          u&=\IntClass{a+c}{b+d}
        \end{aligned}
        \right.
      \right)\dsep{}\\
      &\left(
        v\in\mathbb Z\land
        \exists a',b',c',d'\in\mathbb N\,
        \left\{
        \begin{aligned}
          z&=\IntClass{a'}{b'},\\
          w&=\IntClass{c'}{d'},\\
          v&=\IntClass{a'+c'}{b'+d'}
        \end{aligned}
        \right.
      \right)
      \vdash u=v
    \end{aligned}%
  }{%
    \left(
      u\in\mathbb Z\land
      \exists a,b,c,d\in\mathbb N\,
      \left\{
      \begin{aligned}
        z&=\IntClass{a}{b},\\
        w&=\IntClass{c}{d},\\
        u&=\IntClass{a+c}{b+d}
      \end{aligned}
      \right.
    \right)\dsep
    \left(
      v\in\mathbb Z\land
      \exists a',b',c',d'\in\mathbb N\,
      \left\{
      \begin{aligned}
        z&=\IntClass{a'}{b'},\\
        w&=\IntClass{c'}{d'},\\
        v&=\IntClass{a'+c'}{b'+d'}
      \end{aligned}
      \right.
    \right)
    \vdash u=v%
  }[IntegerAdditionDescentUniqueness]
    \proofstepwidestarbreak[1]{%
      u\in\mathbb Z\land
      \exists a,b,c,d\in\mathbb N\,
      \left\{
      \begin{aligned}
        z&=\IntClass{a}{b},\\
        w&=\IntClass{c}{d},\\
        u&=\IntClass{a+c}{b+d}
      \end{aligned}
      \right.}{\rA}
    \proofstepwidestarbreak[2]{%
      v\in\mathbb Z\land
      \exists a',b',c',d'\in\mathbb N\,
      \left\{
      \begin{aligned}
        z&=\IntClass{a'}{b'},\\
        w&=\IntClass{c'}{d'},\\
        v&=\IntClass{a'+c'}{b'+d'}
      \end{aligned}
      \right.}{\rA}
    \proofstepwidestarbreak[1]{%
      \exists a,b,c,d\in\mathbb N\,
      \left\{
      \begin{aligned}
        z&=\IntClass{a}{b},\\
        w&=\IntClass{c}{d},\\
        u&=\IntClass{a+c}{b+d}
      \end{aligned}
      \right.}{\rAEb{1}}
    \proofstepwidestarbreak[2]{%
      \exists a',b',c',d'\in\mathbb N\,
      \left\{
      \begin{aligned}
        z&=\IntClass{a'}{b'},\\
        w&=\IntClass{c'}{d'},\\
        v&=\IntClass{a'+c'}{b'+d'}
      \end{aligned}
      \right.}{\rAEb{2}}
    \proofstepwidestarbreak[5]{%
      a,b,c,d\in\mathbb N\land
      \left\{
      \begin{aligned}
        z&=\IntClass{a}{b},\\
        w&=\IntClass{c}{d},\\
        u&=\IntClass{a+c}{b+d}
      \end{aligned}
      \right.}{\rA}
    \proofstepwidestarbreak[6]{%
      a',b',c',d'\in\mathbb N\land
      \left\{
      \begin{aligned}
        z&=\IntClass{a'}{b'},\\
        w&=\IntClass{c'}{d'},\\
        v&=\IntClass{a'+c'}{b'+d'}
      \end{aligned}
      \right.}{\rA}
    \proofstepwidestarbreak[5,6]{%
      a,b,c,d,a',b',c',d'\in\mathbb N}{%
      \rAI{\rAEa{5},\rAEa{6}}}
    \proofstepwidestar[5]{z=\IntClass{a}{b}}{%
      \rAEa{\rAEb{5}}}
    \proofstepwidestar[5]{w=\IntClass{c}{d}}{%
      \rAEa{\rAEb{\rAEb{5}}}}
    \proofstepwidestar[5]{u=\IntClass{a+c}{b+d}}{%
      \rAEb{\rAEb{\rAEb{5}}}}
    \proofstepwidestar[6]{z=\IntClass{a'}{b'}}{%
      \rAEa{\rAEb{6}}}
    \proofstepwidestar[6]{w=\IntClass{c'}{d'}}{%
      \rAEa{\rAEb{\rAEb{6}}}}
    \proofstepwidestar[6]{v=\IntClass{a'+c'}{b'+d'}}{%
      \rAEb{\rAEb{\rAEb{6}}}}
    \proofstepwidestar[5,6]{%
      \IntClass{a}{b}=\IntClass{a'}{b'}}{\rIE{8,11}}
    \proofstepwidestar[5,6]{%
      \IntClass{c}{d}=\IntClass{c'}{d'}}{\rIE{9,12}}
    \proofstepwidestar[5,6]{%
      \IntClass{a+c}{b+d}=\IntClass{a'+c'}{b'+d'}}{%
      \FormulaRefAuto{IntegerAdditionCompatible}{7,14,15}}
    \proofstepwidestar[5,6]{u=v}{\rIE{10,13,16}}
    \proofstepwidestar[2,5]{u=v}{\rEE{4,6,17}}
    \proofstepwidestar[1,2]{u=v}{\rEE{3,5,18}}
  \closeproofpartwideindR

  \proofpartwideindR[Eindeutiger Existenzschluss]{%
    \begin{aligned}[t]
      &z,w\in\mathbb Z\vdash
        \exists!u\in\mathbb Z\,
        \exists p,q,r,s\in\mathbb N\,{}\\[-2pt]
      &\qquad
        \left\{
        \begin{aligned}
          z&=\IntClass{p}{q},\\
          w&=\IntClass{r}{s},\\
          u&=\IntClass{p+r}{q+s}.
        \end{aligned}
        \right.
    \end{aligned}%
  }{%
    z,w\in\mathbb Z\vdash
    \exists!u\in\mathbb Z\,
    \exists p,q,r,s\in\mathbb N\,
    \left\{
    \begin{aligned}
      z&=\IntClass{p}{q},\\
      w&=\IntClass{r}{s},\\
      u&=\IntClass{p+r}{q+s}.
    \end{aligned}
    \right.%
  }[IntegerAdditionDescentConclusion]
    \proofstepwidestar[1]{z,w\in\mathbb Z}{\rA}
    \proofstepwidestarbreak[1]{%
      \exists u\in\mathbb Z\,
      \exists p,q,r,s\in\mathbb N\,
      \left\{
      \begin{aligned}
        z&=\IntClass{p}{q},\\
        w&=\IntClass{r}{s},\\
        u&=\IntClass{p+r}{q+s}.
      \end{aligned}
      \right.}{%
      \FormulaRefAuto{IntegerAdditionDescentExistence}{1}}
    \proofstepwidestarbreak[3]{%
      u\in\mathbb Z\land
      \exists p,q,r,s\in\mathbb N\,
      \left\{
      \begin{aligned}
        z&=\IntClass{p}{q},\\
        w&=\IntClass{r}{s},\\
        u&=\IntClass{p+r}{q+s}
      \end{aligned}
      \right.}{\rA}
    \proofstepwidestarbreak[4]{%
      v\in\mathbb Z\land
      \exists p,q,r,s\in\mathbb N\,
      \left\{
      \begin{aligned}
        z&=\IntClass{p}{q},\\
        w&=\IntClass{r}{s},\\
        v&=\IntClass{p+r}{q+s}
      \end{aligned}
      \right.}{\rA}
    \proofstepwidestarbreak[3,4]{u=v}{%
      \FormulaRefAuto{IntegerAdditionDescentUniqueness}{3,4}}
    \proofstepwidestarbreak[1]{%
      \exists!u\in\mathbb Z\,
      \exists p,q,r,s\in\mathbb N\,
      \left\{
      \begin{aligned}
        z&=\IntClass{p}{q},\\
        w&=\IntClass{r}{s},\\
        u&=\IntClass{p+r}{q+s}.
      \end{aligned}
      \right.}{\UEI{2,3,4,5}}
  \closeproofpartwideindR
\end{tabproofsplitwide}
Der eindeutige Existenzsatz ist der Quotientenabstieg der Addition; die
folgende Klassenregel bezeichnet dieses eindeutig bestimmte Ergebnis.

\FormulaDefDeltaK[Addition ganzer Zahlen]{%
  a,b,c,d\in\mathbb N\vdash
  \IntClass{a}{b}+\IntClass{c}{d}
  \coloneqq\IntClass{a+c}{b+d}%
}{IntegerAdditionDef}{%
  \DeltaRow{Mengen}{a\dsep b\dsep c\dsep d}%
  \DeltaRow{Neue Symbole}{+}%
}

\FormulaThmDeltaK[Abgeschlossenheit der Addition ganzer Zahlen]{%
  z\in\mathbb Z\dsep w\in\mathbb Z\vdash z+w\in\mathbb Z%
}{IntAddClosure}{%
  \DeltaRow{Mengen}{z\dsep w\dsep a\dsep b\dsep c\dsep d}%
}
\begin{tabproofsplitwide}
  \proofpartwideindR[Abgeschlossenheit für feste Repräsentanten]{%
    \begin{aligned}[t]
      &a,b,c,d\in\mathbb N\land
        z=\IntClass{a}{b}\land w=\IntClass{c}{d}\\
      &\vdash z+w\in\mathbb Z
    \end{aligned}
  }{%
    a,b,c,d\in\mathbb N\land
    z=\IntClass{a}{b}\land w=\IntClass{c}{d}
    \vdash z+w\in\mathbb Z%
  }[IntAddClosureForRepresentatives]
    \proofstepwidestar[1]{a,b,c,d\in\mathbb N\land
      z=\IntClass{a}{b}\land w=\IntClass{c}{d}}{\rA}
    \proofstepwidestar[1]{z+w=\IntClass{a+c}{b+d}}{%
      \FormulaRefAuto{IntegerAdditionDef}{1}}
    \proofstepwidestar[1]{a+c\in\mathbb N}{%
      \FormulaRefAuto{PeanoAddClosure}{1}}
    \proofstepwidestar[1]{b+d\in\mathbb N}{%
      \FormulaRefAuto{PeanoAddClosure}{1}}
    \proofstepwidestar[1]{\IntClass{a+c}{b+d}\in\mathbb Z}{%
      \FormulaRefAuto{IntegerClassInZ}{3,4}}
    \proofstepwidestar[1]{z+w\in\mathbb Z}{\rIE{2,5}}
  \closeproofpartwideindR

  \proofpartwideindR[Elimination der Repräsentanten]{%
    z\in\mathbb Z\dsep w\in\mathbb Z\vdash z+w\in\mathbb Z%
  }{%
    z\in\mathbb Z\dsep w\in\mathbb Z\vdash z+w\in\mathbb Z%
  }[IntAddClosureRepresentativeElimination]
    \proofstepwidestar[1]{z\in\mathbb Z}{\rA}
    \proofstepwidestar[2]{w\in\mathbb Z}{\rA}
    \proofstepwidestar[1]{\exists a,b\in\mathbb N\,
      z=\IntClass{a}{b}}{\FormulaRefAuto{IntegerPairRepresentative}{1}}
    \proofstepwidestar[2]{\exists c,d\in\mathbb N\,
      w=\IntClass{c}{d}}{\FormulaRefAuto{IntegerPairRepresentative}{2}}
    \proofstepwidestar[5]{a,b,c,d\in\mathbb N\land
      z=\IntClass{a}{b}\land w=\IntClass{c}{d}}{\rA}
    \proofstepwidestar[5]{z+w\in\mathbb Z}{%
      \FormulaRefAuto{IntAddClosureForRepresentatives}{5}}
    \proofstepwidestar[1,2]{z+w\in\mathbb Z}{\rEE{3,4,5,6}}
  \closeproofpartwideindR
\end{tabproofsplitwide}

\FormulaThmDeltaK[Assoziativität der Addition ganzer Zahlen]{%
  z,w,u\in\mathbb Z\vdash(z+w)+u=z+(w+u)%
}{IntAddAssociative}{%
  \DeltaRow{Mengen}{z\dsep w\dsep u\dsep
    a\dsep b\dsep c\dsep d\dsep e\dsep f}%
}
\begin{tabproofsplitwide}
  \proofpartwideindR[Rechnung mit Repräsentanten]{%
    \begin{aligned}[t]
      &a,b,c,d,e,f\in\mathbb N\vdash{}\\
      &(\IntClass{a}{b}+\IntClass{c}{d})+\IntClass{e}{f}
       =\IntClass{a}{b}+(\IntClass{c}{d}+\IntClass{e}{f})
    \end{aligned}
  }{%
    a,b,c,d,e,f\in\mathbb N\vdash
    (\IntClass{a}{b}+\IntClass{c}{d})+\IntClass{e}{f}
    =\IntClass{a}{b}+(\IntClass{c}{d}+\IntClass{e}{f})
  }[IntAddAssociativeOnClasses]
    \proofstepwidestar[1]{a,b,c,d,e,f\in\mathbb N}{\rA}
    \proofstepwidestar[1]{%
      \IntClass{a}{b}+(\IntClass{c}{d}+\IntClass{e}{f})
      =\IntClass{a+(c+e)}{b+(d+f)}}{%
      \FormulaRefAuto{IntegerAdditionDef}{1}}
    \proofstepwidestar[1]{(a+c)+e=a+(c+e)}{%
      \FormulaRefAuto{PeanoAddAssociative}{1}}
    \proofstepwidestar[1]{(b+d)+f=b+(d+f)}{%
      \FormulaRefAuto{PeanoAddAssociative}{1}}
    \proofstepwide[1]{%
      (\IntClass{a}{b}+\IntClass{c}{d})+\IntClass{e}{f}}{=}{%
      \IntClass{(a+c)+e}{(b+d)+f}}{%
      \FormulaRefAuto{IntegerAdditionDef}{1}}
    \proofstepwide[1]{}{=}{%
      \IntClass{a+(c+e)}{b+(d+f)}}{\rIE{3,4}}
    \proofstepwide[1]{}{=}{%
      \IntClass{a}{b}+(\IntClass{c}{d}+\IntClass{e}{f})}{%
      \FormulaRefAuto{a=b\vdash b=a}{2}}
    \proofstepwidestar[1]{%
      (\IntClass{a}{b}+\IntClass{c}{d})+\IntClass{e}{f}
      =\IntClass{a}{b}+(\IntClass{c}{d}+\IntClass{e}{f})}{%
      \rChain{5,6,7}}
  \closeproofpartwideindR

  \proofpartwide{Schluss}
    \proofstepwidestar[1]{z,w,u\in\mathbb Z}{\rA}
    \proofstepwidestar[1]{z\in\mathbb Z}{\rAEa{1}}
    \proofstepwidestar[1]{w\in\mathbb Z}{\rAEa{\rAEb{1}}}
    \proofstepwidestar[1]{u\in\mathbb Z}{\rAEb{\rAEb{1}}}
    \proofstepwidestar[1]{\exists a,b\in\mathbb N\,
      z=\IntClass{a}{b}}{\FormulaRefAuto{IntegerPairRepresentative}{2}}
    \proofstepwidestar[1]{\exists c,d\in\mathbb N\,
      w=\IntClass{c}{d}}{\FormulaRefAuto{IntegerPairRepresentative}{3}}
    \proofstepwidestar[1]{\exists e,f\in\mathbb N\,
      u=\IntClass{e}{f}}{\FormulaRefAuto{IntegerPairRepresentative}{4}}
    \proofstepwidestar[8]{a,b,c,d,e,f\in\mathbb N\land
      z=\IntClass{a}{b}\land w=\IntClass{c}{d}\land
      u=\IntClass{e}{f}}{\rA}
    \proofstepwidestar[8]{%
      (\IntClass{a}{b}+\IntClass{c}{d})+\IntClass{e}{f}
      =\IntClass{a}{b}+(\IntClass{c}{d}+\IntClass{e}{f})}{%
      \FormulaRefAuto{IntAddAssociativeOnClasses}{8}}
    \proofstepwidestar[8]{(z+w)+u=z+(w+u)}{\rIE{8,9}}
    \proofstepwidestar[1]{(z+w)+u=z+(w+u)}{%
      \rEE{5,6,7,8,10}}
  \closeproofpartwide
\end{tabproofsplitwide}

\FormulaThmDeltaK[Kommutativität der Addition ganzer Zahlen]{%
  z,w\in\mathbb Z\vdash z+w=w+z%
}{IntAddCommutative}{%
  \DeltaRow{Mengen}{z\dsep w\dsep a\dsep b\dsep c\dsep d}%
}
\begin{tabproofsplitwide}
  \proofpartwideindR[Vertauschung fester Repräsentanten]{%
    \begin{aligned}[t]
      &a,b,c,d\in\mathbb N\land
        z=\IntClass{a}{b}\land w=\IntClass{c}{d}\\
      &\vdash z+w=w+z
    \end{aligned}
  }{%
    a,b,c,d\in\mathbb N\land
    z=\IntClass{a}{b}\land w=\IntClass{c}{d}
    \vdash z+w=w+z%
  }[IntAddCommutativeOnRepresentatives]
    \proofstepwidestar[1]{a,b,c,d\in\mathbb N\land
      z=\IntClass{a}{b}\land w=\IntClass{c}{d}}{\rA}
    \proofstepwidestar[1]{w+z=\IntClass{c+a}{d+b}}{%
      \FormulaRefAuto{IntegerAdditionDef}{1}}
    \proofstepwidestar[1]{a+c=c+a}{%
      \FormulaRefAuto{PeanoAddCommutative}{1}}
    \proofstepwidestar[1]{b+d=d+b}{%
      \FormulaRefAuto{PeanoAddCommutative}{1}}
    \proofstepwide[1]{z+w}{=}{\IntClass{a+c}{b+d}}{%
      \FormulaRefAuto{IntegerAdditionDef}{1}}
    \proofstepwide[1]{}{=}{\IntClass{c+a}{d+b}}{\rIE{3,4}}
    \proofstepwide[1]{}{=}{w+z}{%
      \FormulaRefAuto{a=b\vdash b=a}{2}}
    \proofstepwidestar[1]{z+w=w+z}{\rChain{5,6,7}}
  \closeproofpartwideindR

  \proofpartwideindR[Getrennte Repräsentantenwahl]{%
    z,w\in\mathbb Z\vdash z+w=w+z%
  }{%
    z,w\in\mathbb Z\vdash z+w=w+z%
  }[IntAddCommutativeRepresentativeElimination]
    \proofstepwidestar[1]{z,w\in\mathbb Z}{\rA}
    \proofstepwidestar[1]{z\in\mathbb Z}{\rAEa{1}}
    \proofstepwidestar[1]{w\in\mathbb Z}{\rAEb{1}}
    \proofstepwidestar[1]{\exists a,b\in\mathbb N\,
      z=\IntClass{a}{b}}{\FormulaRefAuto{IntegerPairRepresentative}{2}}
    \proofstepwidestar[1]{\exists c,d\in\mathbb N\,
      w=\IntClass{c}{d}}{\FormulaRefAuto{IntegerPairRepresentative}{3}}
    \proofstepwidestar[6]{a,b,c,d\in\mathbb N\land
      z=\IntClass{a}{b}\land w=\IntClass{c}{d}}{\rA}
    \proofstepwidestar[6]{z+w=w+z}{%
      \FormulaRefAuto{IntAddCommutativeOnRepresentatives}{6}}
    \proofstepwidestar[1]{z+w=w+z}{\rEE{4,5,6,7}}
  \closeproofpartwideindR
\end{tabproofsplitwide}

\FormulaThmDeltaK[Nullgesetze der Addition ganzer Zahlen]{%
  z\in\mathbb Z\vdash
  z+\IntZero=z\land\IntZero+z=z%
}{IntAddZero}{%
  \DeltaRow{Mengen}{z\dsep a\dsep b}%
}
\begin{tabproofsplitwide}
  \proofpartwideindR[Rechtsneutrales Element für einen Repräsentanten]{%
    \begin{aligned}[t]
      &a,b\in\mathbb N\land z=\IntClass{a}{b}\\
      &\vdash z+\IntZero=z
    \end{aligned}
  }{%
    a,b\in\mathbb N\land z=\IntClass{a}{b}
    \vdash z+\IntZero=z%
  }[IntAddRightZeroForRepresentative]
    \proofstepwidestar[1]{a,b\in\mathbb N\land
      z=\IntClass{a}{b}}{\rA}
    \proofstepwidestar[]{\IntZero=\IntClass{0}{0}}{%
      \FormulaRefAuto{IntegerZeroClass}}
    \proofstepwidestar[1]{a+0=a}{%
      \FormulaRefAuto{PeanoAddRightZero}{1}}
    \proofstepwidestar[1]{b+0=b}{%
      \FormulaRefAuto{PeanoAddRightZero}{1}}
    \proofstepwide[1]{z+\IntZero}{=}{\IntClass{a+0}{b+0}}{%
      \FormulaRefAuto{IntegerAdditionDef}{1,2}}
    \proofstepwide[1]{}{=}{\IntClass{a}{b}}{\rIE{3,4}}
    \proofstepwide[1]{}{=}{z}{%
      \FormulaRefAuto{a=b\vdash b=a}{\rAEb{1}}}
    \proofstepwidestar[1]{z+\IntZero=z}{\rChain{5,6,7}}
  \closeproofpartwideindR

  \proofpartwideindR[Elimination des Repräsentanten]{%
    z\in\mathbb Z\vdash
    z+\IntZero=z\land\IntZero+z=z%
  }{%
    z\in\mathbb Z\vdash
    z+\IntZero=z\land\IntZero+z=z%
  }[IntAddZeroRepresentativeElimination]
    \proofstepwidestar[1]{z\in\mathbb Z}{\rA}
    \proofstepwidestar[1]{\exists a,b\in\mathbb N\,
      z=\IntClass{a}{b}}{\FormulaRefAuto{IntegerPairRepresentative}{1}}
    \proofstepwidestar[3]{a,b\in\mathbb N\land
      z=\IntClass{a}{b}}{\rA}
    \proofstepwidestar[3]{z+\IntZero=z}{%
      \FormulaRefAuto{IntAddRightZeroForRepresentative}{3}}
    \proofstepwidestar[3]{\IntZero+z=z}{%
      \FormulaRefAuto{IntAddCommutative}{%
        \FormulaRefAuto{IntZeroCarrier},1,
        \FormulaRefAuto{a=b\vdash b=a}{4}}}
    \proofstepwidestar[3]{%
      z+\IntZero=z\land\IntZero+z=z}{\rAI{4,5}}
    \proofstepwidestar[1]{%
      z+\IntZero=z\land\IntZero+z=z}{\rEE{2,3,6}}
  \closeproofpartwideindR
\end{tabproofsplitwide}

\FormulaThmDeltaK[Additive Inversenregeln]{%
  z\in\mathbb Z\vdash
  z+(-z)=\IntZero\land(-z)+z=\IntZero%
}{IntAddInverse}{%
  \DeltaRow{Mengen}{z\dsep a\dsep b}%
}
\begin{tabproofsplitwide}
  \proofpartwideindR[Inverse Rechnung auf einer Klasse]{%
    a,b\in\mathbb N\vdash
    \IntClass{a}{b}+(-\IntClass{a}{b})=\IntZero%
  }{%
    a,b\in\mathbb N\vdash
    \IntClass{a}{b}+(-\IntClass{a}{b})=\IntZero%
  }[IntAddInverseOnClass]
    \proofstepwidestar[1]{a,b\in\mathbb N}{\rA}
    \proofstepwidestar[1]{-\IntClass{a}{b}=\IntClass{b}{a}}{%
      \FormulaRefAuto{IntegerNegationDef}{1}}
    \proofstepwidestar[1]{%
      \IntClass{a}{b}+(-\IntClass{a}{b})=\IntClass{a+b}{b+a}}{%
      \FormulaRefAuto{IntegerAdditionDef}{1,2}}
    \proofstepwidestar[1]{a+b\in\mathbb N}{%
      \FormulaRefAuto{PeanoAddClosure}{1}}
    \proofstepwidestar[1]{b+a\in\mathbb N}{%
      \FormulaRefAuto{PeanoAddClosure}{1}}
    \proofstepwidestar[]{0\in\mathbb N}{\FormulaRefAuto{PeanoZero}}
    \proofstepwidestar[1]{(a+b)+0=(b+a)+0}{%
      \FormulaRefAuto{PeanoAddCongruence}{%
        4,5,6,6,\FormulaRefAuto{PeanoAddCommutative}{1},\rII}}
    \proofstepwidestar[1]{\IntClass{a+b}{b+a}=\IntClass{0}{0}}{%
      \FormulaRefAuto{IntegerClassEquality}{4,5,6,6,7}}
    \proofstepwidestar[]{\IntClass{0}{0}=\IntZero}{%
      \FormulaRefAuto{a=b\vdash b=a}{%
        \FormulaRefAuto{IntegerZeroClass}}}
    \proofstepwidestar[1]{%
      \IntClass{a}{b}+(-\IntClass{a}{b})=\IntZero}{\rChain{3,8,9}}
  \closeproofpartwideindR

  \proofpartwide{Schluss}
    \proofstepwidestarbreakkeep[1]{z\in\mathbb Z}{\rA}
    \proofstepwidestarbreakkeep[1]{\exists a,b\in\mathbb N\,
      z=\IntClass{a}{b}}{\FormulaRefAuto{IntegerPairRepresentative}{1}}
    \proofstepwidestarbreakkeep[3]{a,b\in\mathbb N\land z=\IntClass{a}{b}}{\rA}
    \proofstepwidestarbreakkeep[3]{z+(-z)=\IntZero}{%
      \FormulaRefAuto{IntAddInverseOnClass}{3}}
    \proofstepwidestarbreakkeep[1]{-z\in\mathbb Z}{\FormulaRefAuto{IntNegClosure}{1}}
    \proofstepwidestarbreakkeep[1,3]{(-z)+z=z+(-z)}{%
      \FormulaRefAuto{IntAddCommutative}{5,1}}
    \proofstepwidestarbreakkeep[1,3]{(-z)+z=\IntZero}{\rChain{6,4}}
    \proofstepwidestarbreakkeep[1,3]{%
      z+(-z)=\IntZero\land(-z)+z=\IntZero}{\rAI{4,7}}
    \proofstepwidestarbreakkeep[1]{%
      z+(-z)=\IntZero\land(-z)+z=\IntZero}{\rEE{2,3,8}}
  \closeproofpartwide
\end{tabproofsplitwide}

\section{Multiplikation}

Metasprachlich lesen wir \(\IntClass{a}{b}\) als \(a-b\) und
\(\IntClass{c}{d}\) als \(c-d\). Die vertraute Ausmultiplizierung
\[
  (a-b)(c-d)=(ac+bd)-(ad+bc)
\]
legt daher das Differenzenpaar \((ac+bd,ad+bc)\) als Produktrepräsentanten
nahe. Da eine ganze Zahl viele Differenzenpaare besitzt, ist diese Formel
aber erst dann eine zulässige Definition auf den Klassen, wenn ein Wechsel
beider Repräsentanten dieselbe Ganzzahlklasse liefert. Genau diese
Repräsentantenunabhängigkeit wird im folgenden Satz bewiesen; erst danach
steigt die Operation auf den Quotienten \(\mathbb Z\) ab.

\FormulaThmDeltaK[Repräsentantenunabhängigkeit der Multiplikation]{%
  \begin{aligned}[t]
    &a,b,a',b',c,d,c',d'\in\mathbb N\dsep
      \IntClass{a}{b}=\IntClass{a'}{b'}\dsep
      \IntClass{c}{d}=\IntClass{c'}{d'}\vdash{}\\
    &\IntClass{a\cdot c+b\cdot d}{a\cdot d+b\cdot c}
      =
      \IntClass{a'\cdot c'+b'\cdot d'}{a'\cdot d'+b'\cdot c'}
  \end{aligned}%
}{IntegerMultiplicationCompatible}{%
  \DeltaRow{Mengen}{a\dsep b\dsep a'\dsep b'\dsep
    c\dsep d\dsep c'\dsep d'}%
}
\begin{tabproofsplitwide}
  \proofpartwideindR[Wechsel des ersten Repräsentanten]{%
    \begin{aligned}[t]
      &a,b,a',b',c,d\in\mathbb N\dsep
      \IntClass{a}{b}=\IntClass{a'}{b'}\vdash{}\\
      &\IntClass{a c+b d}{a d+b c}
       =\IntClass{a'c+b'd}{a'd+b'c}
    \end{aligned}
  }{%
    a,b,a',b',c,d\in\mathbb N\dsep
    \IntClass{a}{b}=\IntClass{a'}{b'}
    \vdash
    \IntClass{a c+b d}{a d+b c}
    =\IntClass{a'c+b'd}{a'd+b'c}
  }[IntegerMultiplicationCompatibleLeft]
    \proofstepwidestar[1]{a,b,a',b',c,d\in\mathbb N}{\rA}
    \proofstepwidestar[2]{\IntClass{a}{b}=\IntClass{a'}{b'}}{\rA}
    \proofstepwidestar[1,2]{a+b'=b+a'}{%
      \FormulaRefAuto{IntegerClassEquality}{1,2}}
    \proofstepwidestar[1,2]{ac+b'c=bc+a'c}{%
      \FormulaRefAuto{PeanoScaledSumEqualityRight}{1,3}}
    \proofstepwidestar[1,2]{bd+a'd=ad+b'd}{%
      \FormulaRefAuto{a=b\vdash b=a}{%
        \FormulaRefAuto{PeanoScaledSumEqualityRight}{1,3}}}
    \proofstepwidestar[1]{a'd+b'c=b'c+a'd}{%
      \FormulaRefAuto{PeanoAddCommutative}{1}}
    \proofstepwidestar[1]{bc+ad=ad+bc}{%
      \FormulaRefAuto{PeanoAddCommutative}{1}}
    \proofstepwidestarbreakkeep[1]{ac,b'c,bc,a'c\in\mathbb N}{%
      \rAIRepeat{4}{\FormulaRefAuto{PeanoMulClosure}}{1}}
    \proofstepwidestarbreakkeep[1]{bd,a'd,ad,b'd\in\mathbb N}{%
      \rAIRepeat{4}{\FormulaRefAuto{PeanoMulClosure}}{1}}
    \proofstepwidestarbreakkeep[1,2]{%
      (ac+bd)+(b'c+a'd)=(bc+ad)+(a'c+b'd)}{%
      \FormulaRefAuto{PeanoSumEqualitiesAddition}{\rAI{8,9},4,5}}
    \proofstepwidestarbreakkeep[1,2]{%
      (ac+bd)+(a'd+b'c)=(ad+bc)+(a'c+b'd)}{%
      \rIE{6,10,7}}
    \proofstepwidestarbreakkeep[1,2]{%
      \IntClass{ac+bd}{ad+bc}=\IntClass{a'c+b'd}{a'd+b'c}}{%
      \FormulaRefAuto{IntegerClassEquality}{1,11}}
  \closeproofpartwideindR

  \proofpartwideindR[Wechsel des zweiten Repräsentanten]{%
    \begin{aligned}[t]
      &a,b,c,d,c',d'\in\mathbb N\dsep
      \IntClass{c}{d}=\IntClass{c'}{d'}\vdash{}\\
      &\IntClass{a c+b d}{a d+b c}
       =\IntClass{a c'+b d'}{a d'+b c'}
    \end{aligned}
  }{%
    a,b,c,d,c',d'\in\mathbb N\dsep
    \IntClass{c}{d}=\IntClass{c'}{d'}
    \vdash
    \IntClass{a c+b d}{a d+b c}
    =\IntClass{a c'+b d'}{a d'+b c'}
  }[IntegerMultiplicationCompatibleRight]
    \proofstepwidestar[1]{a,b,c,d,c',d'\in\mathbb N}{\rA}
    \proofstepwidestar[2]{\IntClass{c}{d}=\IntClass{c'}{d'}}{\rA}
    \proofstepwidestar[1,2]{c+d'=d+c'}{%
      \FormulaRefAuto{IntegerClassEquality}{1,2}}
    \proofstepwidestar[1,2]{ac+ad'=ad+ac'}{%
      \FormulaRefAuto{PeanoScaledSumEqualityLeft}{1,3}}
    \proofstepwidestar[1,2]{bd+bc'=bc+bd'}{%
      \FormulaRefAuto{a=b\vdash b=a}{%
        \FormulaRefAuto{PeanoScaledSumEqualityLeft}{1,3}}}
    \proofstepwidestarbreakkeep[1]{ac,ad',ad,ac'\in\mathbb N}{%
      \rAIRepeat{4}{\FormulaRefAuto{PeanoMulClosure}}{1}}
    \proofstepwidestarbreakkeep[1]{bd,bc',bc,bd'\in\mathbb N}{%
      \rAIRepeat{4}{\FormulaRefAuto{PeanoMulClosure}}{1}}
    \proofstepwidestar[1,2]{%
      (ac+bd)+(ad'+bc')=(ad+bc)+(ac'+bd')}{%
      \FormulaRefAuto{PeanoSumEqualitiesAddition}{\rAI{6,7},4,5}}
    \proofstepwidestar[1,2]{%
      \IntClass{ac+bd}{ad+bc}=\IntClass{ac'+bd'}{ad'+bc'}}{%
      \FormulaRefAuto{IntegerClassEquality}{1,8}}
  \closeproofpartwideindR

  \proofpartwide{Zusammenführung}
    \proofstepwidestar[1]{a,b,a',b',c,d,c',d'\in\mathbb N}{\rA}
    \proofstepwidestar[2]{\IntClass{a}{b}=\IntClass{a'}{b'}}{\rA}
    \proofstepwidestar[3]{\IntClass{c}{d}=\IntClass{c'}{d'}}{\rA}
    \proofstepwidestar[1,2]{%
      \IntClass{ac+bd}{ad+bc}
      =\IntClass{a'c+b'd}{a'd+b'c}}{%
      \FormulaRefAuto{IntegerMultiplicationCompatibleLeft}{1,2}}
    \proofstepwidestar[1,3]{%
      \IntClass{a'c+b'd}{a'd+b'c}
      =\IntClass{a'c'+b'd'}{a'd'+b'c'}}{%
      \FormulaRefAuto{IntegerMultiplicationCompatibleRight}{1,3}}
    \proofstepwidestar[1,2,3]{%
      \IntClass{ac+bd}{ad+bc}
      =\IntClass{a'c'+b'd'}{a'd'+b'c'}}{%
      \rChain{4,5}}
  \closeproofpartwide
\end{tabproofsplitwide}

\FormulaThmDeltaK[Quotientenabstieg der Multiplikation]{%
  \begin{aligned}[t]
    &z,w\in\mathbb Z\vdash
    \exists! u\in\mathbb Z\,
    \exists p\in\mathbb N\,\exists q\in\mathbb N\,
    \exists r\in\mathbb N\,\exists s\in\mathbb N\,{}\\
    &\qquad
    \bigl(z=\IntClass{p}{q}\land
      w=\IntClass{r}{s}\land
      u=\IntClass{p r+q s}{p s+q r}\bigr)
  \end{aligned}%
}{IntegerMultiplicationQuotientDescent}{%
  \DeltaRow{Mengen}{z\dsep w\dsep u\dsep v\dsep
    a\dsep b\dsep c\dsep d\dsep
    a'\dsep b'\dsep c'\dsep d'\dsep
    p\dsep q\dsep r\dsep s}%
}
\begin{tabproofsplitwide}
  \proofpartwideindR[Abgeschlossenheit der Ergebniskoordinaten]{%
    a,b,c,d\in\mathbb N\vdash
    ac+bd\in\mathbb N\land ad+bc\in\mathbb N%
  }{%
    a,b,c,d\in\mathbb N\vdash
    ac+bd\in\mathbb N\land ad+bc\in\mathbb N%
  }[IntegerMultiplicationDescentCoordinates]
    \proofstepwidestarbreakkeep[1]{a,b,c,d\in\mathbb N}{\rA}
    \proofstepwidestarbreakkeep[1]{ac\in\mathbb N}{%
      \FormulaRefAuto{PeanoMulClosure}{1}}
    \proofstepwidestarbreakkeep[1]{bd\in\mathbb N}{%
      \FormulaRefAuto{PeanoMulClosure}{1}}
    \proofstepwidestarbreakkeep[1]{ad\in\mathbb N}{%
      \FormulaRefAuto{PeanoMulClosure}{1}}
    \proofstepwidestarbreakkeep[1]{bc\in\mathbb N}{%
      \FormulaRefAuto{PeanoMulClosure}{1}}
    \proofstepwidestarbreakkeep[1]{ac+bd\in\mathbb N}{%
      \FormulaRefAuto{PeanoAddClosure}{2,3}}
    \proofstepwidestarbreakkeep[1]{ad+bc\in\mathbb N}{%
      \FormulaRefAuto{PeanoAddClosure}{4,5}}
    \proofstepwidestarbreakkeep[1]{%
      ac+bd\in\mathbb N\land ad+bc\in\mathbb N}{\rAI{6,7}}
  \closeproofpartwideindR

  \proofpartwideindR[Ergebniskandidat fester Repräsentanten]{%
    \begin{aligned}[t]
      &a,b,c,d\in\mathbb N\dsep
      z=\IntClass{a}{b}\dsep w=\IntClass{c}{d}\vdash{}\\
      &\IntClass{ac+bd}{ad+bc}\in\mathbb Z\land
      \exists p\in\mathbb N\,\exists q\in\mathbb N\,
      \exists r\in\mathbb N\,\exists s\in\mathbb N\,{}\\
      &\quad
      \bigl(z=\IntClass{p}{q}\land
        w=\IntClass{r}{s}\land
        \IntClass{ac+bd}{ad+bc}
          =\IntClass{p r+q s}{p s+q r}\bigr)
    \end{aligned}%
  }{%
    a,b,c,d\in\mathbb N\dsep
    z=\IntClass{a}{b}\dsep w=\IntClass{c}{d}\vdash
    \IntClass{ac+bd}{ad+bc}\in\mathbb Z\land
    \exists p\in\mathbb N\,\exists q\in\mathbb N\,
    \exists r\in\mathbb N\,\exists s\in\mathbb N\,
    \bigl(z=\IntClass{p}{q}\land
      w=\IntClass{r}{s}\land
      \IntClass{ac+bd}{ad+bc}
        =\IntClass{p r+q s}{p s+q r}\bigr)%
  }[IntegerMultiplicationDescentCandidate]
    \proofstepwidestarbreakkeep[1]{a,b,c,d\in\mathbb N}{\rA}
    \proofstepwidestarbreakkeep[2]{z=\IntClass{a}{b}}{\rA}
    \proofstepwidestarbreakkeep[3]{w=\IntClass{c}{d}}{\rA}
    \proofstepwidestarbreakkeep[1]{%
      ac+bd\in\mathbb N\land ad+bc\in\mathbb N}{%
      \FormulaRefAuto{IntegerMultiplicationDescentCoordinates}{1}}
    \proofstepwidestarbreakkeep[1]{\IntClass{ac+bd}{ad+bc}\in\mathbb Z}{%
      \FormulaRefAuto{IntegerClassInZ}{\rAEa{4},\rAEb{4}}}
    \proofstepwidestarbreakkeep[]{%
      \IntClass{ac+bd}{ad+bc}=\IntClass{ac+bd}{ad+bc}}{\rII}
    \proofstepwidestarbreakkeep[1,2,3]{%
      \begin{aligned}[t]
        &\exists p,q,r,s\in\mathbb N\,{}\\
        &\bigl(z=\IntClass{p}{q}\land w=\IntClass{r}{s}\land{}\\
        &\qquad\IntClass{ac+bd}{ad+bc}
          =\IntClass{p r+q s}{p s+q r}\bigr)
      \end{aligned}}{%
      \rEI{\rAI{1,\rAI{2,\rAI{3,6}}}}}
    \proofstepwidestarbreakkeep[1,2,3]{%
      \begin{aligned}[t]
        &\IntClass{ac+bd}{ad+bc}\in\mathbb Z\land
          \exists p,q,r,s\in\mathbb N\,{}\\
        &\bigl(z=\IntClass{p}{q}\land w=\IntClass{r}{s}\land{}\\
        &\qquad\IntClass{ac+bd}{ad+bc}
          =\IntClass{p r+q s}{p s+q r}\bigr)
      \end{aligned}}{\rAI{5,7}}
  \closeproofpartwideindR

  \proofpartwideindR[Existenz durch getrennte Repräsentantenwahl]{%
    \begin{aligned}[t]
      &z,w\in\mathbb Z\vdash
      \exists u\in\mathbb Z\,
      \exists p\in\mathbb N\,\exists q\in\mathbb N\,
      \exists r\in\mathbb N\,\exists s\in\mathbb N\,{}\\
      &\qquad
      \bigl(z=\IntClass{p}{q}\land
        w=\IntClass{r}{s}\land
        u=\IntClass{p r+q s}{p s+q r}\bigr)
    \end{aligned}%
  }{%
    z,w\in\mathbb Z\vdash
    \exists u\in\mathbb Z\,
    \exists p\in\mathbb N\,\exists q\in\mathbb N\,
    \exists r\in\mathbb N\,\exists s\in\mathbb N\,
    \bigl(z=\IntClass{p}{q}\land
      w=\IntClass{r}{s}\land
      u=\IntClass{p r+q s}{p s+q r}\bigr)%
  }[IntegerMultiplicationDescentExistence]
    \proofstepwidestarbreakkeep[1]{z,w\in\mathbb Z}{\rA}
    \proofstepwidestarbreakkeep[1]{z\in\mathbb Z}{\rAEa{1}}
    \proofstepwidestarbreakkeep[1]{w\in\mathbb Z}{\rAEb{1}}
    \proofstepwidestarbreakkeep[1]{%
      \exists a\in\mathbb N\,\exists b\in\mathbb N\,
      z=\IntClass{a}{b}}{%
      \FormulaRefAuto{IntegerPairRepresentative}{2}}
    \proofstepwidestarbreakkeep[1]{%
      \exists c\in\mathbb N\,\exists d\in\mathbb N\,
      w=\IntClass{c}{d}}{%
      \FormulaRefAuto{IntegerPairRepresentative}{3}}
    \proofstepwidestarbreakkeep[6]{%
      a,b\in\mathbb N\land z=\IntClass{a}{b}}{\rA}
    \proofstepwidestarbreakkeep[7]{%
      c,d\in\mathbb N\land w=\IntClass{c}{d}}{\rA}
    \proofstepwidestarbreakkeep[6]{a,b\in\mathbb N}{\rAEa{6}}
    \proofstepwidestarbreakkeep[6]{z=\IntClass{a}{b}}{\rAEb{6}}
    \proofstepwidestarbreakkeep[7]{c,d\in\mathbb N}{\rAEa{7}}
    \proofstepwidestarbreakkeep[7]{w=\IntClass{c}{d}}{\rAEb{7}}
    \proofstepwidestarbreakkeep[6,7]{a,b,c,d\in\mathbb N}{\rAI{8,10}}
    \proofstepwidestarbreakkeep[6,7]{%
      \begin{aligned}[t]
        &\IntClass{ac+bd}{ad+bc}\in\mathbb Z\land
          \exists p,q,r,s\in\mathbb N\,{}\\
        &\bigl(z=\IntClass{p}{q}\land w=\IntClass{r}{s}\land{}\\
        &\qquad\IntClass{ac+bd}{ad+bc}
          =\IntClass{p r+q s}{p s+q r}\bigr)
      \end{aligned}}{%
      \FormulaRefAuto{IntegerMultiplicationDescentCandidate}{12,9,11}}
    \proofstepwidestarbreakkeep[6,7]{%
      \begin{aligned}[t]
        &\exists u\in\mathbb Z\,\exists p,q,r,s\in\mathbb N\,{}\\
        &\bigl(z=\IntClass{p}{q}\land w=\IntClass{r}{s}\land
          u=\IntClass{p r+q s}{p s+q r}\bigr)
      \end{aligned}}{\rEI{13}}
    \proofstepwidestarbreakkeep[6]{%
      \begin{aligned}[t]
        &\exists u\in\mathbb Z\,\exists p,q,r,s\in\mathbb N\,{}\\
        &\bigl(z=\IntClass{p}{q}\land w=\IntClass{r}{s}\land
          u=\IntClass{p r+q s}{p s+q r}\bigr)
      \end{aligned}}{\rEE{5,7,14}}
    \proofstepwidestarbreakkeep[1]{%
      \begin{aligned}[t]
        &\exists u\in\mathbb Z\,\exists p,q,r,s\in\mathbb N\,{}\\
        &\bigl(z=\IntClass{p}{q}\land w=\IntClass{r}{s}\land
          u=\IntClass{p r+q s}{p s+q r}\bigr)
      \end{aligned}}{\rEE{4,6,15}}
  \closeproofpartwideindR

  \proofpartwideindR[Eindeutigkeit für feste Repräsentanten]{%
    \begin{aligned}[t]
      &a,b,c,d,a',b',c',d'\in\mathbb N\dsep{}\\
      &z=\IntClass{a}{b}\land w=\IntClass{c}{d}
        \land u=\IntClass{ac+bd}{ad+bc}\dsep{}\\
      &z=\IntClass{a'}{b'}\land w=\IntClass{c'}{d'}
        \land v=\IntClass{a'c'+b'd'}{a'd'+b'c'}
      \vdash u=v
    \end{aligned}%
  }{%
    a,b,c,d,a',b',c',d'\in\mathbb N\dsep
    z=\IntClass{a}{b}\land w=\IntClass{c}{d}
      \land u=\IntClass{ac+bd}{ad+bc}\dsep
    z=\IntClass{a'}{b'}\land w=\IntClass{c'}{d'}
      \land v=\IntClass{a'c'+b'd'}{a'd'+b'c'}
    \vdash u=v%
  }[IntegerMultiplicationDescentFixedUniqueness]
    \proofstepwidestarbreakkeep[1]{%
      a,b,c,d,a',b',c',d'\in\mathbb N}{\rA}
    \proofstepwidestarbreakkeep[2]{%
      z=\IntClass{a}{b}\land w=\IntClass{c}{d}
      \land u=\IntClass{ac+bd}{ad+bc}}{\rA}
    \proofstepwidestarbreakkeep[3]{%
      z=\IntClass{a'}{b'}\land w=\IntClass{c'}{d'}
      \land v=\IntClass{a'c'+b'd'}{a'd'+b'c'}}{\rA}
    \proofstepwidestarbreakkeep[2]{z=\IntClass{a}{b}}{\rAEa{2}}
    \proofstepwidestarbreakkeep[2]{w=\IntClass{c}{d}}{\rAEa{\rAEb{2}}}
    \proofstepwidestarbreakkeep[2]{u=\IntClass{ac+bd}{ad+bc}}{\rAEb{\rAEb{2}}}
    \proofstepwidestarbreakkeep[3]{z=\IntClass{a'}{b'}}{\rAEa{3}}
    \proofstepwidestarbreakkeep[3]{w=\IntClass{c'}{d'}}{\rAEa{\rAEb{3}}}
    \proofstepwidestarbreakkeep[3]{%
      v=\IntClass{a'c'+b'd'}{a'd'+b'c'}}{\rAEb{\rAEb{3}}}
    \proofstepwidestarbreakkeep[2,3]{%
      \IntClass{a}{b}=\IntClass{a'}{b'}}{\rIE{4,7}}
    \proofstepwidestarbreakkeep[2,3]{%
      \IntClass{c}{d}=\IntClass{c'}{d'}}{\rIE{5,8}}
    \proofstepwidestarbreakkeep[1,2,3]{%
      \IntClass{ac+bd}{ad+bc}
        =\IntClass{a'c'+b'd'}{a'd'+b'c'}}{%
      \FormulaRefAuto{IntegerMultiplicationCompatible}{1,10,11}}
    \proofstepwidestarbreakkeep[1,2,3]{u=v}{\rIE{6,9,12}}
  \closeproofpartwideindR

  \proofpartwideindR[Eindeutigkeit beliebiger Ergebniszeugen]{%
    \begin{aligned}[t]
      &u\in\mathbb Z\land\exists a,b,c,d\in\mathbb N\,
        \bigl(z=\IntClass{a}{b}\land w=\IntClass{c}{d}\land{}\\
      &\qquad u=\IntClass{ac+bd}{ad+bc}\bigr)\dsep{}\\
      &v\in\mathbb Z\land\exists a',b',c',d'\in\mathbb N\,
        \bigl(z=\IntClass{a'}{b'}\land w=\IntClass{c'}{d'}\land{}\\
      &\qquad v=\IntClass{a'c'+b'd'}{a'd'+b'c'}\bigr)
        \vdash u=v
    \end{aligned}%
  }{%
    \bigl(u\in\mathbb Z\land
    \exists a,b,c,d\in\mathbb N\,
    (z=\IntClass{a}{b}\land w=\IntClass{c}{d}
     \land u=\IntClass{ac+bd}{ad+bc})\bigr)\dsep
    \bigl(v\in\mathbb Z\land
    \exists a',b',c',d'\in\mathbb N\,
    (z=\IntClass{a'}{b'}\land w=\IntClass{c'}{d'}
     \land v=\IntClass{a'c'+b'd'}{a'd'+b'c'})\bigr)
    \vdash u=v%
  }[IntegerMultiplicationDescentUniqueness]
    \proofstepwidestarbreakkeep[1]{%
      \begin{aligned}[t]
        &u\in\mathbb Z\land\exists a,b,c,d\in\mathbb N\,{}\\
        &(z=\IntClass{a}{b}\land w=\IntClass{c}{d}
          \land u=\IntClass{ac+bd}{ad+bc})
      \end{aligned}}{\rA}
    \proofstepwidestarbreakkeep[2]{%
      \begin{aligned}[t]
        &v\in\mathbb Z\land\exists a',b',c',d'\in\mathbb N\,{}\\
        &(z=\IntClass{a'}{b'}\land w=\IntClass{c'}{d'}\land{}\\
        &\qquad v=\IntClass{a'c'+b'd'}{a'd'+b'c'})
      \end{aligned}}{\rA}
    \proofstepwidestarbreakkeep[1]{%
      \exists a,b,c,d\in\mathbb N\,
      (z=\IntClass{a}{b}\land w=\IntClass{c}{d}
       \land u=\IntClass{ac+bd}{ad+bc})}{\rAEb{1}}
    \proofstepwidestarbreakkeep[2]{%
      \begin{aligned}[t]
        &\exists a',b',c',d'\in\mathbb N\,{}\\
        &(z=\IntClass{a'}{b'}\land w=\IntClass{c'}{d'}\land{}\\
        &\qquad v=\IntClass{a'c'+b'd'}{a'd'+b'c'})
      \end{aligned}}{\rAEb{2}}
    \proofstepwidestarbreakkeep[5]{%
      \begin{aligned}[t]
        &a,b,c,d\in\mathbb N\land z=\IntClass{a}{b}\land{}\\
        &w=\IntClass{c}{d}\land u=\IntClass{ac+bd}{ad+bc}
      \end{aligned}}{\rA}
    \proofstepwidestarbreakkeep[6]{%
      \begin{aligned}[t]
        &a',b',c',d'\in\mathbb N\land z=\IntClass{a'}{b'}\land{}\\
        &w=\IntClass{c'}{d'}\land
          v=\IntClass{a'c'+b'd'}{a'd'+b'c'}
      \end{aligned}}{\rA}
    \proofstepwidestarbreakkeep[5]{a,b,c,d\in\mathbb N}{\rAEa{5}}
    \proofstepwidestarbreakkeep[5]{%
      z=\IntClass{a}{b}\land w=\IntClass{c}{d}
      \land u=\IntClass{ac+bd}{ad+bc}}{\rAEb{5}}
    \proofstepwidestarbreakkeep[6]{a',b',c',d'\in\mathbb N}{\rAEa{6}}
    \proofstepwidestarbreakkeep[6]{%
      z=\IntClass{a'}{b'}\land w=\IntClass{c'}{d'}
      \land v=\IntClass{a'c'+b'd'}{a'd'+b'c'}}{\rAEb{6}}
    \proofstepwidestarbreakkeep[5,6]{%
      a,b,c,d,a',b',c',d'\in\mathbb N}{\rAI{7,9}}
    \proofstepwidestarbreakkeep[5,6]{u=v}{%
      \FormulaRefAuto{IntegerMultiplicationDescentFixedUniqueness}{%
        11,8,10}}
    \proofstepwidestarbreakkeep[5]{u=v}{\rEE{4,6,12}}
    \proofstepwidestarbreakkeep[1,2]{u=v}{\rEE{3,5,13}}
  \closeproofpartwideindR

  \proofpartwideindR[Abschluss des Abstiegs]{%
    \begin{aligned}[t]
      &z,w\in\mathbb Z\vdash
      \exists! u\in\mathbb Z\,
      \exists p\in\mathbb N\,\exists q\in\mathbb N\,
      \exists r\in\mathbb N\,\exists s\in\mathbb N\,{}\\
      &\qquad
      \bigl(z=\IntClass{p}{q}\land w=\IntClass{r}{s}
      \land u=\IntClass{p r+q s}{p s+q r}\bigr)
    \end{aligned}%
  }{%
    z,w\in\mathbb Z\vdash
    \exists! u\in\mathbb Z\,
    \exists p\in\mathbb N\,\exists q\in\mathbb N\,
    \exists r\in\mathbb N\,\exists s\in\mathbb N\,
    \bigl(z=\IntClass{p}{q}\land w=\IntClass{r}{s}
    \land u=\IntClass{p r+q s}{p s+q r}\bigr)%
  }[IntegerMultiplicationDescentConclusion]
    \proofstepwidestarbreakkeep[1]{z,w\in\mathbb Z}{\rA}
    \proofstepwidestarbreakkeep[1]{%
      \begin{aligned}[t]
        &\exists u\in\mathbb Z\,\exists p,q,r,s\in\mathbb N\,{}\\
        &(z=\IntClass{p}{q}\land w=\IntClass{r}{s}\land
          u=\IntClass{p r+q s}{p s+q r})
      \end{aligned}}{%
      \FormulaRefAuto{IntegerMultiplicationDescentExistence}{1}}
    \proofstepwidestarbreakkeep[3]{%
      \begin{aligned}[t]
        &u\in\mathbb Z\land\exists p,q,r,s\in\mathbb N\,{}\\
        &(z=\IntClass{p}{q}\land w=\IntClass{r}{s}\land
          u=\IntClass{p r+q s}{p s+q r})
      \end{aligned}}{\rA}
    \proofstepwidestarbreakkeep[4]{%
      \begin{aligned}[t]
        &v\in\mathbb Z\land\exists p,q,r,s\in\mathbb N\,{}\\
        &(z=\IntClass{p}{q}\land w=\IntClass{r}{s}\land
          v=\IntClass{p r+q s}{p s+q r})
      \end{aligned}}{\rA}
    \proofstepwidestarbreakkeep[3,4]{u=v}{%
      \FormulaRefAuto{IntegerMultiplicationDescentUniqueness}{3,4}}
    \proofstepwidestarbreakkeep[1]{%
      \begin{aligned}[t]
        &\exists! u\in\mathbb Z\,\exists p,q,r,s\in\mathbb N\,{}\\
        &(z=\IntClass{p}{q}\land w=\IntClass{r}{s}\land
          u=\IntClass{p r+q s}{p s+q r})
      \end{aligned}}{\UEI{2,3,4,5}}
  \closeproofpartwideindR
\end{tabproofsplitwide}

Der Quotientenabstieg liefert damit auch für die Multiplikation unabhängig
von beiden gewählten Differenzenpaaren genau ein Ergebnis in \(\mathbb Z\).

\FormulaDefDeltaK[Multiplikation ganzer Zahlen]{%
  a,b,c,d\in\mathbb N\vdash
  \IntClass{a}{b}\cdot\IntClass{c}{d}
  \coloneqq
  \IntClass{a\cdot c+b\cdot d}{a\cdot d+b\cdot c}%
}{IntegerMultiplicationDef}{%
  \DeltaRow{Mengen}{a\dsep b\dsep c\dsep d}%
  \DeltaRow{Neue Symbole}{\cdot}%
}

\FormulaThmDeltaK[Abgeschlossenheit der Multiplikation ganzer Zahlen]{%
  z\in\mathbb Z\dsep w\in\mathbb Z
  \vdash z\cdot w\in\mathbb Z%
}{IntMulClosure}{%
  \DeltaRow{Mengen}{z\dsep w\dsep a\dsep b\dsep c\dsep d}%
}
\begin{tabproofsplitwide}
  \proofpartwideindR[Abgeschlossenheit für feste Repräsentanten]{%
    \begin{aligned}[t]
      &a,b,c,d\in\mathbb N\land
        z=\IntClass{a}{b}\land w=\IntClass{c}{d}\\
      &\vdash z\cdot w\in\mathbb Z
    \end{aligned}
  }{%
    a,b,c,d\in\mathbb N\land
    z=\IntClass{a}{b}\land w=\IntClass{c}{d}
    \vdash z\cdot w\in\mathbb Z%
  }[IntMulClosureForRepresentatives]
    \proofstepwidestarbreakkeep[1]{a,b,c,d\in\mathbb N\land
      z=\IntClass{a}{b}\land w=\IntClass{c}{d}}{\rA}
    \proofstepwidestarbreakkeep[1]{%
      z\cdot w=\IntClass{ac+bd}{ad+bc}}{%
      \FormulaRefAuto{IntegerMultiplicationDef}{1}}
    \proofstepwidestarbreakkeep[1]{ac+bd\in\mathbb N}{%
      \FormulaRefAuto{PeanoAddClosure}{%
        \FormulaRefAuto{PeanoMulClosure}{1},
        \FormulaRefAuto{PeanoMulClosure}{1}}}
    \proofstepwidestarbreakkeep[1]{ad+bc\in\mathbb N}{%
      \FormulaRefAuto{PeanoAddClosure}{%
        \FormulaRefAuto{PeanoMulClosure}{1},
        \FormulaRefAuto{PeanoMulClosure}{1}}}
    \proofstepwidestarbreakkeep[1]{\IntClass{ac+bd}{ad+bc}\in\mathbb Z}{%
      \FormulaRefAuto{IntegerClassInZ}{3,4}}
    \proofstepwidestarbreakkeep[1]{z\cdot w\in\mathbb Z}{\rIE{2,5}}
  \closeproofpartwideindR

  \proofpartwideindR[Elimination der Repräsentanten]{%
    z\in\mathbb Z\dsep w\in\mathbb Z
    \vdash z\cdot w\in\mathbb Z%
  }{%
    z\in\mathbb Z\dsep w\in\mathbb Z
    \vdash z\cdot w\in\mathbb Z%
  }[IntMulClosureRepresentativeElimination]
    \proofstepwidestarbreakkeep[1]{z\in\mathbb Z}{\rA}
    \proofstepwidestarbreakkeep[2]{w\in\mathbb Z}{\rA}
    \proofstepwidestarbreakkeep[1]{\exists a,b\in\mathbb N\,
      z=\IntClass{a}{b}}{\FormulaRefAuto{IntegerPairRepresentative}{1}}
    \proofstepwidestarbreakkeep[2]{\exists c,d\in\mathbb N\,
      w=\IntClass{c}{d}}{\FormulaRefAuto{IntegerPairRepresentative}{2}}
    \proofstepwidestarbreakkeep[5]{a,b,c,d\in\mathbb N\land
      z=\IntClass{a}{b}\land w=\IntClass{c}{d}}{\rA}
    \proofstepwidestarbreakkeep[5]{z\cdot w\in\mathbb Z}{%
      \FormulaRefAuto{IntMulClosureForRepresentatives}{5}}
    \proofstepwidestarbreakkeep[1,2]{z\cdot w\in\mathbb Z}{\rEE{3,4,5,6}}
  \closeproofpartwideindR
\end{tabproofsplitwide}

\FormulaThmDeltaK[Kommutativität der Multiplikation ganzer Zahlen]{%
  z,w\in\mathbb Z\vdash z\cdot w=w\cdot z%
}{IntMulCommutative}{%
  \DeltaRow{Mengen}{z\dsep w\dsep a\dsep b\dsep c\dsep d}%
}
\begin{tabproofsplitwide}
  \proofpartwideindR[Produkte gewählter Repräsentanten]{%
    \begin{aligned}[t]
      &a,b,c,d\in\mathbb N\land
        z=\IntClass{a}{b}\land w=\IntClass{c}{d}\\
      &\vdash
        z\cdot w=\IntClass{ac+bd}{ad+bc}\land
        w\cdot z=\IntClass{ca+db}{cb+da}
    \end{aligned}
  }{%
    a,b,c,d\in\mathbb N\land
    z=\IntClass{a}{b}\land w=\IntClass{c}{d}
    \vdash
    z\cdot w=\IntClass{ac+bd}{ad+bc}\land
    w\cdot z=\IntClass{ca+db}{cb+da}
  }[IntMulCommutativeRepresentativeProducts]
    \proofstepwidestar[1]{a,b,c,d\in\mathbb N\land
      z=\IntClass{a}{b}\land w=\IntClass{c}{d}}{\rA}
    \proofstepwidestar[1]{z\cdot w=\IntClass{ac+bd}{ad+bc}}{%
      \FormulaRefAuto{IntegerMultiplicationDef}{1}}
    \proofstepwidestar[1]{w\cdot z=\IntClass{ca+db}{cb+da}}{%
      \FormulaRefAuto{IntegerMultiplicationDef}{1}}
    \proofstepwidestar[1]{%
      z\cdot w=\IntClass{ac+bd}{ad+bc}\land
      w\cdot z=\IntClass{ca+db}{cb+da}}{\rAI{2,3}}
  \closeproofpartwideindR

  \proofpartwideindR[Vertauschung der Klassenkoordinaten]{%
    a,b,c,d\in\mathbb N\vdash
    \IntClass{ac+bd}{ad+bc}=\IntClass{ca+db}{cb+da}%
  }{%
    a,b,c,d\in\mathbb N\vdash
    \IntClass{ac+bd}{ad+bc}=\IntClass{ca+db}{cb+da}%
  }[IntMulCommutativeClassCoordinates]
    \proofstepwidestar[1]{a,b,c,d\in\mathbb N}{\rA}
    \proofstepwidestar[1]{ac=ca}{%
      \FormulaRefAuto{PeanoMulCommutative}{1}}
    \proofstepwidestar[1]{bd=db}{%
      \FormulaRefAuto{PeanoMulCommutative}{1}}
    \proofstepwidestar[1]{ad=da}{%
      \FormulaRefAuto{PeanoMulCommutative}{1}}
    \proofstepwidestar[1]{bc=cb}{%
      \FormulaRefAuto{PeanoMulCommutative}{1}}
    \proofstepwidestar[1]{ac+bd=ca+db}{\rIE{2,3}}
    \proofstepwidestar[1]{ad+bc=cb+da}{%
      \FormulaRefAuto{PeanoAddCommutative}{\rIE{4,5}}}
    \proofstepwidestar[1]{\IntClass{ac+bd}{ad+bc}=
      \IntClass{ca+db}{cb+da}}{\rIE{6,7}}
  \closeproofpartwideindR

  \proofpartwideindR[Elimination der Repräsentanten]{%
    z,w\in\mathbb Z\vdash z\cdot w=w\cdot z%
  }{%
    z,w\in\mathbb Z\vdash z\cdot w=w\cdot z%
  }[IntMulCommutativeRepresentativeElimination]
    \proofstepwidestar[1]{z,w\in\mathbb Z}{\rA}
    \proofstepwidestar[1]{z\in\mathbb Z}{\rAEa{1}}
    \proofstepwidestar[1]{w\in\mathbb Z}{\rAEb{1}}
    \proofstepwidestar[1]{\exists a,b\in\mathbb N\,
      z=\IntClass{a}{b}}{\FormulaRefAuto{IntegerPairRepresentative}{2}}
    \proofstepwidestar[1]{\exists c,d\in\mathbb N\,
      w=\IntClass{c}{d}}{\FormulaRefAuto{IntegerPairRepresentative}{3}}
    \proofstepwidestar[6]{a,b,c,d\in\mathbb N\land
      z=\IntClass{a}{b}\land w=\IntClass{c}{d}}{\rA}
    \proofstepwidestar[6]{%
      z\cdot w=\IntClass{ac+bd}{ad+bc}\land
      w\cdot z=\IntClass{ca+db}{cb+da}}{%
      \FormulaRefAuto{IntMulCommutativeRepresentativeProducts}{6}}
    \proofstepwide[6]{z\cdot w}{=}{\IntClass{ac+bd}{ad+bc}}{\rAEa{7}}
    \proofstepwide[6]{}{=}{\IntClass{ca+db}{cb+da}}{%
      \FormulaRefAuto{IntMulCommutativeClassCoordinates}{6}}
    \proofstepwide[6]{}{=}{w\cdot z}{%
      \FormulaRefAuto{a=b\vdash b=a}{\rAEb{7}}}
    \proofstepwidestar[6]{z\cdot w=w\cdot z}{\rChain{8,9,10}}
    \proofstepwidestar[1]{z\cdot w=w\cdot z}{\rEE{4,5,6,11}}
  \closeproofpartwideindR
\end{tabproofsplitwide}

\FormulaThmDeltaK[Einheitsgesetze der Multiplikation ganzer Zahlen]{%
  z\in\mathbb Z\vdash
  z\cdot\IntOne=z\land\IntOne\cdot z=z%
}{IntMulOne}{%
  \DeltaRow{Mengen}{z\dsep a\dsep b}%
}
\begin{tabproofsplitwide}
  \proofpartwideindR[Rechte Einheit am Repräsentanten]{%
    a,b\in\mathbb N\land z=\IntClass{a}{b}
    \vdash z\cdot\IntOne=z%
  }{%
    a,b\in\mathbb N\land z=\IntClass{a}{b}
    \vdash z\cdot\IntOne=z%
  }[IntMulRightIdentityRepresentative]
    \proofstepwidestar[1]{%
      a,b\in\mathbb N\land z=\IntClass{a}{b}}{\rA}
    \proofstepwidestar[]{\IntOne=\IntEmbed(1)}{%
      \FormulaRefAuto{IntegerOneDef}}
    \proofstepwidestar[]{\IntEmbed(1)=\IntClass{1}{0}}{%
      \FormulaRefAuto{NaturalIntegerEmbeddingValue}[thm]{%
        \FormulaRefAuto{PeanoOneInN}}}
    \proofstepwidestar[1]{a\cdot1=a}{%
      \FormulaRefAuto{PeanoMulRightOne}{1}}
    \proofstepwidestar[1]{b\cdot1=b}{%
      \FormulaRefAuto{PeanoMulRightOne}{1}}
    \proofstepwidestar[1]{a\cdot0=0}{%
      \FormulaRefAuto{PeanoMulRightZero}{1}}
    \proofstepwidestar[1]{b\cdot0=0}{%
      \FormulaRefAuto{PeanoMulRightZero}{1}}
    \proofstepwide[1]{a\cdot1+b\cdot0}{=}{a+0}{\rIE{4,7}}
    \proofstepwide[1]{}{=}{a}{\FormulaRefAuto{PeanoAddRightZero}{1}}
    \proofstepwidestar[1]{a\cdot1+b\cdot0=a}{\rChain{8,9}}
    \proofstepwide[1]{a\cdot0+b\cdot1}{=}{0+b}{\rIE{6,5}}
    \proofstepwide[1]{}{=}{b}{\FormulaRefAuto{PeanoAddLeftZero}{1}}
    \proofstepwidestar[1]{a\cdot0+b\cdot1=b}{\rChain{11,12}}
    \proofstepwide[1]{z\cdot\IntOne}{=}{%
      \IntClass{a\cdot1+b\cdot0}{a\cdot0+b\cdot1}}{%
      \FormulaRefAuto{IntegerMultiplicationDef}{1,2,3}}
    \proofstepwide[1]{}{=}{\IntClass{a}{b}}{\rIE{10,13}}
    \proofstepwide[1]{}{=}{z}{%
      \FormulaRefAuto{a=b\vdash b=a}{\rAEb{1}}}
    \proofstepwidestar[1]{z\cdot\IntOne=z}{\rChain{14,15,16}}
  \closeproofpartwideindR

  \proofpartwideindR[Beidseitige Einheit]{%
    z\in\mathbb Z\vdash
    z\cdot\IntOne=z\land\IntOne\cdot z=z%
  }{%
    z\in\mathbb Z\vdash
    z\cdot\IntOne=z\land\IntOne\cdot z=z%
  }[IntMulTwoSidedIdentityConclusion]
    \proofstepwidestar[1]{z\in\mathbb Z}{\rA}
    \proofstepwidestar[1]{\exists a,b\in\mathbb N\,
      z=\IntClass{a}{b}}{%
      \FormulaRefAuto{IntegerPairRepresentative}{1}}
    \proofstepwidestar[3]{%
      a,b\in\mathbb N\land z=\IntClass{a}{b}}{\rA}
    \proofstepwidestar[3]{z\cdot\IntOne=z}{%
      \FormulaRefAuto{IntMulRightIdentityRepresentative}{3}}
    \proofstepwidestar[3]{\IntOne\cdot z=z}{%
      \FormulaRefAuto{IntMulCommutative}{%
        \FormulaRefAuto{IntOneCarrier},1,
        \FormulaRefAuto{a=b\vdash b=a}{4}}}
    \proofstepwidestar[3]{z\cdot\IntOne=z\land
      \IntOne\cdot z=z}{\rAI{4,5}}
    \proofstepwidestar[1]{z\cdot\IntOne=z\land
      \IntOne\cdot z=z}{\rEE{2,3,6}}
  \closeproofpartwideindR
\end{tabproofsplitwide}

\FormulaThmDeltaK[Assoziativität der Multiplikation ganzer Zahlen]{%
  z,w,u\in\mathbb Z\vdash
  (z\cdot w)\cdot u=z\cdot(w\cdot u)%
}{IntMulAssociative}{%
  \DeltaRow{Mengen}{z\dsep w\dsep u\dsep
    a\dsep b\dsep c\dsep d\dsep e\dsep f}%
}
\begin{tabproofsplitwide}
  \proofpartwideindR[Positive Koordinate: linke Seite ausmultiplizieren]{%
    \begin{aligned}[t]
      &a,b,c,d,e,f\in\mathbb N\vdash{}\\
      &(ac+bd)e+(ad+bc)f
      =\bigl((ac)e+(bd)e\bigr)+\bigl((ad)f+(bc)f\bigr)
    \end{aligned}
  }{%
    a,b,c,d,e,f\in\mathbb N\vdash
    (ac+bd)e+(ad+bc)f
    =\bigl((ac)e+(bd)e\bigr)+\bigl((ad)f+(bc)f\bigr)
  }[IntMulAssociativePositiveLeftExpansion]
    \proofstepwidestarbreakkeep[1]{a,b,c,d,e,f\in\mathbb N}{\rA}
    \proofstepwide[1]{(ac+bd)e}{=}{(ac)e+(bd)e}{%
      \FormulaRefAuto{PeanoMulRightDistributive}{1}}
    \proofstepwide[1]{(ad+bc)f}{=}{(ad)f+(bc)f}{%
      \FormulaRefAuto{PeanoMulRightDistributive}{1}}
    \proofstepwidestarbreakkeep[1]{%
      (ac+bd)e+(ad+bc)f
      =\bigl((ac)e+(bd)e\bigr)+\bigl((ad)f+(bc)f\bigr)}{\rIE{2,3}}
  \closeproofpartwideindR

  \proofpartwideindR[Positive Koordinate: rechte Seite ausmultiplizieren]{%
    \begin{aligned}[t]
      &a,b,c,d,e,f\in\mathbb N\vdash{}\\
      &a(ce+df)+b(cf+de)
      =\bigl(a(ce)+a(df)\bigr)+\bigl(b(cf)+b(de)\bigr)
    \end{aligned}
  }{%
    a,b,c,d,e,f\in\mathbb N\vdash
    a(ce+df)+b(cf+de)
    =\bigl(a(ce)+a(df)\bigr)+\bigl(b(cf)+b(de)\bigr)
  }[IntMulAssociativePositiveRightExpansion]
    \proofstepwidestarbreakkeep[1]{a,b,c,d,e,f\in\mathbb N}{\rA}
    \proofstepwide[1]{a(ce+df)}{=}{a(ce)+a(df)}{%
      \FormulaRefAuto{PeanoMulLeftDistributive}{1}}
    \proofstepwide[1]{b(cf+de)}{=}{b(cf)+b(de)}{%
      \FormulaRefAuto{PeanoMulLeftDistributive}{1}}
    \proofstepwidestarbreakkeep[1]{%
      a(ce+df)+b(cf+de)
      =\bigl(a(ce)+a(df)\bigr)+\bigl(b(cf)+b(de)\bigr)}{\rIE{2,3}}
  \closeproofpartwideindR

  \proofpartwideindR[Positive Koordinate: ausmultiplizierte Terme ordnen]{%
    \begin{aligned}[t]
      &a,b,c,d,e,f\in\mathbb N\vdash{}\\
      &\bigl((ac)e+(bd)e\bigr)+\bigl((ad)f+(bc)f\bigr)\\
      &\qquad=
      \bigl(a(ce)+a(df)\bigr)+\bigl(b(cf)+b(de)\bigr)
    \end{aligned}
  }{%
    \begin{aligned}[t]
      &a,b,c,d,e,f\in\mathbb N\vdash{}\\
      &\bigl((ac)e+(bd)e\bigr)+\bigl((ad)f+(bc)f\bigr)
      =
      \bigl(a(ce)+a(df)\bigr)+\bigl(b(cf)+b(de)\bigr)
    \end{aligned}
  }[IntMulAssociativePositiveExpandedComparison]
    \proofstepwidestarbreakkeep[1]{a,b,c,d,e,f\in\mathbb N}{\rA}
    \proofstepwidestarbreakkeep[1]{(ac)e=a(ce)}{%
      \FormulaRefAuto{PeanoMulAssociative}{1}}
    \proofstepwidestarbreakkeep[1]{(bd)e=b(de)}{%
      \FormulaRefAuto{PeanoMulAssociative}{1}}
    \proofstepwidestarbreakkeep[1]{(ad)f=a(df)}{%
      \FormulaRefAuto{PeanoMulAssociative}{1}}
    \proofstepwidestarbreakkeep[1]{(bc)f=b(cf)}{%
      \FormulaRefAuto{PeanoMulAssociative}{1}}
    \proofstepwidestarbreakkeep[1]{%
      \bigl((ac)e+(bd)e\bigr)+\bigl((ad)f+(bc)f\bigr)
      =
      \bigl(a(ce)+b(de)\bigr)+\bigl(a(df)+b(cf)\bigr)}{%
      \rIE{2,3,4,5}}
    \proofstepwidestarbreakkeep[1]{%
      \bigl(a(ce)+b(de)\bigr)+\bigl(a(df)+b(cf)\bigr)
      =
      \bigl(a(ce)+a(df)\bigr)+\bigl(b(de)+b(cf)\bigr)}{%
      \FormulaRefAuto{PeanoAddInterchange}{1}}
    \proofstepwidestarbreakkeep[1]{b(de)+b(cf)=b(cf)+b(de)}{%
      \FormulaRefAuto{PeanoAddCommutative}{1}}
    \proofstepwidestarbreakkeep[1]{%
      \bigl(a(ce)+a(df)\bigr)+\bigl(b(de)+b(cf)\bigr)
      =
      \bigl(a(ce)+a(df)\bigr)+\bigl(b(cf)+b(de)\bigr)}{\rIE{8}}
    \proofstepwidestarbreakkeep[1]{%
      \bigl((ac)e+(bd)e\bigr)+\bigl((ad)f+(bc)f\bigr)
      =
      \bigl(a(ce)+a(df)\bigr)+\bigl(b(cf)+b(de)\bigr)}{%
      \rChain{6,7,9}}
  \closeproofpartwideindR

  \proofpartwideindR[Gleichheit der positiven Koordinaten]{%
    \begin{aligned}[t]
      &a,b,c,d,e,f\in\mathbb N\vdash{}\\
      &(ac+bd)e+(ad+bc)f
      =a(ce+df)+b(cf+de)
    \end{aligned}
  }{%
    a,b,c,d,e,f\in\mathbb N\vdash
    (ac+bd)e+(ad+bc)f
    =a(ce+df)+b(cf+de)
  }[IntMulAssociativePositiveCoordinate]
    \proofstepwidestarbreakkeep[1]{a,b,c,d,e,f\in\mathbb N}{\rA}
    \proofstepwidestarbreakkeep[1]{%
      (ac+bd)e+(ad+bc)f
      =\bigl((ac)e+(bd)e\bigr)+\bigl((ad)f+(bc)f\bigr)}{%
      \FormulaRefAuto{IntMulAssociativePositiveLeftExpansion}{1}}
    \proofstepwidestarbreakkeep[1]{%
      \bigl((ac)e+(bd)e\bigr)+\bigl((ad)f+(bc)f\bigr)
      =
      \bigl(a(ce)+a(df)\bigr)+\bigl(b(cf)+b(de)\bigr)}{%
      \FormulaRefAuto{IntMulAssociativePositiveExpandedComparison}{1}}
    \proofstepwidestarbreakkeep[1]{%
      a(ce+df)+b(cf+de)
      =\bigl(a(ce)+a(df)\bigr)+\bigl(b(cf)+b(de)\bigr)}{%
      \FormulaRefAuto{IntMulAssociativePositiveRightExpansion}{1}}
    \proofstepwidestarbreakkeep[1]{%
      \bigl(a(ce)+a(df)\bigr)+\bigl(b(cf)+b(de)\bigr)
      =a(ce+df)+b(cf+de)}{%
      \FormulaRefAuto{a=b\vdash b=a}{4}}
    \proofstepwidestarbreakkeep[1]{%
      (ac+bd)e+(ad+bc)f
      =a(ce+df)+b(cf+de)}{%
      \rChain{2,3,5}}
  \closeproofpartwideindR

  \proofpartwideindR[Negative Koordinate: linke Seite ausmultiplizieren]{%
    \begin{aligned}[t]
      &a,b,c,d,e,f\in\mathbb N\vdash{}\\
      &(ac+bd)f+(ad+bc)e
      =\bigl((ac)f+(bd)f\bigr)+\bigl((ad)e+(bc)e\bigr)
    \end{aligned}
  }{%
    a,b,c,d,e,f\in\mathbb N\vdash
    (ac+bd)f+(ad+bc)e
    =\bigl((ac)f+(bd)f\bigr)+\bigl((ad)e+(bc)e\bigr)
  }[IntMulAssociativeNegativeLeftExpansion]
    \proofstepwidestarbreakkeep[1]{a,b,c,d,e,f\in\mathbb N}{\rA}
    \proofstepwide[1]{(ac+bd)f}{=}{(ac)f+(bd)f}{%
      \FormulaRefAuto{PeanoMulRightDistributive}{1}}
    \proofstepwide[1]{(ad+bc)e}{=}{(ad)e+(bc)e}{%
      \FormulaRefAuto{PeanoMulRightDistributive}{1}}
    \proofstepwidestarbreakkeep[1]{%
      (ac+bd)f+(ad+bc)e
      =\bigl((ac)f+(bd)f\bigr)+\bigl((ad)e+(bc)e\bigr)}{\rIE{2,3}}
  \closeproofpartwideindR

  \proofpartwideindR[Negative Koordinate: rechte Seite ausmultiplizieren]{%
    \begin{aligned}[t]
      &a,b,c,d,e,f\in\mathbb N\vdash{}\\
      &a(cf+de)+b(ce+df)
      =\bigl(a(cf)+a(de)\bigr)+\bigl(b(ce)+b(df)\bigr)
    \end{aligned}
  }{%
    a,b,c,d,e,f\in\mathbb N\vdash
    a(cf+de)+b(ce+df)
    =\bigl(a(cf)+a(de)\bigr)+\bigl(b(ce)+b(df)\bigr)
  }[IntMulAssociativeNegativeRightExpansion]
    \proofstepwidestarbreakkeep[1]{a,b,c,d,e,f\in\mathbb N}{\rA}
    \proofstepwide[1]{a(cf+de)}{=}{a(cf)+a(de)}{%
      \FormulaRefAuto{PeanoMulLeftDistributive}{1}}
    \proofstepwide[1]{b(ce+df)}{=}{b(ce)+b(df)}{%
      \FormulaRefAuto{PeanoMulLeftDistributive}{1}}
    \proofstepwidestarbreakkeep[1]{%
      a(cf+de)+b(ce+df)
      =\bigl(a(cf)+a(de)\bigr)+\bigl(b(ce)+b(df)\bigr)}{\rIE{2,3}}
  \closeproofpartwideindR

  \proofpartwideindR[Negative Koordinate: ausmultiplizierte Terme ordnen]{%
    \begin{aligned}[t]
      &a,b,c,d,e,f\in\mathbb N\vdash{}\\
      &\bigl((ac)f+(bd)f\bigr)+\bigl((ad)e+(bc)e\bigr)\\
      &\qquad=
      \bigl(a(cf)+a(de)\bigr)+\bigl(b(ce)+b(df)\bigr)
    \end{aligned}
  }{%
    \begin{aligned}[t]
      &a,b,c,d,e,f\in\mathbb N\vdash{}\\
      &\bigl((ac)f+(bd)f\bigr)+\bigl((ad)e+(bc)e\bigr)
      =
      \bigl(a(cf)+a(de)\bigr)+\bigl(b(ce)+b(df)\bigr)
    \end{aligned}
  }[IntMulAssociativeNegativeExpandedComparison]
    \proofstepwidestarbreakkeep[1]{a,b,c,d,e,f\in\mathbb N}{\rA}
    \proofstepwidestarbreakkeep[1]{(ac)f=a(cf)}{%
      \FormulaRefAuto{PeanoMulAssociative}{1}}
    \proofstepwidestarbreakkeep[1]{(bd)f=b(df)}{%
      \FormulaRefAuto{PeanoMulAssociative}{1}}
    \proofstepwidestarbreakkeep[1]{(ad)e=a(de)}{%
      \FormulaRefAuto{PeanoMulAssociative}{1}}
    \proofstepwidestarbreakkeep[1]{(bc)e=b(ce)}{%
      \FormulaRefAuto{PeanoMulAssociative}{1}}
    \proofstepwidestarbreakkeep[1]{%
      \bigl((ac)f+(bd)f\bigr)+\bigl((ad)e+(bc)e\bigr)
      =
      \bigl(a(cf)+b(df)\bigr)+\bigl(a(de)+b(ce)\bigr)}{%
      \rIE{2,3,4,5}}
    \proofstepwidestarbreakkeep[1]{%
      \bigl(a(cf)+b(df)\bigr)+\bigl(a(de)+b(ce)\bigr)
      =
      \bigl(a(cf)+a(de)\bigr)+\bigl(b(df)+b(ce)\bigr)}{%
      \FormulaRefAuto{PeanoAddInterchange}{1}}
    \proofstepwidestarbreakkeep[1]{b(df)+b(ce)=b(ce)+b(df)}{%
      \FormulaRefAuto{PeanoAddCommutative}{1}}
    \proofstepwidestarbreakkeep[1]{%
      \bigl(a(cf)+a(de)\bigr)+\bigl(b(df)+b(ce)\bigr)
      =
      \bigl(a(cf)+a(de)\bigr)+\bigl(b(ce)+b(df)\bigr)}{\rIE{8}}
    \proofstepwidestarbreakkeep[1]{%
      \bigl((ac)f+(bd)f\bigr)+\bigl((ad)e+(bc)e\bigr)
      =
      \bigl(a(cf)+a(de)\bigr)+\bigl(b(ce)+b(df)\bigr)}{%
      \rChain{6,7,9}}
  \closeproofpartwideindR

  \proofpartwideindR[Gleichheit der negativen Koordinaten]{%
    \begin{aligned}[t]
      &a,b,c,d,e,f\in\mathbb N\vdash{}\\
      &(ac+bd)f+(ad+bc)e
      =a(cf+de)+b(ce+df)
    \end{aligned}
  }{%
    a,b,c,d,e,f\in\mathbb N\vdash
    (ac+bd)f+(ad+bc)e
    =a(cf+de)+b(ce+df)
  }[IntMulAssociativeNegativeCoordinate]
    \proofstepwidestarbreakkeep[1]{a,b,c,d,e,f\in\mathbb N}{\rA}
    \proofstepwidestarbreakkeep[1]{%
      (ac+bd)f+(ad+bc)e
      =\bigl((ac)f+(bd)f\bigr)+\bigl((ad)e+(bc)e\bigr)}{%
      \FormulaRefAuto{IntMulAssociativeNegativeLeftExpansion}{1}}
    \proofstepwidestarbreakkeep[1]{%
      \bigl((ac)f+(bd)f\bigr)+\bigl((ad)e+(bc)e\bigr)
      =
      \bigl(a(cf)+a(de)\bigr)+\bigl(b(ce)+b(df)\bigr)}{%
      \FormulaRefAuto{IntMulAssociativeNegativeExpandedComparison}{1}}
    \proofstepwidestarbreakkeep[1]{%
      a(cf+de)+b(ce+df)
      =\bigl(a(cf)+a(de)\bigr)+\bigl(b(ce)+b(df)\bigr)}{%
      \FormulaRefAuto{IntMulAssociativeNegativeRightExpansion}{1}}
    \proofstepwidestarbreakkeep[1]{%
      \bigl(a(cf)+a(de)\bigr)+\bigl(b(ce)+b(df)\bigr)
      =a(cf+de)+b(ce+df)}{%
      \FormulaRefAuto{a=b\vdash b=a}{4}}
    \proofstepwidestarbreakkeep[1]{%
      (ac+bd)f+(ad+bc)e
      =a(cf+de)+b(ce+df)}{%
      \rChain{2,3,5}}
  \closeproofpartwideindR

  \proofpartwideindR[Rechnung mit Klassen]{%
    \begin{aligned}[t]
      &a,b,c,d,e,f\in\mathbb N\vdash{}\\
      &(\IntClass{a}{b}\cdot\IntClass{c}{d})\cdot\IntClass{e}{f}
      =\IntClass{a}{b}\cdot
       (\IntClass{c}{d}\cdot\IntClass{e}{f})
    \end{aligned}
  }{%
    a,b,c,d,e,f\in\mathbb N\vdash
    (\IntClass{a}{b}\cdot\IntClass{c}{d})\cdot\IntClass{e}{f}
    =\IntClass{a}{b}\cdot
     (\IntClass{c}{d}\cdot\IntClass{e}{f})
  }[IntMulAssociativeOnClasses]
    \proofstepwidestarbreakkeep[1]{a,b,c,d,e,f\in\mathbb N}{\rA}
    \proofstepwidestarbreakkeep[1]{%
      \begin{aligned}[t]
        &\IntClass{a}{b}\cdot
          (\IntClass{c}{d}\cdot\IntClass{e}{f})={}\\
        &\IntClass{a(ce+df)+b(cf+de)}
                 {a(cf+de)+b(ce+df)}
      \end{aligned}}{%
      \FormulaRefAuto{IntegerMultiplicationDef}{1}}
    \proofstepwidestarbreakkeep[1]{%
      (ac+bd)e+(ad+bc)f=a(ce+df)+b(cf+de)}{%
      \FormulaRefAuto{IntMulAssociativePositiveCoordinate}{1}}
    \proofstepwidestarbreakkeep[1]{%
      (ac+bd)f+(ad+bc)e=a(cf+de)+b(ce+df)}{%
      \FormulaRefAuto{IntMulAssociativeNegativeCoordinate}{1}}
    \proofstepwidestarbreakkeep[1]{%
      \begin{aligned}[t]
        &(\IntClass{a}{b}\cdot\IntClass{c}{d})
          \cdot\IntClass{e}{f}={}\\
        &\IntClass{(ac+bd)e+(ad+bc)f}
                 {(ac+bd)f+(ad+bc)e}
      \end{aligned}}{%
      \FormulaRefAuto{IntegerMultiplicationDef}{1}}
    \proofstepwidestarbreakkeep[1]{%
      \begin{aligned}[t]
        &\IntClass{(ac+bd)e+(ad+bc)f}
                  {(ac+bd)f+(ad+bc)e}={}\\
        &\IntClass{a(ce+df)+b(cf+de)}
                  {a(cf+de)+b(ce+df)}
      \end{aligned}}{\rIE{3,4}}
    \proofstepwidestarbreakkeep[1]{%
      \begin{aligned}[t]
        &\IntClass{a(ce+df)+b(cf+de)}
                 {a(cf+de)+b(ce+df)}={}\\
        &\IntClass{a}{b}\cdot
          (\IntClass{c}{d}\cdot\IntClass{e}{f})
      \end{aligned}}{%
      \FormulaRefAuto{a=b\vdash b=a}{2}}
    \proofstepwidestarbreakkeep[1]{%
      \begin{aligned}[t]
        &(\IntClass{a}{b}\cdot\IntClass{c}{d})
          \cdot\IntClass{e}{f}={}\\
        &\IntClass{a}{b}\cdot
          (\IntClass{c}{d}\cdot\IntClass{e}{f})
      \end{aligned}}{\rChain{5,6,7}}
  \closeproofpartwideindR

  \proofpartwide{Schluss}
    \proofstepwidestarbreakkeep[1]{z,w,u\in\mathbb Z}{\rA}
    \proofstepwidestarbreakkeep[1]{z\in\mathbb Z}{\rAEa{1}}
    \proofstepwidestarbreakkeep[1]{w\in\mathbb Z}{\rAEa{\rAEb{1}}}
    \proofstepwidestarbreakkeep[1]{u\in\mathbb Z}{\rAEb{\rAEb{1}}}
    \proofstepwidestarbreakkeep[1]{\exists a,b\in\mathbb N\,
      z=\IntClass{a}{b}}{\FormulaRefAuto{IntegerPairRepresentative}{2}}
    \proofstepwidestarbreakkeep[1]{\exists c,d\in\mathbb N\,
      w=\IntClass{c}{d}}{\FormulaRefAuto{IntegerPairRepresentative}{3}}
    \proofstepwidestarbreakkeep[1]{\exists e,f\in\mathbb N\,
      u=\IntClass{e}{f}}{\FormulaRefAuto{IntegerPairRepresentative}{4}}
    \proofstepwidestarbreakkeep[8]{a,b,c,d,e,f\in\mathbb N\land
      z=\IntClass{a}{b}\land w=\IntClass{c}{d}\land
      u=\IntClass{e}{f}}{\rA}
    \proofstepwidestarbreakkeep[8]{%
      \begin{aligned}[t]
        &(\IntClass{a}{b}\cdot\IntClass{c}{d})
          \cdot\IntClass{e}{f}={}\\
        &\IntClass{a}{b}\cdot
          (\IntClass{c}{d}\cdot\IntClass{e}{f})
      \end{aligned}}{%
      \FormulaRefAuto{IntMulAssociativeOnClasses}{8}}
    \proofstepwidestarbreakkeep[8]{%
      (z\cdot w)\cdot u=z\cdot(w\cdot u)}{\rIE{8,9}}
    \proofstepwidestarbreakkeep[1]{(z\cdot w)\cdot u=z\cdot(w\cdot u)}{%
      \rEE{5,6,7,8,10}}
  \closeproofpartwide
\end{tabproofsplitwide}

\FormulaThmDeltaK[Linksdistributivität ganzer Zahlen]{%
  z,w,u\in\mathbb Z\vdash
  z\cdot(w+u)=z\cdot w+z\cdot u%
}{IntMulLeftDistributive}{%
  \DeltaRow{Mengen}{z\dsep w\dsep u\dsep
    a\dsep b\dsep c\dsep d\dsep e\dsep f}%
}
\begin{tabproofsplitwide}
  \proofpartwideindR[Positive Koordinate distributiv ausmultiplizieren]{%
    \begin{aligned}[t]
      &a,b,c,d,e,f\in\mathbb N\vdash{}\\
      &a(c+e)+b(d+f)=(ac+bd)+(ae+bf)
    \end{aligned}
  }{%
    a,b,c,d,e,f\in\mathbb N\vdash
    a(c+e)+b(d+f)=(ac+bd)+(ae+bf)
  }[IntMulLeftDistributivePositiveCoordinate]
    \proofstepwidestarbreakkeep[1]{a,b,c,d,e,f\in\mathbb N}{\rA}
    \proofstepwidestarbreakkeep[1]{%
      a(c+e)+b(d+f)=(ac+bd)+(ae+bf)}{%
      \FormulaRefAuto{PeanoBilinearAddInterchange}{1}}
  \closeproofpartwideindR

  \proofpartwideindR[Negative Koordinate distributiv ausmultiplizieren]{%
    \begin{aligned}[t]
      &a,b,c,d,e,f\in\mathbb N\vdash{}\\
      &a(d+f)+b(c+e)=(ad+bc)+(af+be)
    \end{aligned}
  }{%
    a,b,c,d,e,f\in\mathbb N\vdash
    a(d+f)+b(c+e)=(ad+bc)+(af+be)
  }[IntMulLeftDistributiveNegativeCoordinate]
    \proofstepwidestarbreakkeep[1]{a,b,c,d,e,f\in\mathbb N}{\rA}
    \proofstepwidestarbreakkeep[1]{%
      a(d+f)+b(c+e)=(ad+bc)+(af+be)}{%
      \FormulaRefAuto{PeanoBilinearAddInterchange}{1}}
  \closeproofpartwideindR

  \proofpartwideindR[Distributivrechnung auf Klassen]{%
    \begin{aligned}[t]
      &a,b,c,d,e,f\in\mathbb N\vdash{}\\
      &\IntClass{a}{b}\cdot
       (\IntClass{c}{d}+\IntClass{e}{f})
       =
       \IntClass{a}{b}\cdot\IntClass{c}{d}
       +\IntClass{a}{b}\cdot\IntClass{e}{f}
    \end{aligned}
  }{%
    a,b,c,d,e,f\in\mathbb N\vdash
    \IntClass{a}{b}\cdot
    (\IntClass{c}{d}+\IntClass{e}{f})
    =
    \IntClass{a}{b}\cdot\IntClass{c}{d}
    +\IntClass{a}{b}\cdot\IntClass{e}{f}
  }[IntMulLeftDistributiveOnClasses]
    \proofstepwidestarbreakkeep[1]{a,b,c,d,e,f\in\mathbb N}{\rA}
    \proofstepwidestarbreakkeep[1]{%
      \begin{aligned}[t]
        &\IntClass{a}{b}\cdot\IntClass{c}{d}
          +\IntClass{a}{b}\cdot\IntClass{e}{f}={}\\
        &\IntClass{(ac+bd)+(ae+bf)}
                 {(ad+bc)+(af+be)}
      \end{aligned}}{%
      \FormulaRefAuto{IntegerMultiplicationDef}{%
        \FormulaRefAuto{IntegerAdditionDef}{1}}}
    \proofstepwidestarbreakkeep[1]{%
      a(c+e)+b(d+f)=(ac+bd)+(ae+bf)}{%
      \FormulaRefAuto{IntMulLeftDistributivePositiveCoordinate}{1}}
    \proofstepwidestarbreakkeep[1]{%
      a(d+f)+b(c+e)=(ad+bc)+(af+be)}{%
      \FormulaRefAuto{IntMulLeftDistributiveNegativeCoordinate}{1}}
    \proofstepwidestarbreakkeep[1]{%
      \begin{aligned}[t]
        &\IntClass{a}{b}\cdot
          (\IntClass{c}{d}+\IntClass{e}{f})={}\\
        &\IntClass{a(c+e)+b(d+f)}
                 {a(d+f)+b(c+e)}
      \end{aligned}}{%
      \FormulaRefAuto{IntegerAdditionDef}{%
        \FormulaRefAuto{IntegerMultiplicationDef}{1}}}
    \proofstepwidestarbreakkeep[1]{%
      \IntClass{a(c+e)+b(d+f)}
               {a(d+f)+b(c+e)}
      =\IntClass{(ac+bd)+(ae+bf)}
               {(ad+bc)+(af+be)}}{\rIE{3,4}}
    \proofstepwidestarbreakkeep[1]{%
      \begin{aligned}[t]
        &\IntClass{(ac+bd)+(ae+bf)}
                 {(ad+bc)+(af+be)}={}\\
        &\IntClass{a}{b}\cdot\IntClass{c}{d}
          +\IntClass{a}{b}\cdot\IntClass{e}{f}
      \end{aligned}}{%
      \FormulaRefAuto{a=b\vdash b=a}{2}}
    \proofstepwidestarbreakkeep[1]{%
      \begin{aligned}[t]
        &\IntClass{a}{b}\cdot
          (\IntClass{c}{d}+\IntClass{e}{f})={}\\
        &\IntClass{a}{b}\cdot\IntClass{c}{d}
          +\IntClass{a}{b}\cdot\IntClass{e}{f}
      \end{aligned}}{\rChain{5,6,7}}
  \closeproofpartwideindR

  \proofpartwide{Schluss}
    \proofstepwidestarbreakkeep[1]{z,w,u\in\mathbb Z}{\rA}
    \proofstepwidestarbreakkeep[1]{z\in\mathbb Z}{\rAEa{1}}
    \proofstepwidestarbreakkeep[1]{w\in\mathbb Z}{\rAEa{\rAEb{1}}}
    \proofstepwidestarbreakkeep[1]{u\in\mathbb Z}{\rAEb{\rAEb{1}}}
    \proofstepwidestarbreakkeep[1]{\exists a,b\in\mathbb N\,
      z=\IntClass{a}{b}}{\FormulaRefAuto{IntegerPairRepresentative}{2}}
    \proofstepwidestarbreakkeep[1]{\exists c,d\in\mathbb N\,
      w=\IntClass{c}{d}}{\FormulaRefAuto{IntegerPairRepresentative}{3}}
    \proofstepwidestarbreakkeep[1]{\exists e,f\in\mathbb N\,
      u=\IntClass{e}{f}}{\FormulaRefAuto{IntegerPairRepresentative}{4}}
    \proofstepwidestarbreakkeep[8]{a,b,c,d,e,f\in\mathbb N\land
      z=\IntClass{a}{b}\land w=\IntClass{c}{d}\land
      u=\IntClass{e}{f}}{\rA}
    \proofstepwidestarbreakkeep[8]{%
      \begin{aligned}[t]
        &\IntClass{a}{b}\cdot
          (\IntClass{c}{d}+\IntClass{e}{f})={}\\
        &\IntClass{a}{b}\cdot\IntClass{c}{d}
          +\IntClass{a}{b}\cdot\IntClass{e}{f}
      \end{aligned}}{%
      \FormulaRefAuto{IntMulLeftDistributiveOnClasses}{8}}
    \proofstepwidestarbreakkeep[8]{%
      z\cdot(w+u)=z\cdot w+z\cdot u}{\rIE{8,9}}
    \proofstepwidestarbreakkeep[1]{z\cdot(w+u)=z\cdot w+z\cdot u}{%
      \rEE{5,6,7,8,10}}
  \closeproofpartwide
\end{tabproofsplitwide}

\FormulaThmDeltaK[Rechtsdistributivität ganzer Zahlen]{%
  z,w,u\in\mathbb Z\vdash
  (z+w)\cdot u=z\cdot u+w\cdot u%
}{IntMulRightDistributive}{%
  \DeltaRow{Mengen}{z\dsep w\dsep u}%
}
\begin{tabproofwide}
  \proofstepwidestar[1]{z,w,u\in\mathbb Z}{\rA}
  \proofstepwidestar[1]{u\cdot z=z\cdot u}{%
    \FormulaRefAuto{IntMulCommutative}{1}}
  \proofstepwidestar[1]{u\cdot w=w\cdot u}{%
    \FormulaRefAuto{IntMulCommutative}{1}}
  \proofstepwide[1]{(z+w)\cdot u}{=}{u\cdot(z+w)}{%
    \FormulaRefAuto{IntMulCommutative}{%
      \FormulaRefAuto{IntAddClosure}{1},1}}
  \proofstepwide[1]{}{=}{u\cdot z+u\cdot w}{%
    \FormulaRefAuto{IntMulLeftDistributive}{1}}
  \proofstepwide[1]{}{=}{z\cdot u+w\cdot u}{\rIE{2,3}}
  \proofstepwidestar[1]{(z+w)\cdot u=z\cdot u+w\cdot u}{%
    \rChain{4,5,6}}
\end{tabproofwide}

\section{Verträglichkeit der Einbettung mit der Arithmetik}

\FormulaThmDeltaK[Additivität der natürlichen Einbettung]{%
  m,n\in\mathbb N\vdash
  \IntEmbed(m+n)=\IntEmbed(m)+\IntEmbed(n)%
}{NaturalIntegerEmbeddingAdditive}{%
  \DeltaRow{Mengen}{m\dsep n}%
  \DeltaRow{Funktion}{\IntEmbed}%
}
\begin{tabproofwide}
  \proofstepwidestar[1]{m,n\in\mathbb N}{\rA}
  \proofstepwidestar[1]{m+n\in\mathbb N}{%
    \FormulaRefAuto{PeanoAddClosure}{1}}
  \proofstepwidestar[1]{\IntEmbed(m)=\IntClass{m}{0}}{%
    \FormulaRefAuto{NaturalIntegerEmbeddingValue}[thm]{1}}
  \proofstepwidestar[1]{\IntEmbed(n)=\IntClass{n}{0}}{%
    \FormulaRefAuto{NaturalIntegerEmbeddingValue}[thm]{1}}
  \proofstepwidestar[]{0+0=0}{%
    \FormulaRefAuto{PeanoAddRightZero}{\FormulaRefAuto{PeanoZero}}}
  \proofstepwidestar[]{0=0+0}{%
    \FormulaRefAuto{a=b\vdash b=a}{5}}
  \proofstepwide[1]{\IntEmbed(m+n)}{=}{\IntClass{m+n}{0}}{%
    \FormulaRefAuto{NaturalIntegerEmbeddingValue}[thm]{2}}
  \proofstepwide[1]{}{=}{\IntClass{m+n}{0+0}}{\rIE{6}}
  \proofstepwide[1]{}{=}{\IntEmbed(m)+\IntEmbed(n)}{%
    \FormulaRefAuto{a=b\vdash b=a}{%
      \FormulaRefAuto{IntegerAdditionDef}{1,3,4}}}
  \proofstepwidestar[1]{\IntEmbed(m+n)=
    \IntEmbed(m)+\IntEmbed(n)}{\rChain{7,8,9}}
\end{tabproofwide}

\FormulaThmDeltaK[Multiplikativität der natürlichen Einbettung]{%
  m,n\in\mathbb N\vdash
  \IntEmbed(m\cdot n)=\IntEmbed(m)\cdot\IntEmbed(n)%
}{NaturalIntegerEmbeddingMultiplicative}{%
  \DeltaRow{Mengen}{m\dsep n}%
  \DeltaRow{Funktion}{\IntEmbed}%
}
\begin{tabproofsplitwide}
  \proofpartwideindR[Produkt der eingebetteten Faktoren]{%
    m,n\in\mathbb N\vdash
    \IntEmbed(m)\cdot\IntEmbed(n)=\IntClass{mn}{0}%
  }{%
    m,n\in\mathbb N\vdash
    \IntEmbed(m)\cdot\IntEmbed(n)=\IntClass{mn}{0}%
  }[NaturalIntegerEmbeddingProductClass]
    \proofstepwidestar[1]{m,n\in\mathbb N}{\rA}
    \proofstepwidestar[1]{mn\in\mathbb N}{%
      \FormulaRefAuto{PeanoMulClosure}{1}}
    \proofstepwidestar[1]{\IntEmbed(m)=\IntClass{m}{0}}{%
      \FormulaRefAuto{NaturalIntegerEmbeddingValue}[thm]{1}}
    \proofstepwidestar[1]{\IntEmbed(n)=\IntClass{n}{0}}{%
      \FormulaRefAuto{NaturalIntegerEmbeddingValue}[thm]{1}}
    \proofstepwidestar[]{0\cdot0=0}{%
      \FormulaRefAuto{PeanoMulLeftZero}{\FormulaRefAuto{PeanoZero}}}
    \proofstepwidestar[1]{m\cdot0=0}{%
      \FormulaRefAuto{PeanoMulRightZero}{1}}
    \proofstepwidestar[1]{0\cdot n=0}{%
      \FormulaRefAuto{PeanoMulLeftZero}{1}}
    \proofstepwidestar[1]{mn+0\cdot0=mn}{%
      \rIE{5,\FormulaRefAuto{PeanoAddRightZero}{2}}}
    \proofstepwidestar[1]{m\cdot0+0\cdot n=0}{%
      \rIE{6,7,
        \FormulaRefAuto{PeanoAddRightZero}{\FormulaRefAuto{PeanoZero}}}}
    \proofstepwide[1]{\IntEmbed(m)\cdot\IntEmbed(n)}{=}{%
      \IntClass{mn+0\cdot0}{m\cdot0+0\cdot n}}{%
      \FormulaRefAuto{IntegerMultiplicationDef}{1,3,4}}
    \proofstepwide[1]{}{=}{\IntClass{mn}{0}}{\rIE{8,9}}
    \proofstepwidestar[1]{%
      \IntEmbed(m)\cdot\IntEmbed(n)=\IntClass{mn}{0}}{%
      \rChain{10,11}}
  \closeproofpartwideindR

  \proofpartwideindR[Vergleich mit dem Produktbild]{%
    m,n\in\mathbb N\vdash
    \IntEmbed(m\cdot n)=\IntEmbed(m)\cdot\IntEmbed(n)%
  }{%
    m,n\in\mathbb N\vdash
    \IntEmbed(m\cdot n)=\IntEmbed(m)\cdot\IntEmbed(n)%
  }[NaturalIntegerEmbeddingMultiplicativeConclusion]
    \proofstepwidestar[1]{m,n\in\mathbb N}{\rA}
    \proofstepwidestar[1]{mn\in\mathbb N}{%
      \FormulaRefAuto{PeanoMulClosure}{1}}
    \proofstepwidestar[1]{%
      \IntEmbed(m)\cdot\IntEmbed(n)=\IntClass{mn}{0}}{%
      \FormulaRefAuto{NaturalIntegerEmbeddingProductClass}{1}}
    \proofstepwide[1]{\IntEmbed(mn)}{=}{\IntClass{mn}{0}}{%
      \FormulaRefAuto{NaturalIntegerEmbeddingValue}[thm]{2}}
    \proofstepwide[1]{}{=}{\IntEmbed(m)\cdot\IntEmbed(n)}{%
      \FormulaRefAuto{a=b\vdash b=a}{3}}
    \proofstepwidestar[1]{\IntEmbed(mn)=
      \IntEmbed(m)\cdot\IntEmbed(n)}{%
      \rChain{4,5}}
  \closeproofpartwideindR
\end{tabproofsplitwide}

\section{Herleitung der arithmetischen Axiome}

Bis hierher wurden Träger, Negation, Addition und Multiplikation am
Quotienten konstruiert und ihre Wohldefiniertheit sowie die elementaren
Rechengesetze am Modell bewiesen. Diese Gesetze werden nun einzeln als
künftige Axiome registriert. Das beendet die Modellphase noch nicht: Ordnung,
Diskretheit, Normalformen und zweiseitige Induktion müssen ebenfalls am
Konstruktionsmodell hergeleitet werden. Bis zu diesem vollständigen
Modellnachweis bleiben Klassenargumente daher zulässig, aber auf die
Herleitung genau dieser noch fehlenden Axiome beschränkt.

Der eigentliche Abstraktionsschnitt folgt erst im Kapitel
\emph{Axiomatische Beschreibung der ganzen Zahlen}. Ab dort werden weder
Paarrepräsentanten noch Äquivalenzklassen verwendet.

Die fett gesetzten Zwischenzeilen gliedern die Einzelaxiome lediglich
metasprachlich. Zusammen beschreiben die folgenden Blöcke einen kommutativen
Ring mit Eins; der abstrakte Ringbegriff wird dafür nicht vorausgesetzt.

\FormulaDefDeltaK[Elementare Axiome der Ganzzahlarithmetik]{%
  \IntegerArithmetic{\mathbb Z,\IntZero,\IntOne,-,+,\cdot}%
}{IntegerArithmeticAxioms}{%
  \DeltaRow{\textbf{Träger und Konstanten}}{}%
  \DeltaRow{Null}{\IntZero\in\mathbb Z}
    [\FormulaRefAuto{IntegerZeroCarrierAxiom}]%
  \DeltaRow{Eins}{\IntOne\in\mathbb Z}
    [\FormulaRefAuto{IntegerOneCarrierAxiom}]%
  \end{DeltaContext}
  \begin{DeltaContext}{ }%
  \DeltaRow{\makecell[l]{\textbf{Additive abelsche}\\
    \textbf{Gruppe}}}{}%
  \DeltaRow{Abschluss Negation}{z\in\mathbb Z\vdash-z\in\mathbb Z}
    [\FormulaRefAuto{IntegerNegClosureAxiom}]%
  \DeltaRow{Abschluss Addition}{z,w\in\mathbb Z\vdash z+w\in\mathbb Z}
    [\FormulaRefAuto{IntegerAddClosureAxiom}]%
  \DeltaRow{Addition: assoziativ}
    {\begin{aligned}[t]
       &z,w,u\in\mathbb Z\vdash{}\\[-2pt]
       &(z+w)+u=z+(w+u)
     \end{aligned}}
    [\FormulaRefAuto{IntegerAddAssociativeAxiom}]%
  \DeltaRow{Addition: kommutativ}
    {z,w\in\mathbb Z\vdash z+w=w+z}
    [\FormulaRefAuto{IntegerAddCommutativeAxiom}]%
  \DeltaRow{Addition: neutral}
    {\begin{aligned}[t]
       &z\in\mathbb Z\vdash{}\\[-2pt]
       &z+\IntZero=z\dsep\IntZero+z=z
     \end{aligned}}
    [\FormulaRefAuto{IntegerAddZeroAxiom}]%
  \DeltaRow{Addition: invers}
    {\begin{aligned}[t]
       &z\in\mathbb Z\vdash{}\\[-2pt]
       &z+(-z)=\IntZero\dsep(-z)+z=\IntZero
     \end{aligned}}
    [\FormulaRefAuto{IntegerAddInverseAxiom}]%
  \end{DeltaContext}
  \begin{DeltaContext}{ }%
  \DeltaRow{\makecell[l]{\textbf{Kommutatives}\\\textbf{Monoid}}}{}%
  \DeltaRow{Abschluss Produkt}
    {z,w\in\mathbb Z\vdash z\cdot w\in\mathbb Z}
    [\FormulaRefAuto{IntegerMulClosureAxiom}]%
  \DeltaRow{Produkt: assoziativ}
    {\begin{aligned}[t]
       &z,w,u\in\mathbb Z\vdash{}\\[-2pt]
       &(z\cdot w)\cdot u=z\cdot(w\cdot u)
     \end{aligned}}
    [\FormulaRefAuto{IntegerMulAssociativeAxiom}]%
  \DeltaRow{Produkt: kommutativ}
    {z,w\in\mathbb Z\vdash z\cdot w=w\cdot z}
    [\FormulaRefAuto{IntegerMulCommutativeAxiom}]%
  \DeltaRow{Produkt: neutral}
    {\begin{aligned}[t]
       &z\in\mathbb Z\vdash{}\\[-2pt]
       &z\cdot\IntOne=z\dsep\IntOne\cdot z=z
     \end{aligned}}
    [\FormulaRefAuto{IntegerMulOneAxiom}]%
  \end{DeltaContext}
  \begin{DeltaContext}{ }%
  \DeltaRow{\textbf{Distributivität}}{}%
  \DeltaRow{\makecell[l]{Links-\\distributivität}}
    {\begin{aligned}[t]
       &z,w,u\in\mathbb Z\vdash{}\\[-2pt]
       &z\cdot(w+u)=z\cdot w+z\cdot u
     \end{aligned}}
    [\FormulaRefAuto{IntegerMulLeftDistributiveAxiom}]%
  \DeltaRow{\makecell[l]{Rechts-\\distributivität}}
    {\begin{aligned}[t]
       &z,w,u\in\mathbb Z\vdash{}\\[-2pt]
       &(z+w)\cdot u=z\cdot u+w\cdot u
     \end{aligned}}
    [\FormulaRefAuto{IntegerMulRightDistributiveAxiom}]%
  \end{DeltaContext}
  \begin{DeltaContext}{ }%
  \DeltaRow{\textbf{Signatur}}{}%
  \DeltaRow{Ganze Zahlen}{\mathbb Z}%
  \DeltaRow{Konstanten}{\IntZero\dsep\IntOne}%
  \DeltaRow{Operationen}{-\dsep+\dsep\cdot}%
}

\FormulaAxiomAutoR[Null liegt in den ganzen Zahlen]{%
  \IntZero\in\mathbb Z%
}{IntegerZeroCarrierAxiom}
\par\noindent\emph{Herkunft aus der Quotientenkonstruktion:}
\FormulaRefAuto{IntZeroCarrier}[thm].\par

\FormulaAxiomAutoR[Eins liegt in den ganzen Zahlen]{%
  \IntOne\in\mathbb Z%
}{IntegerOneCarrierAxiom}
\par\noindent\emph{Herkunft aus der Quotientenkonstruktion:}
\FormulaRefAuto{IntOneCarrier}[thm].\par

\FormulaAxiomAutoR[Abgeschlossenheit unter Negation]{%
  z\in\mathbb Z\vdash-z\in\mathbb Z%
}{IntegerNegClosureAxiom}
\par\noindent\emph{Herkunft aus der Quotientenkonstruktion:}
\FormulaRefAuto{IntNegClosure}[thm].\par

\FormulaAxiomAutoR[Abgeschlossenheit der Ganzzahladdition]{%
  z,w\in\mathbb Z\vdash z+w\in\mathbb Z%
}{IntegerAddClosureAxiom}
\par\noindent\emph{Herkunft aus der Quotientenkonstruktion:}
\FormulaRefAuto{IntAddClosure}[thm].\par

\FormulaAxiomAutoR[Assoziativität der Ganzzahladdition]{%
  z,w,u\in\mathbb Z\vdash(z+w)+u=z+(w+u)%
}{IntegerAddAssociativeAxiom}
\par\noindent\emph{Herkunft aus der Quotientenkonstruktion:}
\FormulaRefAuto{IntAddAssociative}[thm].\par

\FormulaAxiomAutoR[Kommutativität der Ganzzahladdition]{%
  z,w\in\mathbb Z\vdash z+w=w+z%
}{IntegerAddCommutativeAxiom}
\par\noindent\emph{Herkunft aus der Quotientenkonstruktion:}
\FormulaRefAuto{IntAddCommutative}[thm].\par

\FormulaAxiomAutoR[Nullgesetze der Ganzzahladdition]{%
  z\in\mathbb Z\vdash
  z+\IntZero=z\land\IntZero+z=z%
}{IntegerAddZeroAxiom}
\par\noindent\emph{Herkunft aus der Quotientenkonstruktion:}
\FormulaRefAuto{IntAddZero}[thm].\par

\FormulaAxiomAutoR[Additive Inversen in den ganzen Zahlen]{%
  z\in\mathbb Z\vdash
  z+(-z)=\IntZero\land(-z)+z=\IntZero%
}{IntegerAddInverseAxiom}
\par\noindent\emph{Herkunft aus der Quotientenkonstruktion:}
\FormulaRefAuto{IntAddInverse}[thm].\par

\FormulaAxiomAutoR[Abgeschlossenheit der Ganzzahlmultiplikation]{%
  z,w\in\mathbb Z\vdash z\cdot w\in\mathbb Z%
}{IntegerMulClosureAxiom}
\par\noindent\emph{Herkunft aus der Quotientenkonstruktion:}
\FormulaRefAuto{IntMulClosure}[thm].\par

\FormulaAxiomAutoR[Assoziativität der Ganzzahlmultiplikation]{%
  z,w,u\in\mathbb Z\vdash
  (z\cdot w)\cdot u=z\cdot(w\cdot u)%
}{IntegerMulAssociativeAxiom}
\par\noindent\emph{Herkunft aus der Quotientenkonstruktion:}
\FormulaRefAuto{IntMulAssociative}[thm].\par

\FormulaAxiomAutoR[Kommutativität der Ganzzahlmultiplikation]{%
  z,w\in\mathbb Z\vdash z\cdot w=w\cdot z%
}{IntegerMulCommutativeAxiom}
\par\noindent\emph{Herkunft aus der Quotientenkonstruktion:}
\FormulaRefAuto{IntMulCommutative}[thm].\par

\FormulaAxiomAutoR[Einheitsgesetze der Ganzzahlmultiplikation]{%
  z\in\mathbb Z\vdash
  z\cdot\IntOne=z\land\IntOne\cdot z=z%
}{IntegerMulOneAxiom}
\par\noindent\emph{Herkunft aus der Quotientenkonstruktion:}
\FormulaRefAuto{IntMulOne}[thm].\par

\FormulaAxiomAutoR[Linksdistributivität der Ganzzahlmultiplikation]{%
  z,w,u\in\mathbb Z\vdash
  z\cdot(w+u)=z\cdot w+z\cdot u%
}{IntegerMulLeftDistributiveAxiom}
\par\noindent\emph{Herkunft aus der Quotientenkonstruktion:}
\FormulaRefAuto{IntMulLeftDistributive}[thm].\par

\FormulaAxiomAutoR[Rechtsdistributivität der Ganzzahlmultiplikation]{%
  z,w,u\in\mathbb Z\vdash
  (z+w)\cdot u=z\cdot u+w\cdot u%
}{IntegerMulRightDistributiveAxiom}
\par\noindent\emph{Herkunft aus der Quotientenkonstruktion:}
\FormulaRefAuto{IntMulRightDistributive}[thm].\par

\FormulaThmDeltaK[Multiplikation ganzer Zahlen mit null aus den Axiomen]{%
  z\in\mathbb Z\vdash
  z\cdot\IntZero=\IntZero\dsep
  \IntZero\cdot z=\IntZero%
}{IntegerMulZeroFromAxioms}{%
  \DeltaRow{Mengen}{z}%
}
\begin{tabproofwide}
  \proofstepwidestarbreak[1]{z\in\mathbb Z}{\rA}
  \proofstepwidestarbreak[1]{z\cdot\IntZero\in\mathbb Z}{%
    \FormulaRefAuto{IntegerMulClosureAxiom}[ax]{%
      1,\FormulaRefAuto{IntegerZeroCarrierAxiom}[ax]}}
  \proofstepwidestarbreak[]{\IntZero+\IntZero=\IntZero}{%
    \rAEa{\FormulaRefAuto{IntegerAddZeroAxiom}[ax]{%
      \FormulaRefAuto{IntegerZeroCarrierAxiom}[ax]}}}
  \proofstepwidestarbreak[1]{z\cdot(\IntZero+\IntZero)
    =z\cdot\IntZero+z\cdot\IntZero}{%
    \FormulaRefAuto{IntegerMulLeftDistributiveAxiom}[ax]{%
      1,\FormulaRefAuto{IntegerZeroCarrierAxiom}[ax]}}
  \proofstepwidestarbreak[1]{z\cdot\IntZero
    =z\cdot(\IntZero+\IntZero)}{\rIE{3}}
  \proofstepwidestarbreak[1]{z\cdot\IntZero
    =z\cdot\IntZero+z\cdot\IntZero}{\rChain{5,4}}
  \proofstepwidestarbreak[1]{-(z\cdot\IntZero)
    +(z\cdot\IntZero)=\IntZero}{%
    \rAEb{\FormulaRefAuto{IntegerAddInverseAxiom}[ax]{2}}}
  \proofstepwidestarbreak[1]{-(z\cdot\IntZero)\in\mathbb Z}{%
    \FormulaRefAuto{IntegerNegClosureAxiom}[ax]{2}}
  \proofstepwidestarbreak[1]{%
    \bigl(-(z\cdot\IntZero)+(z\cdot\IntZero)\bigr)
      +z\cdot\IntZero
    =-(z\cdot\IntZero)
      +(z\cdot\IntZero+z\cdot\IntZero)}{%
    \FormulaRefAuto{IntegerAddAssociativeAxiom}[ax]{8,2,2}}
  \proofstepwidestarbreak[1]{%
    \bigl(-(z\cdot\IntZero)+(z\cdot\IntZero)\bigr)
      +z\cdot\IntZero
    =\IntZero+z\cdot\IntZero}{\rIE{7}}
  \proofstepwidestarbreak[1]{\IntZero+z\cdot\IntZero
    =z\cdot\IntZero}{%
    \rAEb{\FormulaRefAuto{IntegerAddZeroAxiom}[ax]{2}}}
  \proofstepwidestarbreak[1]{%
    \bigl(-(z\cdot\IntZero)+(z\cdot\IntZero)\bigr)
      +z\cdot\IntZero
    =z\cdot\IntZero}{\rChain{10,11}}
  \proofstepwide[1]{z\cdot\IntZero}{=}{%
    -(z\cdot\IntZero)+(z\cdot\IntZero+z\cdot\IntZero)}{%
    \FormulaRefAuto{a=b\vdash b=a}{\rIE{9,12}}}
  \proofstepwide[1]{}{=}{%
    -(z\cdot\IntZero)+(z\cdot\IntZero)}{%
    \FormulaRefAuto{a=b\vdash b=a}{\rIE{6}}}
  \proofstepwide[1]{}{=}{\IntZero}{%
    \rAEb{\FormulaRefAuto{IntegerAddInverseAxiom}[ax]{2}}}
  \proofstepwidestarbreak[1]{z\cdot\IntZero=\IntZero}{%
    \rChain{13,14,15}}
  \proofstepwidestarbreak[1]{\IntZero\cdot z=z\cdot\IntZero}{%
    \FormulaRefAuto{IntegerMulCommutativeAxiom}[ax]{%
      \FormulaRefAuto{IntegerZeroCarrierAxiom}[ax],1}}
  \proofstepwidestarbreak[1]{\IntZero\cdot z=\IntZero}{\rChain{17,16}}
  \proofstepwidestarbreak[1]{z\cdot\IntZero=\IntZero\land
    \IntZero\cdot z=\IntZero}{\rAI{16,18}}
\end{tabproofwide}

\FormulaThmDeltaK[Negation der Ganzzahlnull]{%
  -\IntZero=\IntZero%
}{IntegerNegativeZero}{%
  \DeltaRow{Mengen}{\IntZero}%
}
\begin{tabproofwide}
  \proofstepwidestar[]{\IntZero\in\mathbb Z}{%
    \FormulaRefAuto{IntegerZeroCarrierAxiom}[ax]}
  \proofstepwidestar[]{-\IntZero\in\mathbb Z}{%
    \FormulaRefAuto{IntegerNegClosureAxiom}[ax]{1}}
  \proofstepwide[]{-\IntZero}{=}{(-\IntZero)+\IntZero}{%
    \FormulaRefAuto{a=b\vdash b=a}{%
      \rAEa{\FormulaRefAuto{IntegerAddZeroAxiom}[ax]{2}}}}
  \proofstepwide[]{}{=}{\IntZero}{%
    \rAEb{\FormulaRefAuto{IntegerAddInverseAxiom}[ax]{1}}}
  \proofstepwidestar[]{-\IntZero=\IntZero}{\rChain{3,4}}
\end{tabproofwide}

\FormulaThmDeltaK[Negatives Vorzeichen im rechten Faktor]{%
  x,y\in\mathbb Z\vdash
  x\cdot(-y)=-(x\cdot y)%
}{IntegerRightNegativeProduct}{%
  \DeltaRow{Mengen}{x\dsep y}%
}
\begin{tabproofwide}
  \proofstepwidestarbreakkeep[1]{x,y\in\mathbb Z}{\rA}
  \proofstepwidestarbreakkeep[1]{x\in\mathbb Z}{\rAEa{1}}
  \proofstepwidestarbreakkeep[1]{y\in\mathbb Z}{\rAEb{1}}
  \proofstepwidestarbreakkeep[1]{-y\in\mathbb Z}{%
    \FormulaRefAuto{IntegerNegClosureAxiom}[ax]{3}}
  \proofstepwidestarbreakkeep[1]{x\cdot y\in\mathbb Z}{%
    \FormulaRefAuto{IntegerMulClosureAxiom}[ax]{2,3}}
  \proofstepwidestarbreakkeep[1]{x\cdot(-y)\in\mathbb Z}{%
    \FormulaRefAuto{IntegerMulClosureAxiom}[ax]{2,4}}
  \proofstepwidestarbreakkeep[1]{-(x\cdot y)\in\mathbb Z}{%
    \FormulaRefAuto{IntegerNegClosureAxiom}[ax]{5}}
  \proofstepwidestarbreakkeep[1]{y+(-y)=\IntZero}{%
    \rAEa{\FormulaRefAuto{IntegerAddInverseAxiom}[ax]{3}}}
  \proofstepwidestarbreakkeep[1]{%
    x\cdot\bigl(y+(-y)\bigr)
      =x\cdot y+x\cdot(-y)}{%
    \FormulaRefAuto{IntegerMulLeftDistributiveAxiom}[ax]{2,3,4}}
  \proofstepwide[1]{x\cdot y+x\cdot(-y)}{=}{%
    x\cdot\bigl(y+(-y)\bigr)}{%
    \FormulaRefAuto{a=b\vdash b=a}{9}}
  \proofstepwide[1]{}{=}{x\cdot\IntZero}{\rIE{8}}
  \proofstepwide[1]{}{=}{\IntZero}{%
    \rAEa{\FormulaRefAuto{IntegerMulZeroFromAxioms}[thm]{2}}}
  \proofstepwidestarbreakkeep[1]{%
    x\cdot y+x\cdot(-y)=\IntZero}{\rChain{10,11,12}}
  \proofstepwidestarbreakkeep[1]{%
    -(x\cdot y)+x\cdot y=\IntZero}{%
    \rAEb{\FormulaRefAuto{IntegerAddInverseAxiom}[ax]{5}}}
  \proofstepwidestarbreakkeep[1]{%
    \IntZero+x\cdot(-y)=x\cdot(-y)}{%
    \rAEb{\FormulaRefAuto{IntegerAddZeroAxiom}[ax]{6}}}
  \proofstepwide[1]{x\cdot(-y)}{=}{\IntZero+x\cdot(-y)}{%
    \FormulaRefAuto{a=b\vdash b=a}{15}}
  \proofstepwidestarbreakkeep[1]{%
    \IntZero=-(x\cdot y)+x\cdot y}{%
    \FormulaRefAuto{a=b\vdash b=a}{14}}
  \proofstepwide[1]{}{=}{%
    \bigl(-(x\cdot y)+x\cdot y\bigr)+x\cdot(-y)}{\rIE{17}}
  \proofstepwide[1]{}{=}{%
    -(x\cdot y)+\bigl(x\cdot y+x\cdot(-y)\bigr)}{%
    \FormulaRefAuto{IntegerAddAssociativeAxiom}[ax]{7,5,6}}
  \proofstepwide[1]{}{=}{-(x\cdot y)+\IntZero}{\rIE{13}}
  \proofstepwide[1]{}{=}{-(x\cdot y)}{%
    \rAEa{\FormulaRefAuto{IntegerAddZeroAxiom}[ax]{7}}}
  \proofstepwidestarbreakkeep[1]{%
    x\cdot(-y)=-(x\cdot y)}{\rChain{16,18,19,20,21}}
\end{tabproofwide}

\FormulaThmDeltaK[Produkt zweier negativer ganzer Zahlen]{%
  x,y\in\mathbb Z\vdash
  (-x)\cdot(-y)=x\cdot y%
}{IntegerNegativeTimesNegative}{%
  \DeltaRow{Mengen}{x\dsep y}%
}
\begin{tabproofwide}
  \proofstepwidestarbreakkeep[1]{x,y\in\mathbb Z}{\rA}
  \proofstepwidestarbreakkeep[1]{x\in\mathbb Z}{\rAEa{1}}
  \proofstepwidestarbreakkeep[1]{y\in\mathbb Z}{\rAEb{1}}
  \proofstepwidestarbreakkeep[1]{-x\in\mathbb Z}{%
    \FormulaRefAuto{IntegerNegClosureAxiom}[ax]{2}}
  \proofstepwidestarbreakkeep[1]{-y\in\mathbb Z}{%
    \FormulaRefAuto{IntegerNegClosureAxiom}[ax]{3}}
  \proofstepwidestarbreakkeep[1]{y\cdot x=x\cdot y}{%
    \FormulaRefAuto{IntegerMulCommutativeAxiom}[ax]{3,2}}
  \proofstepwide[1]{(-x)\cdot y}{=}{y\cdot(-x)}{%
    \FormulaRefAuto{IntegerMulCommutativeAxiom}[ax]{4,3}}
  \proofstepwide[1]{}{=}{-(y\cdot x)}{%
    \FormulaRefAuto{IntegerRightNegativeProduct}[thm]{3,2}}
  \proofstepwide[1]{}{=}{-(x\cdot y)}{\rIE{6}}
  \proofstepwidestarbreakkeep[1]{%
    (-x)\cdot y=-(x\cdot y)}{\rChain{7,8,9}}
  \proofstepwidestarbreakkeep[1]{x\cdot y\in\mathbb Z}{%
    \FormulaRefAuto{IntegerMulClosureAxiom}[ax]{2,3}}
  \proofstepwidestarbreakkeep[1]{%
    -\bigl(-(x\cdot y)\bigr)=x\cdot y}{%
    \FormulaRefAuto{IntegerDoubleNegation}[thm]{11}}
  \proofstepwide[1]{(-x)\cdot(-y)}{=}{%
    -\bigl((-x)\cdot y\bigr)}{%
    \FormulaRefAuto{IntegerRightNegativeProduct}[thm]{4,3}}
  \proofstepwide[1]{}{=}{-\bigl(-(x\cdot y)\bigr)}{\rIE{10}}
  \proofstepwide[1]{}{=}{x\cdot y}{%
    \FormulaRefAuto{IntegerDoubleNegation}[thm]{11}}
  \proofstepwidestarbreakkeep[1]{%
    (-x)\cdot(-y)=x\cdot y}{\rChain{13,14,15}}
\end{tabproofwide}

\chapter{Nachfolger, Vorgänger und Normalformen}

\section{Nachfolger und Vorgänger}

\FormulaDefDeltaK[Nachfolger einer ganzen Zahl]{%
  \begin{aligned}[t]
    &\IntSucc\colon\mathbb Z\to\mathbb Z,\\[-2pt]
    &\forall z\in\mathbb Z\,
      \bigl(\IntSucc(z)\coloneqq z+\IntOne\bigr)
  \end{aligned}%
}{IntegerSuccessorDef}{%
  \DeltaRow{Mengen}{z}%
  \DeltaRow{Funktionensymbol}{\IntSucc}%
  \DeltaRow{Neue Symbole}{\IntSucc}%
}
\begin{tabproof}
  \proofstep{1}{z\in\mathbb Z}{\rA}
  \proofstep{}{\IntOne\in\mathbb Z}{%
    \FormulaRefAuto{IntegerOneCarrierAxiom}[ax]}
  \proofstep{1}{z+\IntOne\in\mathbb Z}{%
    \FormulaRefAuto{IntegerAddClosureAxiom}[ax]{1,2}}
  \proofstep{}{\forall z\in\mathbb Z\,\bigl(z+\IntOne\in\mathbb Z\bigr)}{%
    \rUI{\rRI{1,3}}}
\end{tabproof}

\FormulaDefDeltaK[Vorgänger einer ganzen Zahl]{%
  \begin{aligned}[t]
    &\IntPred\colon\mathbb Z\to\mathbb Z,\\[-2pt]
    &\forall z\in\mathbb Z\,
      \bigl(\IntPred(z)\coloneqq z+(-\IntOne)\bigr)
  \end{aligned}%
}{IntegerPredecessorDef}{%
  \DeltaRow{Mengen}{z}%
  \DeltaRow{Funktionensymbol}{\IntPred}%
  \DeltaRow{Neue Symbole}{\IntPred}%
}
\begin{tabproof}
  \proofstep{1}{z\in\mathbb Z}{\rA}
  \proofstep{}{\IntOne\in\mathbb Z}{%
    \FormulaRefAuto{IntegerOneCarrierAxiom}[ax]}
  \proofstep{}{-\IntOne\in\mathbb Z}{%
    \FormulaRefAuto{IntegerNegClosureAxiom}[ax]{2}}
  \proofstep{1}{z+(-\IntOne)\in\mathbb Z}{%
    \FormulaRefAuto{IntegerAddClosureAxiom}[ax]{1,3}}
  \proofstep{}{%
    \forall z\in\mathbb Z\,\bigl(z+(-\IntOne)\in\mathbb Z\bigr)}{%
    \rUI{\rRI{1,4}}}
\end{tabproof}

\FormulaThmDeltaK[Nachfolgerfunktion auf den ganzen Zahlen]{%
  \IntSucc\colon\mathbb Z\to\mathbb Z%
}{IntegerSuccessorFunction}{%
  \DeltaRow{Funktion}{\IntSucc}%
}
\begin{tabproof}
  \proofstep{}{\IntSucc\colon\mathbb Z\to\mathbb Z}{%
    \FormulaRefAuto{IntegerSuccessorDef}[def]}
\end{tabproof}

\FormulaThmDeltaK[Vorgängerfunktion auf den ganzen Zahlen]{%
  \IntPred\colon\mathbb Z\to\mathbb Z%
}{IntegerPredecessorFunction}{%
  \DeltaRow{Funktion}{\IntPred}%
}
\begin{tabproof}
  \proofstep{}{\IntPred\colon\mathbb Z\to\mathbb Z}{%
    \FormulaRefAuto{IntegerPredecessorDef}[def]}
\end{tabproof}

\FormulaThmDeltaK[Grundgleichung des ganzzahligen Nachfolgers]{%
  z\in\mathbb Z\vdash \IntSucc(z)=z+\IntOne%
}{IntegerSuccessorValue}{%
  \DeltaRow{Mengen}{z}%
  \DeltaRow{Funktion}{\IntSucc}%
}
\begin{tabproof}
  \proofstep{1}{z\in\mathbb Z}{\rA}
  \proofstep{}{%
    \forall u\in\mathbb Z\,
      \bigl(\IntSucc(u)=u+\IntOne\bigr)}{%
    \FormulaRefAuto{IntegerSuccessorDef}[def]}
  \proofstep{1}{\IntSucc(z)=z+\IntOne}{%
    \rRE{\rUE{2},1}}
\end{tabproof}

\FormulaThmDeltaK[Grundgleichung des ganzzahligen Vorgängers]{%
  z\in\mathbb Z\vdash \IntPred(z)=z+(-\IntOne)%
}{IntegerPredecessorValue}{%
  \DeltaRow{Mengen}{z}%
  \DeltaRow{Funktion}{\IntPred}%
}
\begin{tabproof}
  \proofstep{1}{z\in\mathbb Z}{\rA}
  \proofstep{}{%
    \forall u\in\mathbb Z\,
      \bigl(\IntPred(u)=u+(-\IntOne)\bigr)}{%
    \FormulaRefAuto{IntegerPredecessorDef}[def]}
  \proofstep{1}{\IntPred(z)=z+(-\IntOne)}{%
    \rRE{\rUE{2},1}}
\end{tabproof}

\FormulaThmDeltaK[Nachfolger auf Klassen]{%
  a,b\in\mathbb N\vdash
  \IntSucc(\IntClass{a}{b})=\IntClass{\Succ(a)}{b}%
}{IntegerSuccessorClassValue}{%
  \DeltaRow{Mengen}{a\dsep b}%
}
\begin{tabproofwide}
  \proofstepwidestar[1]{a,b\in\mathbb N}{\rA}
  \proofstepwidestar[1]{\IntClass{a}{b}\in\mathbb Z}{%
    \FormulaRefAuto{IntegerClassInZ}{1}}
  \proofstepwidestar[]{%
    \forall z\in\mathbb Z\,\bigl(\IntSucc(z)=z+\IntOne\bigr)}{%
    \FormulaRefAuto{IntegerSuccessorDef}[def]}
  \proofstepwidestar[]{\IntOne=\IntClass{1}{0}}{%
    \FormulaRefAuto{IntegerOneDef}{%
      \FormulaRefAuto{NaturalIntegerEmbeddingValue}[thm]{%
        \FormulaRefAuto{PeanoOneInN}}}}
  \proofstepwidestar[1]{a+1=\Succ(a)}{%
    \FormulaRefAuto{PeanoAddOneSucc}{1}}
  \proofstepwidestar[1]{b+0=b}{\FormulaRefAuto{PeanoAddRightZero}{1}}
  \proofstepwide[1]{\IntSucc(\IntClass{a}{b})}{=}{%
    \IntClass{a}{b}+\IntOne}{\rRE{\rUE{3},2}}
  \proofstepwide[1]{}{=}{\IntClass{a+1}{b+0}}{%
    \FormulaRefAuto{IntegerAdditionDef}{1,4}}
  \proofstepwide[1]{}{=}{\IntClass{\Succ(a)}{b}}{\rIE{5,6}}
  \proofstepwidestar[1]{\IntSucc(\IntClass{a}{b})=
    \IntClass{\Succ(a)}{b}}{\rChain{7,8,9}}
\end{tabproofwide}

\FormulaThmDeltaK[Vorgänger auf Klassen]{%
  a,b\in\mathbb N\vdash
  \IntPred(\IntClass{a}{b})=\IntClass{a}{\Succ(b)}%
}{IntegerPredecessorClassValue}{%
  \DeltaRow{Mengen}{a\dsep b}%
}
\begin{tabproofwide}
  \proofstepwidestar[1]{a,b\in\mathbb N}{\rA}
  \proofstepwidestar[1]{\IntClass{a}{b}\in\mathbb Z}{%
    \FormulaRefAuto{IntegerClassInZ}{1}}
  \proofstepwidestar[]{%
    \forall z\in\mathbb Z\,\bigl(\IntPred(z)=z+(-\IntOne)\bigr)}{%
    \FormulaRefAuto{IntegerPredecessorDef}[def]}
  \proofstepwidestar[]{-\IntOne=\IntClass{0}{1}}{%
    \FormulaRefAuto{IntegerNegationDef}{%
      \FormulaRefAuto{IntegerOneDef},%
      \FormulaRefAuto{NaturalIntegerEmbeddingValue}[thm]}}
  \proofstepwidestar[1]{a+0=a}{\FormulaRefAuto{PeanoAddRightZero}{1}}
  \proofstepwidestar[1]{b+1=\Succ(b)}{\FormulaRefAuto{PeanoAddOneSucc}{1}}
  \proofstepwide[1]{\IntPred(\IntClass{a}{b})}{=}{%
    \IntClass{a}{b}+(-\IntOne)}{\rRE{\rUE{3},2}}
  \proofstepwide[1]{}{=}{\IntClass{a+0}{b+1}}{%
    \FormulaRefAuto{IntegerAdditionDef}{1,4}}
  \proofstepwide[1]{}{=}{\IntClass{a}{\Succ(b)}}{\rIE{5,6}}
  \proofstepwidestar[1]{\IntPred(\IntClass{a}{b})=
    \IntClass{a}{\Succ(b)}}{\rChain{7,8,9}}
\end{tabproofwide}

\FormulaThmDeltaK[Vorgänger nach Nachfolger]{%
  z\in\mathbb Z\vdash
  \IntPred(\IntSucc(z))=z%
}{IntegerPredAfterSucc}{%
  \DeltaRow{Mengen}{z}%
}
\begin{tabproofsplitwide}
  \proofpartwideindR[Auflösen der Definitionen]{%
    z\in\mathbb Z\vdash
    \IntPred(\IntSucc(z))=(z+\IntOne)+(-\IntOne)%
  }{%
    z\in\mathbb Z\vdash
    \IntPred(\IntSucc(z))=(z+\IntOne)+(-\IntOne)%
  }[IntegerPredAfterSuccExpansion]
    \proofstepwidestar[1]{z\in\mathbb Z}{\rA}
    \proofstepwidestar[1]{\IntSucc(z)=z+\IntOne}{%
      \FormulaRefAuto{IntegerSuccessorValue}[thm]{1}}
    \proofstepwidestar[1]{\IntSucc(z)\in\mathbb Z}{%
      \FormulaRefAuto{F\colon A\to B\dsep x\in A\vdash F(x)\in B}{%
        \FormulaRefAuto{IntegerSuccessorFunction},1}}
    \proofstepwidestar[1]{\IntPred(\IntSucc(z))=
      \IntSucc(z)+(-\IntOne)}{%
      \FormulaRefAuto{IntegerPredecessorValue}[thm]{3}}
    \proofstepwidestar[1]{\IntPred(\IntSucc(z))=
      (z+\IntOne)+(-\IntOne)}{\rIE{2,4}}
  \closeproofpartwideindR

  \proofpartwide{Schluss}
    \proofstepwidestar[1]{z\in\mathbb Z}{\rA}
    \proofstepwidestar[]{\IntOne\in\mathbb Z}{%
      \FormulaRefAuto{IntegerOneCarrierAxiom}[ax]}
    \proofstepwidestar[]{\IntOne+(-\IntOne)=\IntZero}{%
      \rAEa{\FormulaRefAuto{IntegerAddInverseAxiom}[ax]{2}}}
    \proofstepwide[1]{\IntPred(\IntSucc(z))}{=}{%
      (z+\IntOne)+(-\IntOne)}{%
      \FormulaRefAuto{IntegerPredAfterSuccExpansion}{1}}
    \proofstepwide[1]{}{=}{z+(\IntOne+(-\IntOne))}{%
      \FormulaRefAuto{IntegerAddAssociativeAxiom}[ax]{1,2,%
        \FormulaRefAuto{IntegerNegClosureAxiom}[ax]{2}}}
    \proofstepwide[1]{}{=}{z+\IntZero}{\rIE{3}}
    \proofstepwide[1]{}{=}{z}{%
      \rAEa{\FormulaRefAuto{IntegerAddZeroAxiom}[ax]{1}}}
    \proofstepwidestar[1]{\IntPred(\IntSucc(z))=z}{%
      \rChain{4,5,6,7}}
  \closeproofpartwide
\end{tabproofsplitwide}

\FormulaThmDeltaK[Nachfolger nach Vorgänger]{%
  z\in\mathbb Z\vdash
  \IntSucc(\IntPred(z))=z%
}{IntegerSuccAfterPred}{%
  \DeltaRow{Mengen}{z}%
}
\begin{tabproofsplitwide}
  \proofpartwideindR[Auflösen der Definitionen]{%
    z\in\mathbb Z\vdash
    \IntSucc(\IntPred(z))=(z+(-\IntOne))+\IntOne%
  }{%
    z\in\mathbb Z\vdash
    \IntSucc(\IntPred(z))=(z+(-\IntOne))+\IntOne%
  }[IntegerSuccAfterPredExpansion]
    \proofstepwidestar[1]{z\in\mathbb Z}{\rA}
    \proofstepwidestar[1]{\IntPred(z)=z+(-\IntOne)}{%
      \FormulaRefAuto{IntegerPredecessorValue}[thm]{1}}
    \proofstepwidestar[1]{\IntPred(z)\in\mathbb Z}{%
      \FormulaRefAuto{F\colon A\to B\dsep x\in A\vdash F(x)\in B}{%
        \FormulaRefAuto{IntegerPredecessorFunction},1}}
    \proofstepwidestar[1]{%
      \IntSucc(\IntPred(z))=\IntPred(z)+\IntOne}{%
      \FormulaRefAuto{IntegerSuccessorValue}[thm]{3}}
    \proofstepwidestar[1]{%
      \IntSucc(\IntPred(z))=(z+(-\IntOne))+\IntOne}{%
      \rIE{2,4}}
  \closeproofpartwideindR

  \proofpartwideindR[Assoziativität und Inverses]{%
    z\in\mathbb Z\vdash\IntSucc(\IntPred(z))=z%
  }{%
    z\in\mathbb Z\vdash\IntSucc(\IntPred(z))=z%
  }[IntegerSuccAfterPredConclusion]
    \proofstepwidestar[1]{z\in\mathbb Z}{\rA}
    \proofstepwidestar[]{\IntOne\in\mathbb Z}{%
      \FormulaRefAuto{IntegerOneCarrierAxiom}[ax]}
    \proofstepwidestar[]{-\IntOne\in\mathbb Z}{%
      \FormulaRefAuto{IntegerNegClosureAxiom}[ax]{2}}
    \proofstepwidestar[]{%
      (-\IntOne)+\IntOne=\IntZero}{%
      \rAEb{\FormulaRefAuto{IntegerAddInverseAxiom}[ax]{2}}}
    \proofstepwide[1]{\IntSucc(\IntPred(z))}{=}{%
      (z+(-\IntOne))+\IntOne}{%
      \FormulaRefAuto{IntegerSuccAfterPredExpansion}{1}}
    \proofstepwide[1]{}{=}{z+((-\IntOne)+\IntOne)}{%
      \FormulaRefAuto{IntegerAddAssociativeAxiom}[ax]{1,3,2}}
    \proofstepwide[1]{}{=}{z+\IntZero}{\rIE{4}}
    \proofstepwide[1]{}{=}{z}{%
      \rAEa{\FormulaRefAuto{IntegerAddZeroAxiom}[ax]{1}}}
    \proofstepwidestar[1]{\IntSucc(\IntPred(z))=z}{%
      \rChain{5,6,7,8}}
  \closeproofpartwideindR
\end{tabproofsplitwide}

\FormulaThmDeltaK[Bijektivität des Ganzzahlnachfolgers]{%
  \IntSucc\colon\mathbb Z\bij\mathbb Z%
}{IntegerSuccessorBijective}{%
  \DeltaRow{Funktionen}{\IntSucc\dsep\IntPred}%
}
\begin{tabproofsplitwide}
  \proofpartwideindR[Injektivität]{%
    x,y\in\mathbb Z\dsep
    \IntSucc(x)=\IntSucc(y)\vdash x=y%
  }{%
    x,y\in\mathbb Z\dsep
    \IntSucc(x)=\IntSucc(y)\vdash x=y%
  }[IntegerSuccessorInjective]
    \proofstepwidestar[1]{x,y\in\mathbb Z}{\rA}
    \proofstepwidestar[2]{\IntSucc(x)=\IntSucc(y)}{\rA}
    \proofstepwidestar[1]{\IntPred(\IntSucc(x))=x}{%
      \FormulaRefAuto{IntegerPredAfterSucc}{1}}
    \proofstepwidestar[1]{\IntPred(\IntSucc(y))=y}{%
      \FormulaRefAuto{IntegerPredAfterSucc}{1}}
    \proofstepwide[1]{x}{=}{\IntPred(\IntSucc(x))}{%
      \FormulaRefAuto{a=b\vdash b=a}{3}}
    \proofstepwide[1,2]{}{=}{\IntPred(\IntSucc(y))}{\rIE{2}}
    \proofstepwide[1]{}{=}{y}{%
      \FormulaRefAuto{IntegerPredAfterSucc}{1}}
    \proofstepwidestar[1,2]{x=y}{\rChain{5,6,7}}
  \closeproofpartwideindR

  \proofpartwideindR[Surjektivität]{%
    y\in\mathbb Z\vdash
    \exists x\in\mathbb Z\,\IntSucc(x)=y%
  }{%
    y\in\mathbb Z\vdash
    \exists x\in\mathbb Z\,\IntSucc(x)=y%
  }[IntegerSuccessorSurjective]
    \proofstepwidestar[1]{y\in\mathbb Z}{\rA}
    \proofstepwidestar[1]{\IntPred(y)\in\mathbb Z}{%
      \FormulaRefAuto{F\colon A\to B\dsep x\in A\vdash F(x)\in B}{%
        \FormulaRefAuto{IntegerPredecessorFunction},1}}
    \proofstepwidestar[1]{\IntSucc(\IntPred(y))=y}{%
      \FormulaRefAuto{IntegerSuccAfterPred}{1}}
    \proofstepwidestar[1]{\exists x\in\mathbb Z\,\IntSucc(x)=y}{%
      \rEI{\rAI{2,3}}}
  \closeproofpartwideindR

  \proofpartwide{Schluss}
    \proofstepwidestar[]{\IntSucc\colon\mathbb Z\to\mathbb Z}{%
      \FormulaRefAuto{IntegerSuccessorFunction}}
    \proofstepwidestar[]{%
      \forall x,y\in\mathbb Z\,
      (\IntSucc(x)=\IntSucc(y)\rightarrow x=y)}{%
      \FormulaRefAuto{IntegerSuccessorInjective}}
    \proofstepwidestar[]{\IntSucc\colon\mathbb Z\inj\mathbb Z}{%
      \FormulaRefAuto{Injektive Funktion}{1,2}}
    \proofstepwidestar[]{%
      \forall y\in\mathbb Z\,\exists x\in\mathbb Z\,
      \IntSucc(x)=y}{\FormulaRefAuto{IntegerSuccessorSurjective}}
    \proofstepwidestar[]{\IntSucc\colon\mathbb Z\sur\mathbb Z}{%
      \FormulaRefAuto{Surjektive Funktion}{1,4}}
    \proofstepwidestar[]{\IntSucc\colon\mathbb Z\bij\mathbb Z}{%
      \FormulaRefAuto{F\colon A\inj B\dsep F\colon A\sur B
        \vdash F\colon A\bij B}{3,5}}
  \closeproofpartwide
\end{tabproofsplitwide}

\section{Normalformen}

\FormulaThmDeltaK[Positiver Strahl und Nachfolger]{%
  n\in\mathbb N\vdash
  \IntSucc(\IntEmbed(n))=\IntEmbed(\Succ(n))%
}{IntegerSuccessorOnEmbedding}{%
  \DeltaRow{Mengen}{n}%
}
\begin{tabproofwide}
  \proofstepwidestarbreakkeep[1]{n\in\mathbb N}{\rA}
  \proofstepwidestarbreakkeep[]{\IntOne=\IntEmbed(1)}{\FormulaRefAuto{IntegerOneDef}}
  \proofstepwidestarbreakkeep[1]{n+1=\Succ(n)}{\FormulaRefAuto{PeanoAddOneSucc}{1}}
  \proofstepwide[1]{\IntSucc(\IntEmbed(n))}{=}{%
    \IntEmbed(n)+\IntOne}{%
    \FormulaRefAuto{IntegerSuccessorValue}[thm]{%
      \FormulaRefAuto{F\colon A\to B\dsep x\in A\vdash F(x)\in B}{%
        \FormulaRefAuto{NaturalIntegerEmbeddingFunction},1}}}
  \proofstepwide[1]{}{=}{\IntEmbed(n)+\IntEmbed(1)}{\rIE{2}}
  \proofstepwide[1]{}{=}{\IntEmbed(n+1)}{%
    \FormulaRefAuto{a=b\vdash b=a}{%
      \FormulaRefAuto{NaturalIntegerEmbeddingAdditive}{1,%
        \FormulaRefAuto{PeanoOneInN}}}}
  \proofstepwide[1]{}{=}{\IntEmbed(\Succ(n))}{\rIE{3}}
  \proofstepwidestarbreakkeep[1]{%
    \IntSucc(\IntEmbed(n))=\IntEmbed(\Succ(n))}{\rChain{4,5,6,7}}
\end{tabproofwide}

\FormulaThmDeltaK[Negativer Strahl und Vorgänger]{%
  n\in\mathbb N\vdash
  \IntPred(-\IntEmbed(n))=-\IntEmbed(\Succ(n))%
}{IntegerPredecessorOnNegativeEmbedding}{%
  \DeltaRow{Mengen}{n}%
}
\begin{tabproofwide}
  \proofstepwidestar[1]{n\in\mathbb N}{\rA}
  \proofstepwidestar[1]{\IntEmbed(n)=\IntClass{n}{0}}{%
    \FormulaRefAuto{NaturalIntegerEmbeddingValue}[thm]{1}}
  \proofstepwidestar[1]{-\IntEmbed(n)=\IntClass{0}{n}}{%
    \FormulaRefAuto{IntegerNegationDef}{1,2}}
  \proofstepwidestar[1]{\IntEmbed(\Succ(n))=\IntClass{\Succ(n)}{0}}{%
    \FormulaRefAuto{NaturalIntegerEmbeddingValue}[thm]{%
      \FormulaRefAuto{PeanoSuccClosure}{1}}}
  \proofstepwidestar[1]{-\IntEmbed(\Succ(n))=\IntClass{0}{\Succ(n)}}{%
    \FormulaRefAuto{IntegerNegationDef}{1,4}}
  \proofstepwide[1]{\IntPred(-\IntEmbed(n))}{=}{%
    \IntClass{0}{\Succ(n)}}{%
    \FormulaRefAuto{IntegerPredecessorClassValue}{1,3}}
  \proofstepwide[1]{}{=}{-\IntEmbed(\Succ(n))}{%
    \FormulaRefAuto{a=b\vdash b=a}{5}}
  \proofstepwidestar[1]{%
    \IntPred(-\IntEmbed(n))=-\IntEmbed(\Succ(n))}{\rChain{6,7}}
\end{tabproofwide}

\FormulaThmDeltaK[Normalform ganzer Zahlen]{%
  z\in\mathbb Z\vdash
  \exists n\in\mathbb N\,
  \bigl(z=\IntEmbed(n)\lor z=-\IntEmbed(n)\bigr)%
}{IntegerSignedNormalForm}{%
  \DeltaRow{Mengen}{z\dsep a\dsep b\dsep n}%
}
\begin{tabproofsplitwide}
  \proofpartwideindR[Fall \(b\leq a\)]{%
    \begin{aligned}[t]
      &a,b\in\mathbb N\dsep b\leq a\vdash{}\\
      &\exists n\in\mathbb N\,
      \IntClass{a}{b}=\IntEmbed(n)
    \end{aligned}
  }{%
    a,b\in\mathbb N\dsep b\leq a
    \vdash\exists n\in\mathbb N\,
    \IntClass{a}{b}=\IntEmbed(n)
  }[IntegerNormalFormNonnegative]
    \proofstepwidestar[1]{a,b\in\mathbb N}{\rA}
    \proofstepwidestar[2]{b\leq a}{\rA}
    \proofstepwidestar[1,2]{\exists n\in\mathbb N\,(b+n=a)}{%
      \FormulaRefAuto{m\in\mathbb{N}\dsep n\in\mathbb{N}\vdash
        m\leq n\leftrightarrow
        \exists k\in\mathbb{N}\,(m+k=n)}{1,2}}
    \proofstepwidestar[4]{n\in\mathbb N\land b+n=a}{\rA}
    \proofstepwidestar[1,4]{a+0=b+n}{%
      \rIE{\FormulaRefAuto{PeanoAddRightZero}{1},\rAEb{4}}}
    \proofstepwidestar[1,4]{\IntClass{a}{b}=\IntClass{n}{0}}{%
      \FormulaRefAuto{IntegerClassEquality}{1,4,5}}
    \proofstepwidestar[4]{\IntEmbed(n)=\IntClass{n}{0}}{%
      \FormulaRefAuto{NaturalIntegerEmbeddingValue}[thm]{4}}
    \proofstepwidestar[1,4]{\IntClass{a}{b}=\IntEmbed(n)}{%
      \FormulaRefAuto{a=b\dsep c=b\vdash a=c}{6,7}}
    \proofstepwidestar[1,2]{\exists n\in\mathbb N\,
      \IntClass{a}{b}=\IntEmbed(n)}{\rEE{3,4,8}}
  \closeproofpartwideindR

  \proofpartwideindR[Fall \(a\leq b\)]{%
    \begin{aligned}[t]
      &a,b\in\mathbb N\dsep a\leq b\vdash{}\\
      &\exists n\in\mathbb N\,
      \IntClass{a}{b}=-\IntEmbed(n)
    \end{aligned}
  }{%
    a,b\in\mathbb N\dsep a\leq b
    \vdash\exists n\in\mathbb N\,
    \IntClass{a}{b}=-\IntEmbed(n)
  }[IntegerNormalFormNonpositive]
    \proofstepwidestar[1]{a,b\in\mathbb N}{\rA}
    \proofstepwidestar[2]{a\leq b}{\rA}
    \proofstepwidestar[1,2]{\exists n\in\mathbb N\,(a+n=b)}{%
      \FormulaRefAuto{m\in\mathbb{N}\dsep n\in\mathbb{N}\vdash
        m\leq n\leftrightarrow
        \exists k\in\mathbb{N}\,(m+k=n)}{1,2}}
    \proofstepwidestar[4]{n\in\mathbb N\land a+n=b}{\rA}
    \proofstepwidestar[1,4]{a+n=b+0}{%
      \rIE{\rAEb{4},\FormulaRefAuto{PeanoAddRightZero}{1}}}
    \proofstepwidestar[1,4]{\IntClass{a}{b}=\IntClass{0}{n}}{%
      \FormulaRefAuto{IntegerClassEquality}{1,4,5}}
    \proofstepwidestar[4]{\IntEmbed(n)=\IntClass{n}{0}}{%
      \FormulaRefAuto{NaturalIntegerEmbeddingValue}[thm]{4}}
    \proofstepwidestar[4]{-\IntEmbed(n)=\IntClass{0}{n}}{%
      \FormulaRefAuto{IntegerNegationDef}{4,7}}
    \proofstepwidestar[1,4]{\IntClass{a}{b}=-\IntEmbed(n)}{%
      \FormulaRefAuto{a=b\dsep c=b\vdash a=c}{6,8}}
    \proofstepwidestar[1,2]{\exists n\in\mathbb N\,
      \IntClass{a}{b}=-\IntEmbed(n)}{\rEE{3,4,9}}
  \closeproofpartwideindR

  \proofpartwide{Schluss}
    \proofstepwidestar[1]{z\in\mathbb Z}{\rA}
    \proofstepwidestar[1]{\exists a,b\in\mathbb N\,
      z=\IntClass{a}{b}}{\FormulaRefAuto{IntegerPairRepresentative}{1}}
    \proofstepwidestar[3]{a,b\in\mathbb N\land z=\IntClass{a}{b}}{\rA}
    \proofstepwidestar[3]{b\leq a\lor a\leq b}{%
      \FormulaRefAuto{PeanoOrderComparison}{3}}
    \proofstepwidestar[5]{b\leq a}{\rA}
    \proofstepwidestar[3,5]{\exists n\in\mathbb N\,
      \IntClass{a}{b}=\IntEmbed(n)}{%
      \FormulaRefAuto{IntegerNormalFormNonnegative}{3,5}}
    \proofstepwidestar[3,5]{\exists n\in\mathbb N\,
      (z=\IntEmbed(n)\lor z=-\IntEmbed(n))}{\rIE{3,6}}
    \proofstepwidestar[8]{a\leq b}{\rA}
    \proofstepwidestar[3,8]{\exists n\in\mathbb N\,
      \IntClass{a}{b}=-\IntEmbed(n)}{%
      \FormulaRefAuto{IntegerNormalFormNonpositive}{3,8}}
    \proofstepwidestar[3,8]{\exists n\in\mathbb N\,
      (z=\IntEmbed(n)\lor z=-\IntEmbed(n))}{\rIE{3,9}}
    \proofstepwidestar[3]{\exists n\in\mathbb N\,
      (z=\IntEmbed(n)\lor z=-\IntEmbed(n))}{%
      \rOE{4,5,7,8,10}}
    \proofstepwidestar[1]{\exists n\in\mathbb N\,
      (z=\IntEmbed(n)\lor z=-\IntEmbed(n))}{\rEE{2,3,11}}
  \closeproofpartwide
\end{tabproofsplitwide}

\chapter{Ordnung der ganzen Zahlen}

\section{Translationsinvarianz der natürlichen Ordnung}

\FormulaThmDeltaK[Translationsinvarianz der natürlichen Ordnung]{%
  a,b,c\in\mathbb N\vdash
  a\leq b\leftrightarrow a+c\leq b+c%
}{NaturalOrderTranslationForIntegers}{%
  \DeltaRow{Mengen}{a\dsep b\dsep c\dsep k}%
}
\begin{tabproofsplitwide}
  \proofpartwideindR[Erhaltung unter Translation]{%
    a,b,c\in\mathbb N\dsep a\leq b
    \vdash a+c\leq b+c%
  }{%
    a,b,c\in\mathbb N\dsep a\leq b
    \vdash a+c\leq b+c%
  }[NaturalOrderTranslationForward]
    \proofstepwidestar[1]{a,b,c\in\mathbb N}{\rA}
    \proofstepwidestar[2]{a\leq b}{\rA}
    \proofstepwidestar[1,2]{\exists k\in\mathbb N\,(a+k=b)}{%
      \FormulaRefAuto{m\in\mathbb{N}\dsep n\in\mathbb{N}\vdash
        m\leq n\leftrightarrow
        \exists k\in\mathbb{N}\,(m+k=n)}{1,2}}
    \proofstepwidestar[4]{k\in\mathbb N\land a+k=b}{\rA}
    \proofstepwide[1,4]{(a+c)+k}{=}{(a+k)+c}{%
      \FormulaRefAuto{PeanoAddInterchange}{1}}
    \proofstepwide[1,4]{(a+k)+c}{=}{b+c}{\rIE{\rAEb{4}}}
    \proofstepwidestar[1,4]{(a+c)+k=b+c}{%
      \rChain{5,6}}
    \proofstepwidestar[1,4]{a+c\leq b+c}{%
      \FormulaRefAuto{m\in\mathbb{N}\dsep n\in\mathbb{N}\vdash
        m\leq n\leftrightarrow
        \exists k\in\mathbb{N}\,(m+k=n)}{1,\rEI{\rAI{\rAEa{4},7}}}}
    \proofstepwidestar[1,2]{a+c\leq b+c}{\rEE{3,4,8}}
  \closeproofpartwideindR

  \proofpartwideindR[Kürzen einer Translation]{%
    a,b,c\in\mathbb N\dsep a+c\leq b+c
    \vdash a\leq b%
  }{%
    a,b,c\in\mathbb N\dsep a+c\leq b+c
    \vdash a\leq b%
  }[NaturalOrderTranslationBackward]
    \proofstepwidestar[1]{a,b,c\in\mathbb N}{\rA}
    \proofstepwidestar[2]{a+c\leq b+c}{\rA}
    \proofstepwidestar[1,2]{\exists k\in\mathbb N\,((a+c)+k=b+c)}{%
      \FormulaRefAuto{m\in\mathbb{N}\dsep n\in\mathbb{N}\vdash
        m\leq n\leftrightarrow
        \exists k\in\mathbb{N}\,(m+k=n)}{1,2}}
    \proofstepwidestar[4]{k\in\mathbb N\land(a+c)+k=b+c}{\rA}
    \proofstepwide[1,4]{(a+k)+c}{=}{(a+c)+k}{%
      \FormulaRefAuto{PeanoAddInterchange}{1}}
    \proofstepwide[1,4]{(a+c)+k}{=}{b+c}{\rAEb{4}}
    \proofstepwidestar[1,4]{(a+k)+c=b+c}{%
      \rChain{5,6}}
    \proofstepwidestar[1,4]{a+k=b}{%
      \FormulaRefAuto{PeanoAddRightCancellation}{%
        \FormulaRefAuto{PeanoAddClosure}{1,4},1,1,7}}
    \proofstepwidestar[1,4]{a\leq b}{%
      \FormulaRefAuto{m\in\mathbb{N}\dsep n\in\mathbb{N}\vdash
        m\leq n\leftrightarrow
        \exists k\in\mathbb{N}\,(m+k=n)}{1,\rEI{\rAI{\rAEa{4},8}}}}
    \proofstepwidestar[1,2]{a\leq b}{\rEE{3,4,9}}
  \closeproofpartwideindR

  \proofpartwide{Schluss}
    \proofstepwidestar[1]{a,b,c\in\mathbb N}{\rA}
    \proofstepwidestar[1]{a\leq b\rightarrow a+c\leq b+c}{%
      \FormulaRefAuto{NaturalOrderTranslationForward}{1}}
    \proofstepwidestar[1]{a+c\leq b+c\rightarrow a\leq b}{%
      \FormulaRefAuto{NaturalOrderTranslationBackward}{1}}
    \proofstepwidestar[1]{a\leq b\leftrightarrow a+c\leq b+c}{%
      \rLRI{2,3}}
  \closeproofpartwide
\end{tabproofsplitwide}

\FormulaThmDeltaK[Addition natürlicher Ungleichungen]{%
  a,b,c,d\in\mathbb N\dsep a\leq b\dsep c\leq d
  \vdash a+c\leq b+d%
}{NaturalOrderAdditionForIntegers}{%
  \DeltaRow{Mengen}{a\dsep b\dsep c\dsep d}%
}
\begin{tabproofsplitwide}
  \proofpartwideindR[Translation der ersten Ungleichung]{%
    a,b,c\in\mathbb N\dsep a\leq b
    \vdash a+c\leq b+c%
  }{%
    a,b,c\in\mathbb N\dsep a\leq b
    \vdash a+c\leq b+c%
  }[NaturalOrderAdditionFirstTranslation]
    \proofstepwidestar[1]{a,b,c\in\mathbb N}{\rA}
    \proofstepwidestar[2]{a\leq b}{\rA}
    \proofstepwidestar[1,2]{a+c\leq b+c}{%
      \FormulaRefAuto{NaturalOrderTranslationForward}{1,2}}
  \closeproofpartwideindR

  \proofpartwideindR[Translation der zweiten Ungleichung]{%
    b,c,d\in\mathbb N\dsep c\leq d
    \vdash b+c\leq b+d%
  }{%
    b,c,d\in\mathbb N\dsep c\leq d
    \vdash b+c\leq b+d%
  }[NaturalOrderAdditionSecondTranslation]
    \proofstepwidestar[1]{b,c,d\in\mathbb N}{\rA}
    \proofstepwidestar[2]{c\leq d}{\rA}
    \proofstepwidestar[1,2]{c+b\leq d+b}{%
      \FormulaRefAuto{NaturalOrderTranslationForward}{1,2}}
    \proofstepwidestar[1]{b+c=c+b}{%
      \FormulaRefAuto{PeanoAddCommutative}{1}}
    \proofstepwidestar[1]{d+b=b+d}{%
      \FormulaRefAuto{PeanoAddCommutative}{1}}
    \proofstepwidestar[1,2]{b+c\leq b+d}{\rIE{4,3,5}}
  \closeproofpartwideindR

  \proofpartwideindR[Verknüpfung durch Transitivität]{%
    a,b,c,d\in\mathbb N\dsep a\leq b\dsep c\leq d
    \vdash a+c\leq b+d%
  }{%
    a,b,c,d\in\mathbb N\dsep a\leq b\dsep c\leq d
    \vdash a+c\leq b+d%
  }[NaturalOrderAdditionTransitiveConclusion]
    \proofstepwidestar[1]{a,b,c,d\in\mathbb N}{\rA}
    \proofstepwidestar[2]{a\leq b}{\rA}
    \proofstepwidestar[3]{c\leq d}{\rA}
    \proofstepwidestar[1,2]{a+c\leq b+c}{%
      \FormulaRefAuto{NaturalOrderAdditionFirstTranslation}{1,2}}
    \proofstepwidestar[1,3]{b+c\leq b+d}{%
      \FormulaRefAuto{NaturalOrderAdditionSecondTranslation}{1,3}}
    \proofstepwidestar[1]{a+c\in\mathbb N}{%
      \FormulaRefAuto{PeanoAddClosure}{1}}
    \proofstepwidestar[1]{b+c\in\mathbb N}{%
      \FormulaRefAuto{PeanoAddClosure}{1}}
    \proofstepwidestar[1]{b+d\in\mathbb N}{%
      \FormulaRefAuto{PeanoAddClosure}{1}}
    \proofstepwidestar[1,2,3]{a+c\leq b+d}{%
      \FormulaRefAuto{%
        k\in\mathbb{N}\dsep m\in\mathbb{N}\dsep n\in\mathbb{N}\dsep
        k\leq m\dsep m\leq n\vdash k\leq n%
      }{6,7,8,4,5}}
  \closeproofpartwideindR
\end{tabproofsplitwide}

\section{Quotientenordnung}

\FormulaThmDeltaK[Repräsentantenunabhängigkeit der Ganzzahlordnung]{%
  \begin{aligned}[t]
    &a,b,a',b',c,d,c',d'\in\mathbb N\dsep
      \IntClass{a}{b}=\IntClass{a'}{b'}\dsep
      \IntClass{c}{d}=\IntClass{c'}{d'}\vdash{}\\
    &(a+d\leq c+b)\leftrightarrow(a'+d'\leq c'+b')
  \end{aligned}%
}{IntegerOrderCompatible}{%
  \DeltaRow{Mengen}{a\dsep b\dsep a'\dsep b'\dsep
    c\dsep d\dsep c'\dsep d'}%
}
\begin{tabproofsplitwide}
  \proofpartwideindR[Hinrichtung]{%
    \begin{aligned}[t]
      &a,b,a',b',c,d,c',d'\in\mathbb N\dsep
      a+b'=b+a'\dsep c+d'=d+c'\dsep a+d\leq c+b\vdash{}\\
      &a'+d'\leq c'+b'
    \end{aligned}
  }{%
    a,b,a',b',c,d,c',d'\in\mathbb N\dsep
    a+b'=b+a'\dsep c+d'=d+c'\dsep a+d\leq c+b
    \vdash a'+d'\leq c'+b'
  }[IntegerOrderCompatibleForward]
    \proofstepwidestar[1]{a,b,a',b',c,d,c',d'\in\mathbb N}{\rA}
    \proofstepwidestar[2]{a+b'=b+a'}{\rA}
    \proofstepwidestar[3]{c+d'=d+c'}{\rA}
    \proofstepwidestar[4]{a+d\leq c+b}{\rA}
    \proofstepwidestar[1,4]{(a+d)+(b'+c')\leq(c+b)+(b'+c')}{%
      \FormulaRefAuto{NaturalOrderTranslationForward}{1,4}}
    \proofstepwidestar[1,2,3]{%
      (a+d)+(b'+c')=(a'+d')+(b+c)}{%
      \FormulaRefAuto{PeanoAddInterchange}{1,2,3}}
    \proofstepwidestar[1]{%
      (c+b)+(b'+c')=(c'+b')+(b+c)}{%
      \FormulaRefAuto{PeanoAddInterchange}{1}}
    \proofstepwidestar[1,2,3,4]{%
      (a'+d')+(b+c)\leq(c'+b')+(b+c)}{\rIE{5,6,7}}
    \proofstepwidestar[1,2,3,4]{a'+d'\leq c'+b'}{%
      \FormulaRefAuto{NaturalOrderTranslationBackward}{1,8}}
  \closeproofpartwideindR

  \proofpartwideindR[Rückrichtung]{%
    \begin{aligned}[t]
      &a,b,a',b',c,d,c',d'\in\mathbb N\dsep
      a+b'=b+a'\dsep c+d'=d+c'\dsep a'+d'\leq c'+b'\vdash{}\\
      &a+d\leq c+b
    \end{aligned}
  }{%
    a,b,a',b',c,d,c',d'\in\mathbb N\dsep
    a+b'=b+a'\dsep c+d'=d+c'\dsep a'+d'\leq c'+b'
    \vdash a+d\leq c+b
  }[IntegerOrderCompatibleBackward]
    \proofstepwidestar[1]{a,b,a',b',c,d,c',d'\in\mathbb N}{\rA}
    \proofstepwidestar[2]{a+b'=b+a'}{\rA}
    \proofstepwidestar[3]{c+d'=d+c'}{\rA}
    \proofstepwidestar[4]{a'+d'\leq c'+b'}{\rA}
    \proofstepwidestar[1,2]{a'+b=b'+a}{%
      \FormulaRefAuto{a=b\vdash b=a}{%
        \FormulaRefAuto{PeanoAddCommutative}{2}}}
    \proofstepwidestar[1,3]{c'+d=d'+c}{%
      \FormulaRefAuto{a=b\vdash b=a}{%
        \FormulaRefAuto{PeanoAddCommutative}{3}}}
    \proofstepwidestar[1,2,3,4]{a+d\leq c+b}{%
      \FormulaRefAuto{IntegerOrderCompatibleForward}{1,5,6,4}}
  \closeproofpartwideindR

  \proofpartwide{Schluss}
    \proofstepwidestar[1]{a,b,a',b',c,d,c',d'\in\mathbb N}{\rA}
    \proofstepwidestar[2]{\IntClass{a}{b}=\IntClass{a'}{b'}}{\rA}
    \proofstepwidestar[3]{\IntClass{c}{d}=\IntClass{c'}{d'}}{\rA}
    \proofstepwidestar[1,2]{a+b'=b+a'}{%
      \FormulaRefAuto{IntegerClassEquality}{1,2}}
    \proofstepwidestar[1,3]{c+d'=d+c'}{%
      \FormulaRefAuto{IntegerClassEquality}{1,3}}
    \proofstepwidestar[1,2,3]{%
      a+d\leq c+b\rightarrow a'+d'\leq c'+b'}{%
      \FormulaRefAuto{IntegerOrderCompatibleForward}{1,4,5}}
    \proofstepwidestar[1,2,3]{%
      a'+d'\leq c'+b'\rightarrow a+d\leq c+b}{%
      \FormulaRefAuto{IntegerOrderCompatibleBackward}{1,4,5}}
    \proofstepwidestar[1,2,3]{%
      (a+d\leq c+b)\leftrightarrow(a'+d'\leq c'+b')}{\rLRI{6,7}}
  \closeproofpartwide
\end{tabproofsplitwide}

\FormulaDefDeltaK[Ordnungsoperator der ganzen Zahlen]{%
  a,b,c,d\in\mathbb N\vdash
  \IntClass{a}{b}\leq_{\mathbb Z}\IntClass{c}{d}
  \coloneqq a+d\leq c+b%
}{IntegerOrderDef}{%
  \DeltaRow{Mengen}{a\dsep b\dsep c\dsep d}%
  \DeltaRow{Relationsoperatoren}{\leq_{\mathbb Z}}%
  \DeltaRow{Neue Symbole}{\leq_{\mathbb Z}}%
}

\FormulaDefDeltaK[Notation der Ganzzahlordnung]{%
  z,w\in\mathbb Z\vdash
  z\leq w\coloneqq z\leq_{\mathbb Z}w%
}{IntegerOrderNotation}{%
  \DeltaRow{Mengen}{z\dsep w}%
  \DeltaRow{Relationsoperatoren}{\leq\dsep\leq_{\mathbb Z}}%
}

\FormulaThmDeltaK[Klassencharakterisierung der Ganzzahlordnung]{%
  a,b,c,d\in\mathbb N\vdash
  \IntClass{a}{b}\leq\IntClass{c}{d}
  \leftrightarrow a+d\leq c+b%
}{IntegerOrderClassChar}{%
  \DeltaRow{Mengen}{a\dsep b\dsep c\dsep d}%
  \DeltaRow{Relationsoperatoren}{\leq\dsep\leq_{\mathbb Z}}%
}
\begin{tabproof}
  \proofstep{1}{a,b,c,d\in\mathbb N}{\rA}
  \proofstep{1}{\IntClass{a}{b}\in\mathbb Z}{%
    \FormulaRefAuto{IntegerClassInZ}[thm]{1}}
  \proofstep{1}{\IntClass{c}{d}\in\mathbb Z}{%
    \FormulaRefAuto{IntegerClassInZ}[thm]{1}}
  \proofstep{1}{%
    \IntClass{a}{b}\leq\IntClass{c}{d}
    \leftrightarrow
    \IntClass{a}{b}\leq_{\mathbb Z}\IntClass{c}{d}}{%
    \FormulaRefAuto{IntegerOrderNotation}[def]{2,3}}
  \proofstep{1}{%
    \IntClass{a}{b}\leq_{\mathbb Z}\IntClass{c}{d}
    \leftrightarrow a+d\leq c+b}{%
    \FormulaRefAuto{IntegerOrderDef}[def]{1}}
  \proofstep{1}{%
    \IntClass{a}{b}\leq\IntClass{c}{d}
    \leftrightarrow a+d\leq c+b}{%
    \FormulaRefAuto{P\leftrightarrow Q\dsep Q\leftrightarrow R
      \vdash P\leftrightarrow R}{4,5}}
\end{tabproof}

\FormulaThmDeltaK[Totale Ordnung der ganzen Zahlen]{%
  \TotOrd{\mathbb Z,\leq_{\mathbb Z}}%
}{IntegerOrderTotal}{%
  \DeltaRow{Mengen}{z\dsep w\dsep u}%
  \DeltaRow{Relationsoperatoren}{\leq_{\mathbb Z}}%
}
\begin{tabproofsplitwide}
  \proofpartwideindR[Reflexivität]{%
    z\in\mathbb Z\vdash z\leq z%
  }{%
    z\in\mathbb Z\vdash z\leq z%
  }[IntegerOrderReflexive]
    \proofstepwidestarbreakkeep[1]{z\in\mathbb Z}{\rA}
    \proofstepwidestarbreakkeep[1]{\exists a,b\in\mathbb N\,
      z=\IntClass{a}{b}}{\FormulaRefAuto{IntegerPairRepresentative}{1}}
    \proofstepwidestarbreakkeep[3]{a,b\in\mathbb N\land z=\IntClass{a}{b}}{\rA}
    \proofstepwidestarbreakkeep[3]{a+b\leq a+b}{%
      \FormulaRefAuto{PeanoLeqReflexive}{%
        \FormulaRefAuto{PeanoAddClosure}{3}}}
    \proofstepwidestarbreakkeep[3]{\IntClass{a}{b}\leq\IntClass{a}{b}}{%
      \FormulaRefAuto{IntegerOrderClassChar}{3,4}}
    \proofstepwidestarbreakkeep[3]{z\leq z}{\rIE{3,5}}
    \proofstepwidestarbreakkeep[1]{z\leq z}{\rEE{2,3,6}}
  \closeproofpartwideindR

  \proofpartwideindR[Antisymmetrie]{%
    z,w\in\mathbb Z\dsep z\leq w\dsep w\leq z
    \vdash z=w%
  }{%
    z,w\in\mathbb Z\dsep z\leq w\dsep w\leq z
    \vdash z=w%
  }[IntegerOrderAntisymmetric]
    \proofstepwidestarbreakkeep[1]{z,w\in\mathbb Z}{\rA}
    \proofstepwidestarbreakkeep[2]{z\leq w}{\rA}
    \proofstepwidestarbreakkeep[3]{w\leq z}{\rA}
    \proofstepwidestarbreakkeep[1]{z\in\mathbb Z}{\rAEa{1}}
    \proofstepwidestarbreakkeep[1]{w\in\mathbb Z}{\rAEb{1}}
    \proofstepwidestarbreakkeep[1]{\exists a,b\in\mathbb N\,
      z=\IntClass{a}{b}}{\FormulaRefAuto{IntegerPairRepresentative}{4}}
    \proofstepwidestarbreakkeep[1]{\exists c,d\in\mathbb N\,
      w=\IntClass{c}{d}}{\FormulaRefAuto{IntegerPairRepresentative}{5}}
    \proofstepwidestarbreakkeep[8]{a,b,c,d\in\mathbb N\land
      z=\IntClass{a}{b}\land w=\IntClass{c}{d}}{\rA}
    \proofstepwidestarbreakkeep[2,8]{a+d\leq c+b}{%
      \FormulaRefAuto{IntegerOrderClassChar}{8,2}}
    \proofstepwidestarbreakkeep[3,8]{c+b\leq a+d}{%
      \FormulaRefAuto{IntegerOrderClassChar}{8,3}}
    \proofstepwidestarbreakkeep[8]{a+d\in\mathbb N}{%
      \FormulaRefAuto{PeanoAddClosure}{8}}
    \proofstepwidestarbreakkeep[8]{c+b\in\mathbb N}{%
      \FormulaRefAuto{PeanoAddClosure}{8}}
    \proofstepwidestarbreakkeep[2,3,8]{a+d=c+b}{%
      \FormulaRefAuto{PeanoLeqAntisymmetric}[thm]{11,12,9,10}}
    \proofstepwidestarbreakkeep[2,3,8]{%
      \IntClass{a}{b}=\IntClass{c}{d}}{%
      \FormulaRefAuto{IntegerClassEquality}{8,13}}
    \proofstepwidestarbreakkeep[2,3,8]{z=w}{\rIE{8,14}}
    \proofstepwidestarbreakkeep[1,2,3]{z=w}{\rEE{6,7,8,15}}
  \closeproofpartwideindR

  \proofpartwideindR[Transitivität]{%
    z,w,u\in\mathbb Z\dsep z\leq w\dsep w\leq u
    \vdash z\leq u%
  }{%
    z,w,u\in\mathbb Z\dsep z\leq w\dsep w\leq u
    \vdash z\leq u%
  }[IntegerOrderTransitive]
    \proofstepwidestarbreakkeep[1]{z,w,u\in\mathbb Z}{\rA}
    \proofstepwidestarbreakkeep[2]{z\leq w}{\rA}
    \proofstepwidestarbreakkeep[3]{w\leq u}{\rA}
    \proofstepwidestarbreakkeep[1]{z\in\mathbb Z}{\rAEa{1}}
    \proofstepwidestarbreakkeep[1]{w\in\mathbb Z}{\rAEa{\rAEb{1}}}
    \proofstepwidestarbreakkeep[1]{u\in\mathbb Z}{\rAEb{\rAEb{1}}}
    \proofstepwidestarbreakkeep[4]{\exists a,b\in\mathbb N\,
      z=\IntClass{a}{b}}{%
      \FormulaRefAuto{IntegerPairRepresentative}{4}}
    \proofstepwidestarbreakkeep[5]{\exists c,d\in\mathbb N\,
      w=\IntClass{c}{d}}{%
      \FormulaRefAuto{IntegerPairRepresentative}{5}}
    \proofstepwidestarbreakkeep[6]{\exists e,f\in\mathbb N\,
      u=\IntClass{e}{f}}{%
      \FormulaRefAuto{IntegerPairRepresentative}{6}}
    \proofstepwidestarbreakkeep[10]{a,b,c,d,e,f\in\mathbb N\land
      z=\IntClass{a}{b}\land w=\IntClass{c}{d}\land
      u=\IntClass{e}{f}}{\rA}
    \proofstepwidestarbreakkeep[2,10]{a+d\leq c+b}{%
      \FormulaRefAuto{IntegerOrderClassChar}{10,2}}
    \proofstepwidestarbreakkeep[3,10]{c+f\leq e+d}{%
      \FormulaRefAuto{IntegerOrderClassChar}{10,3}}
    \proofstepwidestarbreakkeep[2,3,10]{%
      (a+d)+(c+f)\leq(c+b)+(e+d)}{%
      \FormulaRefAuto{NaturalOrderAdditionForIntegers}{10,11,12}}
    \proofstepwidestarbreakkeep[10]{%
      (a+d)+(c+f)=(a+f)+(c+d)}{%
      \FormulaRefAuto{PeanoAddInterchange}{10}}
    \proofstepwidestarbreakkeep[10]{%
      (c+b)+(e+d)=(b+e)+(c+d)}{%
      \FormulaRefAuto{PeanoAddInterchange}{10}}
    \proofstepwidestarbreakkeep[2,3,10]{(a+f)+(c+d)\leq(b+e)+(c+d)}{%
      \rIE{13,14,15}}
    \proofstepwidestarbreakkeep[2,3,10]{a+f\leq b+e}{%
      \FormulaRefAuto{NaturalOrderTranslationBackward}{10,16}}
    \proofstepwidestarbreakkeep[2,3,10]{z\leq u}{%
      \FormulaRefAuto{IntegerOrderClassChar}{10,17}}
    \proofstepwidestarbreakkeep[1,2,3]{z\leq u}{\rEE{7,8,9,10,18}}
  \closeproofpartwideindR

  \proofpartwideindR[Totalität]{%
    z,w\in\mathbb Z\vdash z\leq w\lor w\leq z%
  }{%
    z,w\in\mathbb Z\vdash z\leq w\lor w\leq z%
  }[IntegerOrderComparable]
    \proofstepwidestarbreakkeep[1]{z,w\in\mathbb Z}{\rA}
    \proofstepwidestarbreakkeep[1]{z\in\mathbb Z}{\rAEa{1}}
    \proofstepwidestarbreakkeep[1]{w\in\mathbb Z}{\rAEb{1}}
    \proofstepwidestarbreakkeep[2]{\exists a,b\in\mathbb N\,
      z=\IntClass{a}{b}}{%
      \FormulaRefAuto{IntegerPairRepresentative}{2}}
    \proofstepwidestarbreakkeep[3]{\exists c,d\in\mathbb N\,
      w=\IntClass{c}{d}}{%
      \FormulaRefAuto{IntegerPairRepresentative}{3}}
    \proofstepwidestarbreakkeep[6]{a,b,c,d\in\mathbb N\land
      z=\IntClass{a}{b}\land w=\IntClass{c}{d}}{\rA}
    \proofstepwidestarbreakkeep[6]{a+d\leq c+b\lor c+b\leq a+d}{%
      \FormulaRefAuto{PeanoOrderComparison}{%
        \FormulaRefAuto{PeanoAddClosure}{6},%
        \FormulaRefAuto{PeanoAddClosure}{6}}}
    \proofstepwidestarbreakkeep[6]{z\leq w\lor w\leq z}{%
      \FormulaRefAuto{IntegerOrderClassChar}{6,7}}
    \proofstepwidestarbreakkeep[1]{z\leq w\lor w\leq z}{\rEE{4,5,6,8}}
  \closeproofpartwideindR

  \proofpartwide{Schluss}
    \proofstepwidestarbreakkeep[]{\forall z\in\mathbb Z\,(z\leq z)}{%
      \FormulaRefAuto{IntegerOrderReflexive}}
    \proofstepwidestarbreakkeep[2]{z\in\mathbb Z}{\rA}
    \proofstepwidestarbreakkeep[3]{w\in\mathbb Z}{\rA}
    \proofstepwidestarbreakkeep[4]{u\in\mathbb Z}{\rA}
    \proofstepwidestarbreakkeep[5]{z\leq w}{\rA}
    \proofstepwidestarbreakkeep[6]{w\leq u}{\rA}
    \proofstepwidestarbreakkeep[2,3,4,5,6]{z\leq u}{%
      \FormulaRefAuto{IntegerOrderTransitive}[thm]{2,3,4,5,6}}
    \proofstepwidestarbreakkeep[8]{w\leq z}{\rA}
    \proofstepwidestarbreakkeep[2,3,5,8]{z=w}{%
      \FormulaRefAuto{IntegerOrderAntisymmetric}[thm]{2,3,5,8}}
    \proofstepwidestarbreakkeep[]{\PartOrd{\mathbb Z,\leq_{\mathbb Z}}}{%
      \FormulaRefAuto{Partielle Ordnung}[def]{1,7,9}}
    \proofstepwidestarbreakkeep[2,3]{z\leq w\lor w\leq z}{%
      \FormulaRefAuto{IntegerOrderComparable}[thm]{2,3}}
    \proofstepwidestarbreakkeep[]{\TotOrd{\mathbb Z,\leq_{\mathbb Z}}}{%
      \FormulaRefAuto{Totale Ordnung}[def]{10,11}}
  \closeproofpartwide
\end{tabproofsplitwide}

Die Quotientenordnung ist nun vollständig am Modell verifiziert. Damit können
die für die spätere Axiomatik bestimmten Regeln ohne Vorgriff um die Ordnung
erweitert werden. Der Abstraktionsschnitt ist damit noch nicht erreicht;
zunächst werden auch Positivität, Diskretheit und Induktion am Modell
hergeleitet.

\FormulaAxiomAutoR[Totale Ordnung der Ganzzahlarithmetik]{%
  \TotOrd{\mathbb Z,\leq_{\mathbb Z}}%
}{IntegerOrderTotalAxiom}
\par\noindent\emph{Herkunft aus der Quotientenkonstruktion:}
\FormulaRefAuto{IntegerOrderTotal}[thm].\par

\FormulaAxiomAutoR[Reflexivität der Ganzzahlordnung]{%
  z\in\mathbb Z\vdash z\leq z%
}{IntegerOrderReflexiveAxiom}
\par\noindent\emph{Herkunft aus der Quotientenkonstruktion:}
\FormulaRefAuto{IntegerOrderReflexive}[thm].\par

\FormulaAxiomAutoR[Antisymmetrie der Ganzzahlordnung]{%
  z,w\in\mathbb Z\dsep z\leq w\dsep w\leq z\vdash z=w%
}{IntegerOrderAntisymmetricAxiom}
\par\noindent\emph{Herkunft aus der Quotientenkonstruktion:}
\FormulaRefAuto{IntegerOrderAntisymmetric}[thm].\par

\FormulaAxiomAutoR[Transitivität der Ganzzahlordnung]{%
  z,w,u\in\mathbb Z\dsep z\leq w\dsep w\leq u\vdash z\leq u%
}{IntegerOrderTransitiveAxiom}
\par\noindent\emph{Herkunft aus der Quotientenkonstruktion:}
\FormulaRefAuto{IntegerOrderTransitive}[thm].\par

\FormulaAxiomAutoR[Vergleichbarkeit ganzer Zahlen]{%
  z,w\in\mathbb Z\vdash z\leq w\lor w\leq z%
}{IntegerOrderComparableAxiom}
\par\noindent\emph{Herkunft aus der Quotientenkonstruktion:}
\FormulaRefAuto{IntegerOrderComparable}[thm].\par

\FormulaDefDeltaK[Strikte Ordnung der ganzen Zahlen]{%
  z,w\in\mathbb Z\vdash
  z<w\coloneqq z\leq w\land z\neq w%
}{IntegerStrictOrderDef}{%
  \DeltaRow{Mengen}{z\dsep w}%
  \DeltaRow{Relationsoperatoren}{<}%
}

\FormulaThmDeltaK[Translationsäquivalenz im Quotientenmodell]{%
  z,w,u\in\mathbb Z\vdash
  \bigl(z\leq w\leftrightarrow z+u\leq w+u\bigr)%
}{IntegerOrderTranslationModel}{%
  \DeltaRow{Mengen}{z\dsep w\dsep u\dsep
    a\dsep b\dsep c\dsep d\dsep e\dsep f}%
}
\begin{tabproofsplitwide}
  \proofpartwideindR[Rechnung auf drei festen Repräsentanten]{%
    \begin{aligned}[t]
      &a,b,c,d,e,f\in\mathbb N\vdash{}\\
      &\IntClass{a}{b}\leq\IntClass{c}{d}
       \leftrightarrow
       \IntClass{a}{b}+\IntClass{e}{f}
       \leq\IntClass{c}{d}+\IntClass{e}{f}
    \end{aligned}%
  }{%
    \begin{aligned}[t]
      &a,b,c,d,e,f\in\mathbb N\vdash{}\\
      &\IntClass{a}{b}\leq\IntClass{c}{d}
       \leftrightarrow
       \IntClass{a}{b}+\IntClass{e}{f}
       \leq\IntClass{c}{d}+\IntClass{e}{f}
    \end{aligned}%
  }[IntegerOrderTranslationOnClasses]
    \proofstepwidestarbreakkeep[1]{a,b,c,d,e,f\in\mathbb N}{\rA}
    \proofstepwidestarbreakkeep[1]{%
      \IntClass{a}{b}+\IntClass{e}{f}=\IntClass{a+e}{b+f}}{%
      \FormulaRefAuto{IntegerAdditionDef}{1}}
    \proofstepwidestarbreakkeep[1]{%
      \IntClass{c}{d}+\IntClass{e}{f}=\IntClass{c+e}{d+f}}{%
      \FormulaRefAuto{IntegerAdditionDef}{1}}
    \proofstepwidestarbreakkeep[1]{%
      \IntClass{a}{b}\leq\IntClass{c}{d}
      \leftrightarrow a+d\leq c+b}{%
      \FormulaRefAuto{IntegerOrderClassChar}{1}}
    \proofstepwidestarbreakkeep[1]{%
      \IntClass{a+e}{b+f}\leq\IntClass{c+e}{d+f}
      \leftrightarrow(a+e)+(d+f)\leq(c+e)+(b+f)}{%
      \FormulaRefAuto{IntegerOrderClassChar}{1}}
    \proofstepwidestarbreakkeep[1]{(a+e)+(d+f)=(a+d)+(e+f)}{%
      \FormulaRefAuto{PeanoAddInterchange}{1}}
    \proofstepwidestarbreakkeep[1]{(c+e)+(b+f)=(c+b)+(e+f)}{%
      \FormulaRefAuto{PeanoAddInterchange}{1}}
    \proofstepwidestarbreakkeep[1]{%
      a+d\leq c+b\leftrightarrow
      (a+d)+(e+f)\leq(c+b)+(e+f)}{%
      \FormulaRefAuto{NaturalOrderTranslationForIntegers}{1}}
    \proofstepwidestarbreakkeep[1]{%
      \IntClass{a}{b}+\IntClass{e}{f}
      \leq\IntClass{c}{d}+\IntClass{e}{f}
      \leftrightarrow
      (a+d)+(e+f)\leq(c+b)+(e+f)}{\rIE{2,3,5,6,7}}
    \proofstepwidestarbreakkeep[1]{%
      \IntClass{a}{b}\leq\IntClass{c}{d}
      \leftrightarrow
      \IntClass{a}{b}+\IntClass{e}{f}
      \leq\IntClass{c}{d}+\IntClass{e}{f}}{%
      \FormulaRefAuto{P\leftrightarrow Q\dsep Q\leftrightarrow R
        \vdash P\leftrightarrow R}{%
        \FormulaRefAuto{P\leftrightarrow Q\dsep Q\leftrightarrow R
          \vdash P\leftrightarrow R}{4,8},%
        \FormulaRefAuto{P\leftrightarrow Q\vdash Q\leftrightarrow P}{9}}}
  \closeproofpartwideindR

  \proofpartwide{Elimination der Repräsentanten}
    \proofstepwidestarbreakkeep[1]{z,w,u\in\mathbb Z}{\rA}
    \proofstepwidestarbreakkeep[1]{\exists a,b\in\mathbb N\,
      z=\IntClass{a}{b}}{%
      \FormulaRefAuto{IntegerPairRepresentative}{\rAEa{1}}}
    \proofstepwidestarbreakkeep[1]{\exists c,d\in\mathbb N\,
      w=\IntClass{c}{d}}{%
      \FormulaRefAuto{IntegerPairRepresentative}{\rAEa{\rAEb{1}}}}
    \proofstepwidestarbreakkeep[1]{\exists e,f\in\mathbb N\,
      u=\IntClass{e}{f}}{%
      \FormulaRefAuto{IntegerPairRepresentative}{\rAEb{\rAEb{1}}}}
    \proofstepwidestarbreakkeep[5]{a,b,c,d,e,f\in\mathbb N\land
      z=\IntClass{a}{b}\land w=\IntClass{c}{d}\land
      u=\IntClass{e}{f}}{\rA}
    \proofstepwidestarbreakkeep[5]{%
      \IntClass{a}{b}\leq\IntClass{c}{d}
      \leftrightarrow
      \IntClass{a}{b}+\IntClass{e}{f}
      \leq\IntClass{c}{d}+\IntClass{e}{f}}{%
      \FormulaRefAuto{IntegerOrderTranslationOnClasses}{5}}
    \proofstepwidestarbreakkeep[5]{%
      z\leq w\leftrightarrow z+u\leq w+u}{\rIE{5,6}}
    \proofstepwidestarbreakkeep[1]{%
      z\leq w\leftrightarrow z+u\leq w+u}{%
      \rEE{2,3,4,5,7}}
  \closeproofpartwide
\end{tabproofsplitwide}

\FormulaAxiomAutoR[Translationsäquivalenz der Ganzzahlordnung]{%
  z,w,u\in\mathbb Z\vdash
  \bigl(z\leq w\leftrightarrow z+u\leq w+u\bigr)%
}{IntegerOrderTranslationEquivalenceAxiom}
\par\noindent\emph{Herkunft aus der Quotientenkonstruktion:}
\FormulaRefAuto{IntegerOrderTranslationModel}[thm].\par

\FormulaAxiomAutoR[Translation erhält die Ganzzahlordnung]{%
  z,w,u\in\mathbb Z\dsep z\leq w
  \vdash z+u\leq w+u%
}{IntegerOrderTranslationAxiom}
\par\noindent\emph{Herkunft aus der Quotientenkonstruktion:}
die Hinrichtung von
\FormulaRefAuto{IntegerOrderTranslationModel}[thm].\par

\FormulaThmDeltaK[Negation einer negativen ganzen Zahl]{%
  a\in\mathbb Z\dsep a<\IntZero\vdash\IntZero<-a%
}{IntegerNegativeOrderReversal}{%
  \DeltaRow{Mengen}{a}%
  \DeltaRow{Relationsoperatoren}{<}%
}
\begin{tabproofwide}
  \proofstepwidestarbreakkeep[1]{a\in\mathbb Z}{\rA}
  \proofstepwidestarbreakkeep[2]{a<\IntZero}{\rA}
  \proofstepwidestarbreakkeep[]{\IntZero\in\mathbb Z}{%
    \FormulaRefAuto{IntegerZeroCarrierAxiom}[ax]}
  \proofstepwidestarbreakkeep[1,2]{%
    a\leq\IntZero\land a\neq\IntZero}{%
    \FormulaRefAuto{IntegerStrictOrderDef}[def]{1,3,2}}
  \proofstepwidestarbreakkeep[1]{-a\in\mathbb Z}{%
    \FormulaRefAuto{IntegerNegClosureAxiom}[ax]{1}}
  \proofstepwidestarbreakkeep[1]{a+(-a)=\IntZero}{%
    \rAEa{\FormulaRefAuto{IntegerAddInverseAxiom}[ax]{1}}}
  \proofstepwidestarbreakkeep[1]{\IntZero+(-a)=-a}{%
    \rAEb{\FormulaRefAuto{IntegerAddZeroAxiom}[ax]{5}}}
  \proofstepwidestarbreakkeep[1]{%
    a\leq\IntZero
      \leftrightarrow a+(-a)\leq\IntZero+(-a)}{%
    \FormulaRefAuto{IntegerOrderTranslationEquivalenceAxiom}[ax]{1,3,5}}
  \proofstepwidestarbreakkeep[1,2]{%
    a+(-a)\leq\IntZero+(-a)}{\rRE{8,\rAEa{4}}}
  \proofstepwidestarbreakkeep[1,2]{\IntZero\leq-a}{\rIE{6,7,9}}
  \proofstepwidestarbreakkeep[11]{-a=\IntZero}{\rA}
  \proofstepwidestarbreakkeep[11]{-(-a)=-\IntZero}{\rIE{11}}
  \proofstepwidestarbreakkeep[1]{-(-a)=a}{%
    \FormulaRefAuto{IntegerDoubleNegation}[thm]{1}}
  \proofstepwidestarbreakkeep[]{-\IntZero=\IntZero}{%
    \FormulaRefAuto{IntegerNegativeZero}[thm]}
  \proofstepwidestarbreakkeep[1,11]{a=\IntZero}{\rIE{12,13,14}}
  \proofstepwidestarbreakkeep[1,2,11]{\bot}{%
    \rBI{\rAEb{4},15}}
  \proofstepwidestarbreakkeep[1,2]{-a\neq\IntZero}{\rCI{11,16}}
  \proofstepwidestarbreakkeep[1,2]{\IntZero\neq-a}{%
    \FormulaRefAuto{a\neq b\vdash b\neq a}[thm]{17}}
  \proofstepwidestarbreakkeep[1,2]{%
    \IntZero\leq-a\land\IntZero\neq-a}{\rAI{10,18}}
  \proofstepwidestarbreakkeep[1,2]{\IntZero<-a}{%
    \FormulaRefAuto{IntegerStrictOrderDef}[def]{3,5,19}}
\end{tabproofwide}

\FormulaThmDeltaK[Null und Eins sind im Quotientenmodell verschieden]{%
  \IntZero\neq\IntOne%
}{IntegerZeroNeOneModel}{%
  \DeltaRow{Mengen}{\IntZero\dsep\IntOne}%
}
\begin{tabproof}
  \proofstep{}{0\in\mathbb N}{\FormulaRefAuto{PeanoZero}[ax]}
  \proofstep{}{1\in\mathbb N}{\FormulaRefAuto{PeanoOneInN}[thm]}
  \proofstep{}{0\neq1}{\FormulaRefAuto{PeanoZeroNeOne}[thm]}
  \proofstep{}{\IntEmbed(0)\neq\IntEmbed(1)}{%
    \FormulaRefAuto{NaturalIntegerEmbeddingInjective}[thm]{1,2,3}}
  \proofstep{}{\IntZero=\IntEmbed(0)}{%
    \FormulaRefAuto{IntegerZeroDef}[def]}
  \proofstep{}{\IntOne=\IntEmbed(1)}{%
    \FormulaRefAuto{IntegerOneDef}[def]}
  \proofstep{}{\IntZero\neq\IntOne}{\rIE{5,6,4}}
\end{tabproof}

\FormulaThmDeltaK[Eins ist im Quotientenmodell positiv]{%
  \IntZero<\IntOne%
}{IntegerOnePositiveModel}{%
  \DeltaRow{Mengen}{\IntZero\dsep\IntOne}%
}
\begin{tabproof}
  \proofstep{}{0\in\mathbb N}{\FormulaRefAuto{PeanoZero}[ax]}
  \proofstep{}{1\in\mathbb N}{\FormulaRefAuto{PeanoOneInN}[thm]}
  \proofstep{}{0\leq1}{\FormulaRefAuto{PeanoZeroLeq}[thm]{2}}
  \proofstep{}{\IntEmbed(0)=\IntClass{0}{0}}{%
    \FormulaRefAuto{NaturalIntegerEmbeddingValue}[thm]{1}}
  \proofstep{}{\IntEmbed(1)=\IntClass{1}{0}}{%
    \FormulaRefAuto{NaturalIntegerEmbeddingValue}[thm]{2}}
  \proofstep{}{\IntEmbed(0)\leq\IntEmbed(1)}{%
    \FormulaRefAuto{IntegerOrderClassChar}[thm]{1,2,3,4,5}}
  \proofstep{}{\IntZero\leq\IntOne}{%
    \rIE{6,\FormulaRefAuto{IntegerZeroDef}[def],%
      \FormulaRefAuto{IntegerOneDef}[def]}}
  \proofstep{}{\IntZero\neq\IntOne}{%
    \FormulaRefAuto{IntegerZeroNeOneModel}[thm]}
  \proofstep{}{\IntZero<\IntOne}{%
    \FormulaRefAuto{IntegerStrictOrderDef}[def]{7,8}}
\end{tabproof}

\FormulaAxiomAutoR[Null und Eins sind verschieden]{%
  \IntZero\neq\IntOne%
}{IntegerZeroNeOneAxiom}
\par\noindent\emph{Herkunft aus der Quotientenkonstruktion:}
\FormulaRefAuto{IntegerZeroNeOneModel}[thm].\par

\FormulaAxiomAutoR[Eins ist strikt positiv]{%
  \IntZero<\IntOne%
}{IntegerOnePositiveAxiom}
\par\noindent\emph{Herkunft aus der Quotientenkonstruktion:}
\FormulaRefAuto{IntegerOnePositiveModel}[thm].\par

\FormulaThmDeltaK[Positive ganze Zahlen haben positive natürliche Bilder]{%
  z\in\mathbb Z\dsep\IntZero<z\vdash
  \exists k\in\mathbb N\,(k\neq0\land z=\IntEmbed(k))%
}{IntegerPositiveEmbeddingRepresentativeModel}{%
  \DeltaRow{Mengen}{z\dsep k}%
}
\begin{tabproofsplitwide}
  \proofpartwideindR[Positiver Normalformfall]{%
    \begin{aligned}[t]
      &z\in\mathbb Z\dsep\IntZero<z\dsep k\in\mathbb N\dsep
        z=\IntEmbed(k)\vdash{}\\
      &\exists n\in\mathbb N\,(n\neq0\land z=\IntEmbed(n))
    \end{aligned}%
  }{%
    \begin{aligned}[t]
      &z\in\mathbb Z\dsep\IntZero<z\dsep k\in\mathbb N\dsep
        z=\IntEmbed(k)\vdash{}\\
      &\exists n\in\mathbb N\,(n\neq0\land z=\IntEmbed(n))
    \end{aligned}%
  }[IntegerPositiveEmbeddingRepresentativePositiveCase]
    \proofstepwidestarbreakkeep[1]{z\in\mathbb Z}{\rA}
    \proofstepwidestarbreakkeep[2]{\IntZero<z}{\rA}
    \proofstepwidestarbreakkeep[3]{k\in\mathbb N}{\rA}
    \proofstepwidestarbreakkeep[4]{z=\IntEmbed(k)}{\rA}
    \proofstepwidestarbreakkeep[5]{k=0}{\rA}
    \proofstepwidestarbreakkeep[4,5]{z=\IntZero}{%
      \rIE{4,5,\FormulaRefAuto{IntegerZeroDef}[def]}}
    \proofstepwidestarbreakkeep[1,2]{\IntZero\neq z}{%
      \rAEb{\FormulaRefAuto{IntegerStrictOrderDef}[def]{%
        \FormulaRefAuto{IntZeroCarrier},1,2}}}
    \proofstepwidestarbreakkeep[1,2,4,5]{\bot}{\rBI{7,6}}
    \proofstepwidestarbreakkeep[1,2,3,4]{k\neq0}{\rCI{5,8}}
    \proofstepwidestarbreakkeep[1,2,3,4]{%
      \exists n\in\mathbb N\,(n\neq0\land z=\IntEmbed(n))}{%
      \rEI{\rAI{3,\rAI{9,4}}}}
  \closeproofpartwideindR

  \proofpartwideindR[Negativer Normalformfall ist unmöglich]{%
    \begin{aligned}[t]
      &z\in\mathbb Z\dsep\IntZero<z\dsep k\in\mathbb N\dsep
        z=-\IntEmbed(k)\vdash{}\\
      &\exists n\in\mathbb N\,(n\neq0\land z=\IntEmbed(n))
    \end{aligned}%
  }{%
    \begin{aligned}[t]
      &z\in\mathbb Z\dsep\IntZero<z\dsep k\in\mathbb N\dsep
        z=-\IntEmbed(k)\vdash{}\\
      &\exists n\in\mathbb N\,(n\neq0\land z=\IntEmbed(n))
    \end{aligned}%
  }[IntegerPositiveEmbeddingRepresentativeNegativeCase]
    \proofstepwidestarbreakkeep[1]{z\in\mathbb Z}{\rA}
    \proofstepwidestarbreakkeep[2]{\IntZero<z}{\rA}
    \proofstepwidestarbreakkeep[3]{k\in\mathbb N}{\rA}
    \proofstepwidestarbreakkeep[4]{z=-\IntEmbed(k)}{\rA}
    \proofstepwidestarbreakkeep[1,2]{\IntZero\leq z}{%
      \rAEa{\FormulaRefAuto{IntegerStrictOrderDef}[def]{%
        \FormulaRefAuto{IntZeroCarrier},1,2}}}
    \proofstepwidestarbreakkeep[3]{-\IntEmbed(k)=\IntClass{0}{k}}{%
      \FormulaRefAuto{IntegerNegationDef}[def]{%
        3,\FormulaRefAuto{NaturalIntegerEmbeddingValue}[thm]{3}}}
    \proofstepwidestarbreakkeep[1,2,3,4]{k\leq0}{%
      \FormulaRefAuto{IntegerOrderClassChar}[thm]{%
        \FormulaRefAuto{PeanoZero}[ax],3,4,5,6,%
        \FormulaRefAuto{IntegerZeroClass}[thm]}}
    \proofstepwidestarbreakkeep[3]{0\leq k}{%
      \FormulaRefAuto{PeanoZeroLeq}[thm]{3}}
    \proofstepwidestarbreakkeep[3,7,8]{k=0}{%
      \FormulaRefAuto{PeanoLeqAntisymmetric}[thm]{%
        3,\FormulaRefAuto{PeanoZero}[ax],7,8}}
    \proofstepwidestarbreakkeep[3,4,9]{z=\IntZero}{%
      \rIE{4,9,\FormulaRefAuto{IntegerZeroDef}[def],%
        \FormulaRefAuto{IntegerNegativeZero}[thm]}}
    \proofstepwidestarbreakkeep[1,2]{\IntZero\neq z}{%
      \rAEb{\FormulaRefAuto{IntegerStrictOrderDef}[def]{%
        \FormulaRefAuto{IntZeroCarrier},1,2}}}
    \proofstepwidestarbreakkeep[1,2,3,4]{\bot}{\rBI{11,10}}
    \proofstepwidestarbreakkeep[1,2,3,4]{%
      \exists n\in\mathbb N\,(n\neq0\land z=\IntEmbed(n))}{%
      \FormulaRefAuto{\bot\vdash P}{12}}
  \closeproofpartwideindR

  \proofpartwide{Elimination der Normalform}
    \proofstepwidestarbreakkeep[1]{z\in\mathbb Z}{\rA}
    \proofstepwidestarbreakkeep[2]{\IntZero<z}{\rA}
    \proofstepwidestarbreakkeep[1]{\exists k\in\mathbb N\,
      (z=\IntEmbed(k)\lor z=-\IntEmbed(k))}{%
      \FormulaRefAuto{IntegerSignedNormalForm}[thm]{1}}
    \proofstepwidestarbreakkeep[4]{k\in\mathbb N\land
      (z=\IntEmbed(k)\lor z=-\IntEmbed(k))}{\rA}
    \proofstepwidestarbreakkeep[1,2,4]{z=\IntEmbed(k)\to
      \exists n\in\mathbb N\,(n\neq0\land z=\IntEmbed(n))}{%
      \rRI{z=\IntEmbed(k),%
        \FormulaRefAuto{IntegerPositiveEmbeddingRepresentativePositiveCase}{%
          1,2,\rAEa{4},z=\IntEmbed(k)}}}
    \proofstepwidestarbreakkeep[1,2,4]{z=-\IntEmbed(k)\to
      \exists n\in\mathbb N\,(n\neq0\land z=\IntEmbed(n))}{%
      \rRI{z=-\IntEmbed(k),%
        \FormulaRefAuto{IntegerPositiveEmbeddingRepresentativeNegativeCase}{%
          1,2,\rAEa{4},z=-\IntEmbed(k)}}}
    \proofstepwidestarbreakkeep[1,2,4]{%
      \exists n\in\mathbb N\,(n\neq0\land z=\IntEmbed(n))}{%
      \rOE{\rAEb{4},5,6}}
    \proofstepwidestarbreakkeep[1,2]{%
      \exists n\in\mathbb N\,(n\neq0\land z=\IntEmbed(n))}{%
      \rEE{3,4,7}}
  \closeproofpartwide
\end{tabproofsplitwide}

\FormulaThmDeltaK[Positive natürliche Bilder sind positive ganze Zahlen]{%
  k\in\mathbb N\dsep k\neq0\vdash\IntZero<\IntEmbed(k)%
}{NaturalIntegerEmbeddingPositiveModel}{%
  \DeltaRow{Mengen}{k\dsep h}%
}
\begin{tabproofwide}
  \proofstepwidestarbreakkeep[1]{k\in\mathbb N}{\rA}
  \proofstepwidestarbreakkeep[2]{k\neq0}{\rA}
  \proofstepwidestarbreakkeep[1,2]{\exists h\in\mathbb N\,(k=\Succ(h))}{%
    \FormulaRefAuto{PeanoNonzeroIsSucc}[ax]{1,2}}
  \proofstepwidestarbreakkeep[4]{h\in\mathbb N\land k=\Succ(h)}{\rA}
  \proofstepwidestarbreakkeep[4]{0+\Succ(h)=k}{%
    \rIE{\FormulaRefAuto{PeanoAddLeftZero}[thm]{%
      \FormulaRefAuto{PeanoSuccClosure}[ax]{\rAEa{4}}},\rAEb{4}}}
  \proofstepwidestarbreakkeep[1,4]{0<k}{%
    \FormulaRefAuto{NaturalStrictOrderAdditiveCharacterization}[thm]{%
      \FormulaRefAuto{PeanoZero}[ax],1,\rEI{\rAI{\rAEa{4},5}}}}
  \proofstepwidestarbreakkeep[1,4]{\IntEmbed(0)<\IntEmbed(k)}{%
    \FormulaRefAuto{IntegerStrictOrderDef}[def]{%
      \FormulaRefAuto{IntegerOrderClassChar}[thm]{%
        \FormulaRefAuto{PeanoZero}[ax],1,6},%
      \FormulaRefAuto{NaturalIntegerEmbeddingInjective}[thm]{%
        \FormulaRefAuto{PeanoZero}[ax],1,2}}}
  \proofstepwidestarbreakkeep[1,4]{\IntZero<\IntEmbed(k)}{%
    \rIE{7,\FormulaRefAuto{IntegerZeroDef}[def]}}
  \proofstepwidestarbreakkeep[1,2]{\IntZero<\IntEmbed(k)}{\rEE{3,4,8}}
\end{tabproofwide}

\FormulaThmDeltaK[Produkt von null verschiedener natürlicher Zahlen]{%
  m,n\in\mathbb N\dsep m\neq0\dsep n\neq0
  \vdash m\cdot n\neq0%
}{NaturalNonzeroProductForIntegerModel}{%
  \DeltaRow{Mengen}{m\dsep n\dsep i\dsep j}%
}
\begin{tabproofwide}
  \proofstepwidestarbreakkeep[1]{m,n\in\mathbb N}{\rA}
  \proofstepwidestarbreakkeep[2]{m\neq0}{\rA}
  \proofstepwidestarbreakkeep[3]{n\neq0}{\rA}
  \proofstepwidestarbreakkeep[1,2]{\exists i\in\mathbb N\,(m=\Succ(i))}{%
    \FormulaRefAuto{PeanoNonzeroIsSucc}[ax]{\rAEa{1},2}}
  \proofstepwidestarbreakkeep[1,3]{\exists j\in\mathbb N\,(n=\Succ(j))}{%
    \FormulaRefAuto{PeanoNonzeroIsSucc}[ax]{\rAEb{1},3}}
  \proofstepwidestarbreakkeep[6]{i\in\mathbb N\land m=\Succ(i)}{\rA}
  \proofstepwidestarbreakkeep[7]{j\in\mathbb N\land n=\Succ(j)}{\rA}
  \proofstepwidestarbreakkeep[6,7]{m\cdot n=m\cdot j+m}{%
    \rIE{\rAEb{7},\FormulaRefAuto{PeanoMulSuccRight}[thm]{%
      \rAEa{1},\rAEa{7}}}}
  \proofstepwidestarbreakkeep[1,6,7]{m\cdot j+m=\Succ(m\cdot j+i)}{%
    \rIE{\rAEb{6},\FormulaRefAuto{PeanoAddSuccRight}[thm]{%
      \FormulaRefAuto{PeanoMulClosure}[thm]{\rAEa{1},\rAEa{7}},%
      \rAEa{6}}}}
  \proofstepwidestarbreakkeep[1,6,7]{\Succ(m\cdot j+i)\neq0}{%
    \FormulaRefAuto{PeanoSuccNonzero}[ax]{%
      \FormulaRefAuto{PeanoAddClosure}[thm]{%
        \FormulaRefAuto{PeanoMulClosure}[thm]{\rAEa{1},\rAEa{7}},%
        \rAEa{6}}}}
  \proofstepwidestarbreakkeep[1,6,7]{m\cdot n\neq0}{\rIE{8,9,10}}
  \proofstepwidestarbreakkeep[1,2,3]{m\cdot n\neq0}{%
    \rEE{4,6,\rEE{5,7,11}}}
\end{tabproofwide}

\FormulaThmDeltaK[Produkt positiver ganzer Zahlen im Quotientenmodell]{%
  z,w\in\mathbb Z\dsep\IntZero<z\dsep\IntZero<w
  \vdash\IntZero<z\cdot w%
}{IntegerPositiveProductModel}{%
  \DeltaRow{Mengen}{z\dsep w\dsep m\dsep n}%
}
\begin{tabproofwide}
  \proofstepwidestarbreakkeep[1]{z,w\in\mathbb Z}{\rA}
  \proofstepwidestarbreakkeep[2]{\IntZero<z}{\rA}
  \proofstepwidestarbreakkeep[3]{\IntZero<w}{\rA}
  \proofstepwidestarbreakkeep[1,2]{\exists m\in\mathbb N\,
    (m\neq0\land z=\IntEmbed(m))}{%
    \FormulaRefAuto{IntegerPositiveEmbeddingRepresentativeModel}[thm]{%
      \rAEa{1},2}}
  \proofstepwidestarbreakkeep[1,3]{\exists n\in\mathbb N\,
    (n\neq0\land w=\IntEmbed(n))}{%
    \FormulaRefAuto{IntegerPositiveEmbeddingRepresentativeModel}[thm]{%
      \rAEb{1},3}}
  \proofstepwidestarbreakkeep[6]{m\in\mathbb N\land
    (m\neq0\land z=\IntEmbed(m))}{\rA}
  \proofstepwidestarbreakkeep[7]{n\in\mathbb N\land
    (n\neq0\land w=\IntEmbed(n))}{\rA}
  \proofstepwidestarbreakkeep[6,7]{m\cdot n\neq0}{%
    \FormulaRefAuto{NaturalNonzeroProductForIntegerModel}[thm]{%
      \rAEa{6},\rAEa{7},\rAEa{\rAEb{6}},\rAEa{\rAEb{7}}}}
  \proofstepwidestarbreakkeep[6,7]{\IntZero<\IntEmbed(m\cdot n)}{%
    \FormulaRefAuto{NaturalIntegerEmbeddingPositiveModel}[thm]{%
      \FormulaRefAuto{PeanoMulClosure}[thm]{\rAEa{6},\rAEa{7}},8}}
  \proofstepwidestarbreakkeep[6,7]{z\cdot w=\IntEmbed(m\cdot n)}{%
    \rIE{\rAEb{\rAEb{6}},\rAEb{\rAEb{7}},%
      \FormulaRefAuto{NaturalIntegerEmbeddingMultiplicative}[thm]{%
        \rAEa{6},\rAEa{7}}}}
  \proofstepwidestarbreakkeep[6,7]{\IntZero<z\cdot w}{\rIE{10,9}}
  \proofstepwidestarbreakkeep[1,2,3]{\IntZero<z\cdot w}{%
    \rEE{4,6,\rEE{5,7,11}}}
\end{tabproofwide}

\FormulaAxiomAutoR[Produkt strikt positiver ganzer Zahlen]{%
  z,w\in\mathbb Z\dsep\IntZero<z\dsep\IntZero<w
  \vdash\IntZero<z\cdot w%
}{IntegerPositiveProductAxiom}
\par\noindent\emph{Herkunft aus der Quotientenkonstruktion:}
\FormulaRefAuto{IntegerPositiveProductModel}[thm].\par

\FormulaThmDeltaK[Positive Produkte sind nichtnegativ]{%
  z,w\in\mathbb Z\dsep\IntZero<z\dsep\IntZero<w
  \vdash\IntZero\leq z\cdot w%
}{IntegerPositiveProductNonnegativeFromAxioms}{%
  \DeltaRow{Mengen}{z\dsep w}%
}
\begin{tabproof}
  \proofstep{1}{z,w\in\mathbb Z}{\rA}
  \proofstep{2}{\IntZero<z}{\rA}
  \proofstep{3}{\IntZero<w}{\rA}
  \proofstep{1,2,3}{\IntZero<z\cdot w}{%
    \FormulaRefAuto{IntegerPositiveProductAxiom}[ax]{1,2,3}}
  \proofstep{1,2,3}{\IntZero\leq z\cdot w}{%
    \rAEa{\FormulaRefAuto{IntegerStrictOrderDef}[def]{%
      \FormulaRefAuto{IntegerZeroCarrierAxiom}[ax],%
      \FormulaRefAuto{IntegerMulClosureAxiom}[ax]{1},4}}}
\end{tabproof}

\FormulaThmDeltaK[Positive natürliche Faktoren bewahren strikte Ordnung]{%
  a,c,k\in\mathbb N\dsep k\neq0\dsep a<c
  \vdash a\cdot k<c\cdot k%
}{NaturalPositiveFactorStrictOrderForIntegerModel}{%
  \DeltaRow{Mengen}{a\dsep c\dsep k\dsep h\dsep t}%
}
\begin{tabproofwide}
  \proofstepwidestarbreakkeep[1]{a,c,k\in\mathbb N}{\rA}
  \proofstepwidestarbreakkeep[2]{k\neq0}{\rA}
  \proofstepwidestarbreakkeep[3]{a<c}{\rA}
  \proofstepwidestarbreakkeep[1,3]{\exists h\in\mathbb N\,
    (a+\Succ(h)=c)}{%
    \FormulaRefAuto{NaturalStrictOrderAdditiveCharacterization}[thm]{%
      \rAEa{1},\rAEa{\rAEb{1}},3}}
  \proofstepwidestarbreakkeep[5]{h\in\mathbb N\land a+\Succ(h)=c}{\rA}
  \proofstepwidestarbreakkeep[1,2,5]{\Succ(h)\cdot k\neq0}{%
    \FormulaRefAuto{NaturalNonzeroProductForIntegerModel}[thm]{%
      \FormulaRefAuto{PeanoSuccClosure}[ax]{\rAEa{5}},%
      \rAEb{\rAEb{1}},%
      \FormulaRefAuto{PeanoSuccNonzero}[ax]{\rAEa{5}},2}}
  \proofstepwidestarbreakkeep[1,2,5]{\exists t\in\mathbb N\,
    (\Succ(h)\cdot k=\Succ(t))}{%
    \FormulaRefAuto{PeanoNonzeroIsSucc}[ax]{%
      \FormulaRefAuto{PeanoMulClosure}[thm]{%
        \FormulaRefAuto{PeanoSuccClosure}[ax]{\rAEa{5}},%
        \rAEb{\rAEb{1}}},6}}
  \proofstepwidestarbreakkeep[8]{t\in\mathbb N\land
    \Succ(h)\cdot k=\Succ(t)}{\rA}
  \proofstepwidestarbreakkeep[1,5]{(a+\Succ(h))\cdot k
    =a\cdot k+\Succ(h)\cdot k}{%
    \FormulaRefAuto{PeanoMulRightDistributive}[thm]{%
      \rAEa{1},\FormulaRefAuto{PeanoSuccClosure}[ax]{\rAEa{5}},%
      \rAEb{\rAEb{1}}}}
  \proofstepwidestarbreakkeep[1,5,8]{a\cdot k+\Succ(t)=c\cdot k}{%
    \rIE{\rAEb{5},\rAEb{8},9}}
  \proofstepwidestarbreakkeep[1]{a\cdot k\in\mathbb N}{%
    \FormulaRefAuto{PeanoMulClosure}[thm]{%
      \rAEa{1},\rAEb{\rAEb{1}}}}
  \proofstepwidestarbreakkeep[1]{c\cdot k\in\mathbb N}{%
    \FormulaRefAuto{PeanoMulClosure}[thm]{%
      \rAEa{\rAEb{1}},\rAEb{\rAEb{1}}}}
  \proofstepwidestarbreakkeep[1,5,8]{%
    \exists t\in\mathbb N\,(a\cdot k+\Succ(t)=c\cdot k)}{%
    \rEI{\rAI{\rAEa{8},10}}}
  \proofstepwidestarbreakkeep[1,5,8]{a\cdot k<c\cdot k}{%
    \FormulaRefAuto{NaturalStrictOrderAdditiveCharacterization}[thm]{%
      11,12,13}}
  \proofstepwidestarbreakkeep[1,2,3]{a\cdot k<c\cdot k}{%
    \rEE{4,5,\rEE{7,8,14}}}
\end{tabproofwide}

\FormulaThmDeltaK[Positive natürliche Faktoren bewahren und reflektieren Ordnung]{%
  a,c,k\in\mathbb N\dsep k\neq0
  \vdash\bigl(a\leq c\leftrightarrow a\cdot k\leq c\cdot k\bigr)%
}{NaturalPositiveFactorOrderForIntegerModel}{%
  \DeltaRow{Mengen}{a\dsep c\dsep k\dsep h}%
}
\begin{tabproofsplitwide}
  \proofpartwideindR[Erhaltung]{%
    a,c,k\in\mathbb N\dsep a\leq c
    \vdash a\cdot k\leq c\cdot k%
  }{%
    a,c,k\in\mathbb N\dsep a\leq c
    \vdash a\cdot k\leq c\cdot k%
  }[NaturalPositiveFactorOrderForwardForIntegerModel]
    \proofstepwidestarbreakkeep[1]{a,c,k\in\mathbb N}{\rA}
    \proofstepwidestarbreakkeep[2]{a\leq c}{\rA}
    \proofstepwidestarbreakkeep[1,2]{\exists h\in\mathbb N\,(a+h=c)}{%
      \FormulaRefAuto{NaturalOrderAdditiveCharacterization}[thm]{%
        \rAEa{1},\rAEa{\rAEb{1}},2}}
    \proofstepwidestarbreakkeep[4]{h\in\mathbb N\land a+h=c}{\rA}
    \proofstepwidestarbreakkeep[1,4]{a\cdot k+h\cdot k=c\cdot k}{%
      \rIE{\rAEb{4},\FormulaRefAuto{PeanoMulRightDistributive}[thm]{%
        \rAEa{1},\rAEa{4},\rAEb{\rAEb{1}}}}}
    \proofstepwidestarbreakkeep[1]{a\cdot k\in\mathbb N}{%
      \FormulaRefAuto{PeanoMulClosure}[thm]{%
        \rAEa{1},\rAEb{\rAEb{1}}}}
    \proofstepwidestarbreakkeep[1]{c\cdot k\in\mathbb N}{%
      \FormulaRefAuto{PeanoMulClosure}[thm]{%
        \rAEa{\rAEb{1}},\rAEb{\rAEb{1}}}}
    \proofstepwidestarbreakkeep[1,4]{h\cdot k\in\mathbb N}{%
      \FormulaRefAuto{PeanoMulClosure}[thm]{%
        \rAEa{4},\rAEb{\rAEb{1}}}}
    \proofstepwidestarbreakkeep[1,4]{%
      \exists t\in\mathbb N\,(a\cdot k+t=c\cdot k)}{%
      \rEI{\rAI{8,5}}}
    \proofstepwidestarbreakkeep[1,4]{a\cdot k\leq c\cdot k}{%
      \FormulaRefAuto{NaturalOrderAdditiveCharacterization}[thm]{%
        6,7,9}}
    \proofstepwidestarbreakkeep[1,2]{a\cdot k\leq c\cdot k}{\rEE{3,4,10}}
  \closeproofpartwideindR

  \proofpartwideindR[Reflexion]{%
    a,c,k\in\mathbb N\dsep k\neq0\dsep
    a\cdot k\leq c\cdot k\vdash a\leq c%
  }{%
    a,c,k\in\mathbb N\dsep k\neq0\dsep
    a\cdot k\leq c\cdot k\vdash a\leq c%
  }[NaturalPositiveFactorOrderBackwardForIntegerModel]
    \proofstepwidestarbreakkeep[1]{a,c,k\in\mathbb N}{\rA}
    \proofstepwidestarbreakkeep[2]{k\neq0}{\rA}
    \proofstepwidestarbreakkeep[3]{a\cdot k\leq c\cdot k}{\rA}
    \proofstepwidestarbreakkeep[1]{a<c\lor a=c\lor c<a}{%
      \FormulaRefAuto{PeanoOrderTrichotomy}[thm]{%
        \rAEa{1},\rAEa{\rAEb{1}}}}
    \proofstepwidestarbreakkeep[5]{a<c}{\rA}
    \proofstepwidestarbreakkeep[1,5]{a\leq c}{%
      \FormulaRefAuto{NaturalStrictOrderNotation}[def]{%
        \rAEa{1},\rAEa{\rAEb{1}},5}}
    \proofstepwidestarbreakkeep[7]{a=c}{\rA}
    \proofstepwidestarbreakkeep[1,7]{a\leq c}{%
      \rIE{7,\FormulaRefAuto{PeanoLeqReflexive}[thm]{\rAEa{1}}}}
    \proofstepwidestarbreakkeep[9]{c<a}{\rA}
    \proofstepwidestarbreakkeep[1,2,9]{c\cdot k<a\cdot k}{%
      \FormulaRefAuto{NaturalPositiveFactorStrictOrderForIntegerModel}[thm]{%
        \rAEa{\rAEb{1}},\rAEa{1},\rAEb{\rAEb{1}},2,9}}
    \proofstepwidestarbreakkeep[1,2,9]{%
      c\cdot k\leq a\cdot k\land c\cdot k\neq a\cdot k}{%
      \FormulaRefAuto{NaturalStrictOrderNotation}[def]{%
        \FormulaRefAuto{PeanoMulClosure}[thm]{%
          \rAEa{\rAEb{1}},\rAEb{\rAEb{1}}},%
        \FormulaRefAuto{PeanoMulClosure}[thm]{%
          \rAEa{1},\rAEb{\rAEb{1}}},10}}
    \proofstepwidestarbreakkeep[1,2,9]{c\cdot k\leq a\cdot k}{\rAEa{11}}
    \proofstepwidestarbreakkeep[1,2,9]{c\cdot k\neq a\cdot k}{\rAEb{11}}
    \proofstepwidestarbreakkeep[1]{a\cdot k\in\mathbb N}{%
      \FormulaRefAuto{PeanoMulClosure}[thm]{%
        \rAEa{1},\rAEb{\rAEb{1}}}}
    \proofstepwidestarbreakkeep[1]{c\cdot k\in\mathbb N}{%
      \FormulaRefAuto{PeanoMulClosure}[thm]{%
        \rAEa{\rAEb{1}},\rAEb{\rAEb{1}}}}
    \proofstepwidestarbreakkeep[1,2,3,9]{a\cdot k=c\cdot k}{%
      \FormulaRefAuto{PeanoLeqAntisymmetric}[thm]{%
        14,15,3,12}}
    \proofstepwidestarbreakkeep[1,2,3,9]{c\cdot k=a\cdot k}{%
      \FormulaRefAuto{a=b\vdash b=a}{16}}
    \proofstepwidestarbreakkeep[1,2,3,9]{\bot}{\rBI{13,17}}
    \proofstepwidestarbreakkeep[1,2,3,9]{a\leq c}{%
      \FormulaRefAuto{\bot\vdash P}{18}}
    \proofstepwidestarbreakkeep[1,2,3]{a\leq c}{%
      \rOE{4,5,6,7,8,9,19}}
  \closeproofpartwideindR

  \proofpartwide{Schluss}
    \proofstepwidestarbreakkeep[1]{a,c,k\in\mathbb N}{\rA}
    \proofstepwidestarbreakkeep[2]{k\neq0}{\rA}
    \proofstepwidestarbreakkeep[1]{a\leq c\to a\cdot k\leq c\cdot k}{%
      \rRI{a\leq c,%
        \FormulaRefAuto{NaturalPositiveFactorOrderForwardForIntegerModel}{%
          1,a\leq c}}}
    \proofstepwidestarbreakkeep[1,2]{a\cdot k\leq c\cdot k\to a\leq c}{%
      \rRI{a\cdot k\leq c\cdot k,%
        \FormulaRefAuto{NaturalPositiveFactorOrderBackwardForIntegerModel}{%
          1,2,a\cdot k\leq c\cdot k}}}
    \proofstepwidestarbreakkeep[1,2]{a\leq c\leftrightarrow
      a\cdot k\leq c\cdot k}{\rLRI{3,4}}
  \closeproofpartwide
\end{tabproofsplitwide}

\FormulaThmDeltaK[Positive Faktoren bewahren und reflektieren die Ganzzahlordnung]{%
  a,c,d\in\mathbb Z\dsep\IntZero<d\vdash
  \bigl(a\leq c\leftrightarrow a\cdot d\leq c\cdot d\bigr)%
}{IntegerPositiveFactorOrderModel}{%
  \DeltaRow{Mengen}{a\dsep c\dsep d\dsep x\dsep y\dsep u\dsep v\dsep k}%
}
\begin{tabproofsplitwide}
  \proofpartwideindR[Rechnung auf positiven Repräsentanten]{%
    \begin{aligned}[t]
      &x,y,u,v,k\in\mathbb N\dsep k\neq0\dsep
        a=\IntClass{x}{y}\dsep c=\IntClass{u}{v}\dsep{}\\
      &d=\IntEmbed(k)\vdash
        \bigl(a\leq c\leftrightarrow a\cdot d\leq c\cdot d\bigr)
    \end{aligned}%
  }{%
    \begin{aligned}[t]
      &x,y,u,v,k\in\mathbb N\dsep k\neq0\dsep
        a=\IntClass{x}{y}\dsep c=\IntClass{u}{v}\dsep
        d=\IntEmbed(k)\vdash{}\\
      &a\leq c\leftrightarrow a\cdot d\leq c\cdot d
    \end{aligned}%
  }[IntegerPositiveFactorOrderOnRepresentativesModel]
    \proofstepwidestarbreakkeep[1]{x,y,u,v,k\in\mathbb N}{\rA}
    \proofstepwidestarbreakkeep[2]{k\neq0}{\rA}
    \proofstepwidestarbreakkeep[3]{a=\IntClass{x}{y}}{\rA}
    \proofstepwidestarbreakkeep[4]{c=\IntClass{u}{v}}{\rA}
    \proofstepwidestarbreakkeep[5]{d=\IntEmbed(k)}{\rA}
    \proofstepwidestarbreakkeep[1,3,4]{a\leq c\leftrightarrow x+v\leq u+y}{%
      \FormulaRefAuto{IntegerOrderClassChar}[thm]{1,3,4}}
    \proofstepwidestarbreakkeep[1,2]{x+v\leq u+y\leftrightarrow
      (x+v)\cdot k\leq(u+y)\cdot k}{%
      \FormulaRefAuto{NaturalPositiveFactorOrderForIntegerModel}[thm]{%
        \FormulaRefAuto{PeanoAddClosure}[thm]{1},%
        \FormulaRefAuto{PeanoAddClosure}[thm]{1},%
        \rAEb{\rAEb{\rAEb{\rAEb{1}}}},2}}
    \proofstepwidestarbreakkeep[1]{(x+v)\cdot k=x\cdot k+v\cdot k}{%
      \FormulaRefAuto{PeanoMulRightDistributive}[thm]{1}}
    \proofstepwidestarbreakkeep[1]{(u+y)\cdot k=u\cdot k+y\cdot k}{%
      \FormulaRefAuto{PeanoMulRightDistributive}[thm]{1}}
    \proofstepwidestarbreakkeep[1]{\IntEmbed(k)=\IntClass{k}{0}}{%
      \FormulaRefAuto{NaturalIntegerEmbeddingValue}[thm]{%
        \rAEb{\rAEb{\rAEb{\rAEb{1}}}}}}
    \proofstepwidestarbreakkeep[1]{%
      \IntClass{x}{y}\cdot\IntClass{k}{0}
      =\IntClass{x\cdot k+y\cdot0}{x\cdot0+y\cdot k}}{%
      \FormulaRefAuto{IntegerMultiplicationDef}[def]{1}}
    \proofstepwidestarbreakkeep[1,3,5]{%
      a\cdot d=\IntClass{x\cdot k+y\cdot0}{x\cdot0+y\cdot k}}{%
      \rIE{3,5,10,11}}
    \proofstepwidestarbreakkeep[1]{y\cdot0=0}{%
      \FormulaRefAuto{PeanoMulRightZero}[thm]{1}}
    \proofstepwidestarbreakkeep[1]{x\cdot0=0}{%
      \FormulaRefAuto{PeanoMulRightZero}[thm]{1}}
    \proofstepwidestarbreakkeep[1]{x\cdot k+0=x\cdot k}{%
      \FormulaRefAuto{PeanoAddRightZero}[thm]{1}}
    \proofstepwidestarbreakkeep[1]{0+y\cdot k=y\cdot k}{%
      \FormulaRefAuto{PeanoAddLeftZero}[thm]{1}}
    \proofstepwidestarbreakkeep[1,3,5]{a\cdot d=\IntClass{x\cdot k}{y\cdot k}}{%
      \rIE{12,13,14,15,16}}
    \proofstepwidestarbreakkeep[1]{%
      \IntClass{u}{v}\cdot\IntClass{k}{0}
      =\IntClass{u\cdot k+v\cdot0}{u\cdot0+v\cdot k}}{%
      \FormulaRefAuto{IntegerMultiplicationDef}[def]{1}}
    \proofstepwidestarbreakkeep[1,4,5]{%
      c\cdot d=\IntClass{u\cdot k+v\cdot0}{u\cdot0+v\cdot k}}{%
      \rIE{4,5,10,18}}
    \proofstepwidestarbreakkeep[1]{v\cdot0=0}{%
      \FormulaRefAuto{PeanoMulRightZero}[thm]{1}}
    \proofstepwidestarbreakkeep[1]{u\cdot0=0}{%
      \FormulaRefAuto{PeanoMulRightZero}[thm]{1}}
    \proofstepwidestarbreakkeep[1]{u\cdot k+0=u\cdot k}{%
      \FormulaRefAuto{PeanoAddRightZero}[thm]{1}}
    \proofstepwidestarbreakkeep[1]{0+v\cdot k=v\cdot k}{%
      \FormulaRefAuto{PeanoAddLeftZero}[thm]{1}}
    \proofstepwidestarbreakkeep[1,4,5]{c\cdot d=\IntClass{u\cdot k}{v\cdot k}}{%
      \rIE{19,20,21,22,23}}
    \proofstepwidestarbreakkeep[1,3,4,5]{a\cdot d\leq c\cdot d
      \leftrightarrow x\cdot k+v\cdot k\leq u\cdot k+y\cdot k}{%
      \FormulaRefAuto{IntegerOrderClassChar}[thm]{1,17,24}}
    \proofstepwidestarbreakkeep[1,2,3,4,5]{a\leq c\leftrightarrow
      a\cdot d\leq c\cdot d}{\rIE{6,7,8,9,25}}
  \closeproofpartwideindR

  \proofpartwide{Elimination der Repräsentanten}
    \proofstepwidestarbreakkeep[1]{a,c,d\in\mathbb Z}{\rA}
    \proofstepwidestarbreakkeep[2]{\IntZero<d}{\rA}
    \proofstepwidestarbreakkeep[1]{\exists x,y\in\mathbb N\,
      a=\IntClass{x}{y}}{%
      \FormulaRefAuto{IntegerPairRepresentative}[thm]{\rAEa{1}}}
    \proofstepwidestarbreakkeep[1]{\exists u,v\in\mathbb N\,
      c=\IntClass{u}{v}}{%
      \FormulaRefAuto{IntegerPairRepresentative}[thm]{\rAEa{\rAEb{1}}}}
    \proofstepwidestarbreakkeep[1,2]{\exists k\in\mathbb N\,
      (k\neq0\land d=\IntEmbed(k))}{%
      \FormulaRefAuto{IntegerPositiveEmbeddingRepresentativeModel}[thm]{%
        \rAEb{\rAEb{1}},2}}
    \proofstepwidestarbreakkeep[6]{x,y\in\mathbb N\land a=\IntClass{x}{y}}{\rA}
    \proofstepwidestarbreakkeep[7]{u,v\in\mathbb N\land c=\IntClass{u}{v}}{\rA}
    \proofstepwidestarbreakkeep[8]{k\in\mathbb N\land
      (k\neq0\land d=\IntEmbed(k))}{\rA}
    \proofstepwidestarbreakkeep[6,7,8]{a\leq c\leftrightarrow
      a\cdot d\leq c\cdot d}{%
      \FormulaRefAuto{IntegerPositiveFactorOrderOnRepresentativesModel}{%
        6,7,8}}
    \proofstepwidestarbreakkeep[1,2]{a\leq c\leftrightarrow
      a\cdot d\leq c\cdot d}{%
      \rEE{3,6,\rEE{4,7,\rEE{5,8,9}}}}
  \closeproofpartwide
\end{tabproofsplitwide}

\FormulaAxiomAutoR[Positive Faktoren bewahren und reflektieren Ordnung]{%
  a,c,d\in\mathbb Z\dsep\IntZero<d\vdash
  \bigl(a\leq c\leftrightarrow a\cdot d\leq c\cdot d\bigr)%
}{IntegerPositiveFactorOrderAxiom}
\par\noindent\emph{Herkunft aus der Quotientenkonstruktion:}
\FormulaRefAuto{IntegerPositiveFactorOrderModel}[thm].\par

\FormulaThmDeltaK[Positive Faktorkürzung aus den Ordnungsaxiomen]{%
  a,c,d\in\mathbb Z\dsep\IntZero<d\dsep
  a\cdot d=c\cdot d\vdash a=c%
}{IntegerPositiveFactorCancellationFromAxioms}{%
  \DeltaRow{Mengen}{a\dsep c\dsep d}%
}
\begin{tabproof}
  \proofstep{1}{a,c,d\in\mathbb Z}{\rA}
  \proofstep{2}{\IntZero<d}{\rA}
  \proofstep{3}{a\cdot d=c\cdot d}{\rA}
  \proofstep{1}{a\cdot d\in\mathbb Z}{%
    \FormulaRefAuto{IntegerMulClosureAxiom}[ax]{%
      \rAEa{1},\rAEb{\rAEb{1}}}}
  \proofstep{1}{c\cdot d\in\mathbb Z}{%
    \FormulaRefAuto{IntegerMulClosureAxiom}[ax]{%
      \rAEa{\rAEb{1}},\rAEb{\rAEb{1}}}}
  \proofstep{1}{a\cdot d\leq a\cdot d}{%
    \FormulaRefAuto{IntegerOrderReflexiveAxiom}[ax]{4}}
  \proofstep{1}{c\cdot d\leq c\cdot d}{%
    \FormulaRefAuto{IntegerOrderReflexiveAxiom}[ax]{5}}
  \proofstep{1,3}{a\cdot d\leq c\cdot d}{\rIE{3,6}}
  \proofstep{1,3}{c\cdot d\leq a\cdot d}{\rIE{3,7}}
  \proofstep{1,2,3}{a\leq c}{%
    \rRE{\FormulaRefAuto{IntegerPositiveFactorOrderAxiom}[ax]{1,2},8}}
  \proofstep{1,2,3}{c\leq a}{%
    \rRE{\FormulaRefAuto{IntegerPositiveFactorOrderAxiom}[ax]{1,2},9}}
  \proofstep{1,2,3}{a=c}{%
    \FormulaRefAuto{IntegerOrderAntisymmetricAxiom}[ax]{%
      \rAEa{1},\rAEa{\rAEb{1}},10,11}}
\end{tabproof}

\FormulaThmDeltaK[Positive Faktoren bewahren strikte Ordnung]{%
  a,c,d\in\mathbb Z\dsep\IntZero<d\dsep a<c
  \vdash a\cdot d<c\cdot d%
}{IntegerPositiveFactorStrictOrderFromAxioms}{%
  \DeltaRow{Mengen}{a\dsep c\dsep d}%
}
\begin{tabproofwide}
  \proofstepwidestarbreakkeep[1]{a,c,d\in\mathbb Z}{\rA}
  \proofstepwidestarbreakkeep[2]{\IntZero<d}{\rA}
  \proofstepwidestarbreakkeep[3]{a<c}{\rA}
  \proofstepwidestarbreakkeep[1,3]{a\leq c}{%
    \rAEa{\FormulaRefAuto{IntegerStrictOrderDef}[def]{%
      \rAEa{1},\rAEa{\rAEb{1}},3}}}
  \proofstepwidestarbreakkeep[1,2,3]{a\cdot d\leq c\cdot d}{%
    \rRE{\FormulaRefAuto{IntegerPositiveFactorOrderAxiom}[ax]{1,2},4}}
  \proofstepwidestarbreakkeep[6]{a\cdot d=c\cdot d}{\rA}
  \proofstepwidestarbreakkeep[1,2,6]{a=c}{%
    \FormulaRefAuto{IntegerPositiveFactorCancellationFromAxioms}[thm]{%
      1,2,6}}
  \proofstepwidestarbreakkeep[1,3,6]{\bot}{%
    \rBI{\rAEb{\FormulaRefAuto{IntegerStrictOrderDef}[def]{%
      \rAEa{1},\rAEa{\rAEb{1}},3}},7}}
  \proofstepwidestarbreakkeep[1,2,3]{a\cdot d\neq c\cdot d}{\rCI{6,8}}
  \proofstepwidestarbreakkeep[1]{a\cdot d\in\mathbb Z}{%
    \FormulaRefAuto{IntegerMulClosureAxiom}[ax]{%
      \rAEa{1},\rAEb{\rAEb{1}}}}
  \proofstepwidestarbreakkeep[1]{c\cdot d\in\mathbb Z}{%
    \FormulaRefAuto{IntegerMulClosureAxiom}[ax]{%
      \rAEa{\rAEb{1}},\rAEb{\rAEb{1}}}}
  \proofstepwidestarbreakkeep[1,2,3]{%
    a\cdot d\leq c\cdot d\land a\cdot d\neq c\cdot d}{\rAI{5,9}}
  \proofstepwidestarbreakkeep[1,2,3]{a\cdot d<c\cdot d}{%
    \FormulaRefAuto{IntegerStrictOrderDef}[def]{10,11,12}}
\end{tabproofwide}

\FormulaDefDeltaK[Hergeleitete geordnete Ganzzahlaxiome]{%
  \IntegerArithmetic{%
    \mathbb Z,\IntZero,\IntOne,-,+,\cdot,\leq_{\mathbb Z}}%
}{IntegerOrderedArithmeticAxioms}{%
  \DeltaRow{\textbf{Ring mit Eins}}{}%
  \DeltaRow{Arithmetik}{%
    \IntegerArithmetic{\mathbb Z,\IntZero,\IntOne,-,+,\cdot}}
    [\FormulaRefAuto{IntegerArithmeticAxioms}]%
  \end{DeltaContext}
  \begin{DeltaContext}{ }%
  \DeltaRow{\textbf{Ordnungsstruktur}}{}%
  \DeltaRow{Totale Ordnung}{\TotOrd{\mathbb Z,\leq_{\mathbb Z}}}
    [\FormulaRefAuto{IntegerOrderTotalAxiom}]%
  \end{DeltaContext}
  \begin{DeltaContext}{ }%
  \DeltaRow{\textbf{Entfaltete Regeln}}{}%
  \DeltaRow{Reflexivität}{z\in\mathbb Z\vdash z\leq z}
    [\FormulaRefAuto{IntegerOrderReflexiveAxiom}]%
  \DeltaRow{Antisymmetrie}{%
    \begin{aligned}[t]
      &z,w\in\mathbb Z\dsep z\leq w\dsep{}\\[-2pt]
      &w\leq z\vdash z=w
    \end{aligned}}
    [\FormulaRefAuto{IntegerOrderAntisymmetricAxiom}]%
  \DeltaRow{Transitivität}{%
    \begin{aligned}[t]
      &z,w,u\in\mathbb Z\dsep z\leq w\dsep{}\\[-2pt]
      &w\leq u\vdash z\leq u
    \end{aligned}}
    [\FormulaRefAuto{IntegerOrderTransitiveAxiom}]%
  \DeltaRow{Vergleichbarkeit}{%
    z,w\in\mathbb Z\vdash z\leq w\lor w\leq z}
    [\FormulaRefAuto{IntegerOrderComparableAxiom}]%
  \end{DeltaContext}
  \begin{DeltaContext}{ }%
  \DeltaRow{\textbf{Verträglichkeit}}{}%
  \DeltaRow{Translation}{%
    \begin{aligned}[t]
      &z,w,u\in\mathbb Z\vdash{}\\[-2pt]
      &z\leq w\leftrightarrow z+u\leq w+u
    \end{aligned}}
    [\FormulaRefAuto{IntegerOrderTranslationEquivalenceAxiom}]%
  \DeltaRow{Nichttrivialität}{\IntZero\neq\IntOne}
    [\FormulaRefAuto{IntegerZeroNeOneAxiom}]%
  \DeltaRow{Orientierung}{\IntZero<\IntOne}
    [\FormulaRefAuto{IntegerOnePositiveAxiom}]%
  \DeltaRow{Positives Produkt}{%
    \begin{aligned}[t]
      &z,w\in\mathbb Z\dsep\IntZero<z\dsep{}\\[-2pt]
      &\IntZero<w\vdash\IntZero<z\cdot w
    \end{aligned}}
    [\FormulaRefAuto{IntegerPositiveProductAxiom}]%
  \end{DeltaContext}
  \begin{DeltaContext}{ }%
  \DeltaRow{\textbf{Abgeleitet}}{}%
  \DeltaRow{Faktorordnung}{%
    \begin{aligned}[t]
      &a,c,d\in\mathbb Z\dsep\IntZero<d\vdash{}\\[-2pt]
      &a\leq c\leftrightarrow a\cdot d\leq c\cdot d
    \end{aligned}}
    [\FormulaRefAuto{IntegerPositiveFactorOrderAxiom}]%
  \end{DeltaContext}
  \begin{DeltaContext}{ }%
  \DeltaRow{\textbf{Signatur}}{}%
  \DeltaRow{Ordnung}{\leq_{\mathbb Z}}%
}

Die als aus \(\mathsf{TotOrd}\) entfaltet beziehungsweise als abgeleitete
Schnittstellenregel aufgeführten Zeilen sind keine zusätzlichen
unabhängigen Axiome: Die vier Ordnungsregeln sind bereits in
\(\TotOrd{\mathbb Z,\leq_{\mathbb Z}}\) enthalten, und die positive
Faktorordnung wurde oben hergeleitet. Sie bleiben hier mit ihren eigenen
Registrierschlüsseln sichtbar, damit spätere Beweise kurze und stabile
Verweise verwenden können.

Auch diese Zusammenstellung schließt die Konstruktion noch nicht ab. Die
folgenden Modellargumente liefern insbesondere die diskrete
Nachfolgerorientierung und das zweiseitige Induktionsschema; erst danach wird
die vollständige Axiomatik zur alleinigen Grundlage gemacht.

\subsection{Abgeleitete gemischte Transitivitätsregeln}

\FormulaThmDeltaK[Gemischte Transitivität, links nicht-strikt]{%
  x,y,z\in\mathbb Z\dsep x\leq y\dsep y<z\vdash x<z%
}{IntegerLeqStrictTransitiveFromAxioms}{%
  \DeltaRow{Mengen}{x\dsep y\dsep z}%
}
\begin{tabproofwide}
  \proofstepwidestar[1]{x,y,z\in\mathbb Z}{\rA}
  \proofstepwidestar[2]{x\leq y}{\rA}
  \proofstepwidestar[3]{y<z}{\rA}
  \proofstepwidestar[1,3]{y\leq z\land y\neq z}{%
    \FormulaRefAuto{IntegerStrictOrderDef}[def]{1,3}}
  \proofstepwidestar[1,2,3]{x\leq z}{%
    \FormulaRefAuto{IntegerOrderTransitiveAxiom}[ax]{1,2,\rAEa{4}}}
  \proofstepwidestar[6]{x=z}{\rA}
  \proofstepwidestar[2,6]{z\leq y}{\rIE{6,2}}
  \proofstepwidestar[1,2,3,6]{y=z}{%
    \FormulaRefAuto{IntegerOrderAntisymmetricAxiom}[ax]{%
      1,\rAEa{4},7}}
  \proofstepwidestar[1,2,3,6]{\bot}{\rBI{\rAEb{4},8}}
  \proofstepwidestar[1,2,3]{x\neq z}{\rCI{6,9}}
  \proofstepwidestar[1,2,3]{x\leq z\land x\neq z}{\rAI{5,10}}
  \proofstepwidestar[1,2,3]{x<z}{%
    \FormulaRefAuto{IntegerStrictOrderDef}[def]{1,11}}
\end{tabproofwide}

\FormulaThmDeltaK[Gemischte Transitivität, rechts nicht-strikt]{%
  x,y,z\in\mathbb Z\dsep x<y\dsep y\leq z\vdash x<z%
}{IntegerStrictLeqTransitiveFromAxioms}{%
  \DeltaRow{Mengen}{x\dsep y\dsep z}%
}
\begin{tabproofwide}
  \proofstepwidestar[1]{x,y,z\in\mathbb Z}{\rA}
  \proofstepwidestar[2]{x<y}{\rA}
  \proofstepwidestar[3]{y\leq z}{\rA}
  \proofstepwidestar[1,2]{x\leq y\land x\neq y}{%
    \FormulaRefAuto{IntegerStrictOrderDef}[def]{1,2}}
  \proofstepwidestar[1,2,3]{x\leq z}{%
    \FormulaRefAuto{IntegerOrderTransitiveAxiom}[ax]{1,\rAEa{4},3}}
  \proofstepwidestar[6]{x=z}{\rA}
  \proofstepwidestar[2,6]{z\leq y}{\rIE{6,\rAEa{4}}}
  \proofstepwidestar[1,2,3,6]{y=z}{%
    \FormulaRefAuto{IntegerOrderAntisymmetricAxiom}[ax]{1,3,7}}
  \proofstepwidestar[1,2,3,6]{x=y}{\rIE{6,8}}
  \proofstepwidestar[1,2,3,6]{\bot}{\rBI{\rAEb{4},9}}
  \proofstepwidestar[1,2,3]{x\neq z}{\rCI{6,10}}
  \proofstepwidestar[1,2,3]{x\leq z\land x\neq z}{\rAI{5,11}}
  \proofstepwidestar[1,2,3]{x<z}{%
    \FormulaRefAuto{IntegerStrictOrderDef}[def]{1,12}}
\end{tabproofwide}

\section{Positive ganze Zahlen und Kürzungslemmata}

\FormulaDefDeltaK[Positive ganze Zahlen]{%
  \PositiveIntegers\coloneqq
  \{b\in\mathbb Z\mid \IntZero<b\}%
}{PositiveIntegersDef}{%
  \DeltaRow{Mengen}{\PositiveIntegers}%
  \DeltaRow{Relationsoperatoren}{<}%
  \DeltaRow{Neue Symbole}{\PositiveIntegers}%
}

\FormulaThmDeltaK[Die Ganzzahleins ist positiv]{%
  \IntOne\in\PositiveIntegers%
}{IntegerOnePositive}{%
  \DeltaRow{Mengen}{\PositiveIntegers}%
}
\begin{tabproof}
  \proofstep{}{\IntZero<\IntOne}{%
    \FormulaRefAuto{IntegerOnePositiveAxiom}[ax]}
  \proofstep{}{\IntOne\in\mathbb Z}{%
    \FormulaRefAuto{IntegerOneCarrierAxiom}[ax]}
  \proofstep{}{\IntOne\in\PositiveIntegers}{%
    \FormulaRefAuto{PositiveIntegersDef}[def]{1,2}}
\end{tabproof}

\FormulaThmDeltaK[Positive ganze Zahlen sind von null verschieden]{%
  b\in\PositiveIntegers\vdash b\neq\IntZero%
}{PositiveIntegerNonzero}{%
  \DeltaRow{Mengen}{b}%
}
\begin{tabproof}
  \proofstep{1}{b\in\PositiveIntegers}{\rA}
  \proofstep{1}{b\in\mathbb Z\land\IntZero<b}{%
    \FormulaRefAuto{PositiveIntegersDef}[def]{1}}
  \proofstep{1}{\IntZero\leq b\land\IntZero\neq b}{%
    \FormulaRefAuto{IntegerStrictOrderDef}[def]{2}}
  \proofstep{1}{b\neq\IntZero}{%
    \FormulaRefAuto{a\neq b\vdash b\neq a}{\rAEb{3}}}
\end{tabproof}

\FormulaThmDeltaK[Die Ganzzahleins ist kleinste positive ganze Zahl]{%
  b\in\PositiveIntegers\vdash\IntOne\leq b%
}{PositiveIntegerOneLowerBound}{%
  \DeltaRow{Mengen}{b\dsep k\dsep h}%
  \DeltaRow{Relationsoperatoren}{\leq}%
}
\begin{tabproofwide}
  \proofstepwidestar[1]{b\in\PositiveIntegers}{\rA}
  \proofstepwidestar[1]{b\in\mathbb Z\land\IntZero<b}{%
    \FormulaRefAuto{PositiveIntegersDef}[def]{1}}
  \proofstepwidestar[1]{%
    \exists k\in\mathbb N\,(k\neq0\land b=\IntEmbed(k))}{%
    \FormulaRefAuto{IntegerPositiveEmbeddingRepresentativeModel}[thm]{%
      \rAEa{2},\rAEb{2}}}
  \proofstepwidestar[4]{%
    k\in\mathbb N\land k\neq0\land b=\IntEmbed(k)}{\rA}
  \proofstepwidestar[4]{\exists h\in\mathbb N\,(k=\Succ(h))}{%
    \FormulaRefAuto{PeanoNonzeroIsSucc}[ax]{\rAEa{4},\rAEa{\rAEb{4}}}}
  \proofstepwidestar[6]{h\in\mathbb N\land k=\Succ(h)}{\rA}
  \proofstepwidestar[6]{h+1=\Succ(h)}{%
    \FormulaRefAuto{PeanoAddOneSucc}[thm]{\rAEa{6}}}
  \proofstepwidestar[6]{1+h=h+1}{%
    \FormulaRefAuto{PeanoAddCommutative}[thm]{%
      \FormulaRefAuto{PeanoOneInN}[thm],\rAEa{6}}}
  \proofstepwidestar[6]{1+h=\Succ(h)}{\rChain{8,7}}
  \proofstepwidestar[4,6]{1+h=k}{\rIE{9,\rAEb{6}}}
  \proofstepwidestar[4,6]{h\in\mathbb N\land 1+h=k}{%
    \rAI{\rAEa{6},10}}
  \proofstepwidestar[4,6]{\exists h\in\mathbb N\,(1+h=k)}{\rEI{11}}
  \proofstepwidestar[4,6]{1\leq k}{%
    \FormulaRefAuto{NaturalOrderAdditiveCharacterization}[thm]{%
      \FormulaRefAuto{PeanoOneInN}[thm],\rAEa{4},%
      12}}
  \proofstepwidestar[4,6]{\IntEmbed(1)\leq\IntEmbed(k)}{%
    \FormulaRefAuto{IntegerOrderClassChar}[thm]{%
      \FormulaRefAuto{PeanoOneInN}[thm],\rAEa{4},13}}
  \proofstepwidestar[]{\IntOne=\IntEmbed(1)}{%
    \FormulaRefAuto{IntegerOneDef}[def]}
  \proofstepwidestar[4]{b=\IntEmbed(k)}{\rAEb{\rAEb{4}}}
  \proofstepwidestar[4,6]{\IntOne\leq b}{%
    \rIE{15,14,16}}
  \proofstepwidestar[4]{\IntOne\leq b}{\rEE{5,6,17}}
  \proofstepwidestar[1]{\IntOne\leq b}{\rEE{3,4,18}}
\end{tabproofwide}

\FormulaThmDeltaK[Abgeschlossenheit positiver ganzer Zahlen unter Multiplikation]{%
  b\in\PositiveIntegers\dsep d\in\PositiveIntegers
  \vdash b\cdot d\in\PositiveIntegers%
}{PositiveIntegerMulClosure}{%
  \DeltaRow{Mengen}{b\dsep d}%
}
\begin{tabproof}
  \proofstep{1}{b\in\PositiveIntegers}{\rA}
  \proofstep{2}{d\in\PositiveIntegers}{\rA}
  \proofstep{1}{b\in\mathbb Z\land\IntZero<b}{%
    \FormulaRefAuto{PositiveIntegersDef}[def]{1}}
  \proofstep{2}{d\in\mathbb Z\land\IntZero<d}{%
    \FormulaRefAuto{PositiveIntegersDef}[def]{2}}
  \proofstep{1,2}{b\cdot d\in\mathbb Z}{%
    \FormulaRefAuto{IntegerMulClosureAxiom}[ax]{\rAEa{3},\rAEa{4}}}
  \proofstep{1,2}{\IntZero<b\cdot d}{%
    \FormulaRefAuto{IntegerPositiveProductAxiom}[ax]{3,4}}
  \proofstep{1,2}{b\cdot d\in\PositiveIntegers}{%
    \FormulaRefAuto{PositiveIntegersDef}[def]{5,6}}
\end{tabproof}

\FormulaThmDeltaK[Positive ganzzahlige Faktoren bewahren und reflektieren die Ordnung]{%
  a,c\in\mathbb Z\dsep d\in\PositiveIntegers
  \vdash
  \bigl(a\leq c\leftrightarrow a\cdot d\leq c\cdot d\bigr)%
}{IntegerPositiveFactorOrder}{%
  \DeltaRow{Mengen}{a\dsep c\dsep d}%
}
\begin{tabproof}
  \proofstep{1}{a,c\in\mathbb Z}{\rA}
  \proofstep{2}{d\in\PositiveIntegers}{\rA}
  \proofstep{2}{d\in\mathbb Z\land\IntZero<d}{%
    \FormulaRefAuto{PositiveIntegersDef}[def]{2}}
  \proofstep{1,2}{a\leq c\leftrightarrow
    a\cdot d\leq c\cdot d}{%
    \FormulaRefAuto{IntegerPositiveFactorOrderAxiom}[ax]{1,3}}
\end{tabproof}

\FormulaThmDeltaK[Kürzung eines positiven ganzzahligen Faktors]{%
  a,c\in\mathbb Z\dsep d\in\PositiveIntegers\dsep
  a\cdot d=c\cdot d\vdash a=c%
}{IntegerPositiveFactorCancellation}{%
  \DeltaRow{Mengen}{a\dsep c\dsep d}%
}
\begin{tabproof}
  \proofstep{1}{a,c\in\mathbb Z}{\rA}
  \proofstep{2}{d\in\PositiveIntegers}{\rA}
  \proofstep{3}{a\cdot d=c\cdot d}{\rA}
  \proofstep{2}{d\in\mathbb Z\land\IntZero<d}{%
    \FormulaRefAuto{PositiveIntegersDef}[def]{2}}
  \proofstep{1,2,3}{a=c}{%
    \FormulaRefAuto{IntegerPositiveFactorCancellationFromAxioms}[thm]{%
      1,4,3}}
\end{tabproof}

\FormulaThmDeltaK[Positive natürliche Bilder wachsen unter positiven Faktoren]{%
  n\in\mathbb N\dsep n\neq0\dsep b\in\PositiveIntegers
  \vdash \IntEmbed(n)\leq\IntEmbed(n)\cdot b%
}{PositiveIntegerEmbeddingFactorLowerBound}{%
  \DeltaRow{Mengen}{n\dsep b}%
  \DeltaRow{Funktion}{\IntEmbed}%
}
\begin{tabproofwide}
  \proofstepwidestar[1]{n\in\mathbb N}{\rA}
  \proofstepwidestar[2]{n\neq0}{\rA}
  \proofstepwidestar[3]{b\in\PositiveIntegers}{\rA}
  \proofstepwidestar[1]{\IntEmbed(n)\in\mathbb Z}{%
    \FormulaRefAuto{NaturalIntegerEmbeddingFunction}[thm]{1}}
  \proofstepwidestar[1,2]{\IntZero<\IntEmbed(n)}{%
    \FormulaRefAuto{NaturalIntegerEmbeddingPositiveModel}[thm]{1,2}}
  \proofstepwidestar[3]{\IntOne\leq b}{%
    \FormulaRefAuto{PositiveIntegerOneLowerBound}[thm]{3}}
  \proofstepwidestar[3]{b\in\mathbb Z\land\IntZero<b}{%
    \FormulaRefAuto{PositiveIntegersDef}[def]{3}}
  \proofstepwidestar[]{\IntOne\in\mathbb Z}{%
    \FormulaRefAuto{IntegerOneCarrierAxiom}[ax]}
  \proofstepwidestar[1,2,3]{%
    \IntOne\cdot\IntEmbed(n)\leq b\cdot\IntEmbed(n)}{%
    \rRE{\FormulaRefAuto{IntegerPositiveFactorOrderAxiom}[ax]{%
      8,\rAEa{7},4,5},6}}
  \proofstepwidestar[1]{\IntOne\cdot\IntEmbed(n)=\IntEmbed(n)}{%
    \rAEb{\FormulaRefAuto{IntegerMulOneAxiom}[ax]{4}}}
  \proofstepwidestar[1,3]{b\cdot\IntEmbed(n)
    =\IntEmbed(n)\cdot b}{%
    \FormulaRefAuto{IntegerMulCommutativeAxiom}[ax]{\rAEa{7},4}}
  \proofstepwidestar[1,2,3]{\IntEmbed(n)\leq\IntEmbed(n)\cdot b}{%
    \rIE{10,9,11}}
\end{tabproofwide}


\section{Nachfolger und Diskretheit der Ganzzahlordnung}

\FormulaThmDeltaK[Jede ganze Zahl ist kleiner als ihr Nachfolger]{%
  z\in\mathbb Z\vdash z<\IntSucc(z)%
}{IntegerSuccessorPositive}{%
  \DeltaRow{Mengen}{z\dsep a\dsep b}%
}
\begin{tabproofsplitwide}
  \proofpartwideindR[Ordnungsteil]{%
    z\in\mathbb Z\vdash z\leq\IntSucc(z)%
  }{%
    z\in\mathbb Z\vdash z\leq\IntSucc(z)%
  }[IntegerSuccessorLeq]
    \proofstepwidestarbreakkeep[1]{z\in\mathbb Z}{\rA}
    \proofstepwidestarbreakkeep[1]{\exists a,b\in\mathbb N\,
      z=\IntClass{a}{b}}{\FormulaRefAuto{IntegerPairRepresentative}{1}}
    \proofstepwidestarbreakkeep[3]{a,b\in\mathbb N\land z=\IntClass{a}{b}}{\rA}
    \proofstepwidestarbreakkeep[3]{\IntSucc(z)=\IntClass{\Succ(a)}{b}}{%
      \FormulaRefAuto{IntegerSuccessorClassValue}{3}}
    \proofstepwidestarbreakkeep[3]{a\leq\Succ(a)}{%
      \FormulaRefAuto{PeanoLeqAddOne}{3}}
    \proofstepwidestarbreakkeep[3]{a+b\leq\Succ(a)+b}{%
      \FormulaRefAuto{NaturalOrderTranslationForward}{3,5}}
    \proofstepwidestarbreakkeep[3]{z\leq\IntSucc(z)}{%
      \FormulaRefAuto{IntegerOrderClassChar}{3,4,6}}
    \proofstepwidestarbreakkeep[1]{z\leq\IntSucc(z)}{\rEE{2,3,7}}
  \closeproofpartwideindR

  \proofpartwideindR[Ungleichheitsteil]{%
    z\in\mathbb Z\vdash z\neq\IntSucc(z)%
  }{%
    z\in\mathbb Z\vdash z\neq\IntSucc(z)%
  }[IntegerSuccessorDistinct]
    \proofstepwidestarbreakkeep[1]{z\in\mathbb Z}{\rA}
    \proofstepwidestarbreakkeep[2]{z=\IntSucc(z)}{\rA}
    \proofstepwidestarbreakkeep[1]{\exists a,b\in\mathbb N\,
      z=\IntClass{a}{b}}{\FormulaRefAuto{IntegerPairRepresentative}{1}}
    \proofstepwidestarbreakkeep[4]{a,b\in\mathbb N\land z=\IntClass{a}{b}}{\rA}
    \proofstepwidestarbreakkeep[4]{\IntSucc(z)=\IntClass{\Succ(a)}{b}}{%
      \FormulaRefAuto{IntegerSuccessorClassValue}{4}}
    \proofstepwidestarbreakkeep[2,4]{\IntClass{a}{b}=\IntClass{\Succ(a)}{b}}{\rIE{4,5,2}}
    \proofstepwidestarbreakkeep[2,4]{a+b=b+\Succ(a)}{%
      \FormulaRefAuto{IntegerClassEquality}{4,6}}
    \proofstepwidestarbreakkeep[4]{b+\Succ(a)=\Succ(a)+b}{%
      \FormulaRefAuto{PeanoAddCommutative}{4}}
    \proofstepwidestarbreakkeep[2,4]{a+b=\Succ(a)+b}{\rIE{7,8}}
    \proofstepwidestarbreakkeep[2,4]{a=\Succ(a)}{%
      \FormulaRefAuto{PeanoAddRightCancellation}{4,9}}
    \proofstepwidestarbreakkeep[4]{a<\Succ(a)}{%
      \FormulaRefAuto{n\in\mathbb{N}\vdash n<\Succ(n)}{4}}
    \proofstepwidestarbreakkeep[2,4]{\bot}{%
      \FormulaRefAuto{PeanoStrictEqualityImpossible}{4,11,10}}
    \proofstepwidestarbreakkeep[1,2]{\bot}{\rEE{3,4,12}}
    \proofstepwidestarbreakkeep[1]{z\neq\IntSucc(z)}{\rCI{2,13}}
  \closeproofpartwideindR

  \proofpartwide{Schluss}
    \proofstepwidestarbreakkeep[1]{z\in\mathbb Z}{\rA}
    \proofstepwidestarbreakkeep[1]{z\leq\IntSucc(z)}{%
      \FormulaRefAuto{IntegerSuccessorLeq}{1}}
    \proofstepwidestarbreakkeep[1]{z\neq\IntSucc(z)}{%
      \FormulaRefAuto{IntegerSuccessorDistinct}{1}}
    \proofstepwidestarbreakkeep[1]{z<\IntSucc(z)}{%
      \FormulaRefAuto{IntegerStrictOrderDef}{1,\rAI{2,3}}}
  \closeproofpartwide
\end{tabproofsplitwide}

\FormulaAxiomAutoR[Diskrete Nachfolgerorientierung]{%
  z\in\mathbb Z\vdash z<\IntSucc(z)%
}{IntegerSuccessorPositiveAxiom}
\par\noindent\emph{Herkunft aus der Quotientenkonstruktion:}
\FormulaRefAuto{IntegerSuccessorPositive}[thm].\par

\subsection{Ganzzahlschranken mit positiven Faktoren}

\FormulaThmDeltaK[Die Nachfolgerhöhe einer Normalform ist eine strikte Schranke]{%
  a\in\mathbb Z\dsep m\in\mathbb N\dsep
  \bigl(a=\IntEmbed(m)\lor a=-\IntEmbed(m)\bigr)
  \vdash a<\IntEmbed(\Succ(m))%
}{IntegerNormalFormSuccessorBound}{%
  \DeltaRow{Mengen}{a\dsep m}%
  \DeltaRow{Funktion}{\IntEmbed}%
}
\begin{tabproofwide}
  \proofstepwidestar[1]{a\in\mathbb Z}{\rA}
  \proofstepwidestar[2]{m\in\mathbb N}{\rA}
  \proofstepwidestar[3]{a=\IntEmbed(m)\lor a=-\IntEmbed(m)}{\rA}
  \proofstepwidestar[2]{\IntEmbed(m)\in\mathbb Z}{%
    \FormulaRefAuto{NaturalIntegerEmbeddingFunction}[thm]{2}}
  \proofstepwidestar[2]{\IntEmbed(m)<\IntSucc(\IntEmbed(m))}{%
    \FormulaRefAuto{IntegerSuccessorPositiveAxiom}[ax]{4}}
  \proofstepwidestar[2]{\IntSucc(\IntEmbed(m))
    =\IntEmbed(\Succ(m))}{%
    \FormulaRefAuto{IntegerSuccessorOnEmbedding}[thm]{2}}
  \proofstepwidestar[2]{\IntEmbed(m)<\IntEmbed(\Succ(m))}{\rIE{5,6}}
  \proofstepwidestar[8]{a=\IntEmbed(m)}{\rA}
  \proofstepwidestar[2,8]{a<\IntEmbed(\Succ(m))}{\rIE{8,7}}
  \proofstepwidestar[10]{a=-\IntEmbed(m)}{\rA}
  \proofstepwidestar[2]{0\leq m}{%
    \FormulaRefAuto{PeanoZeroLeq}[thm]{2}}
  \proofstepwidestar[2]{\IntEmbed(0)\leq\IntEmbed(m)}{%
    \FormulaRefAuto{IntegerOrderClassChar}[thm]{%
      \FormulaRefAuto{PeanoZero}[ax],2,11}}
  \proofstepwidestar[]{\IntZero=\IntEmbed(0)}{%
    \FormulaRefAuto{IntegerZeroDef}[def]}
  \proofstepwidestar[2]{\IntZero\leq\IntEmbed(m)}{\rIE{13,12}}
  \proofstepwidestar[2]{-\IntEmbed(m)\in\mathbb Z}{%
    \FormulaRefAuto{IntegerNegClosureAxiom}[ax]{4}}
  \proofstepwidestar[2]{%
    \IntZero+(-\IntEmbed(m))
    \leq\IntEmbed(m)+(-\IntEmbed(m))}{%
    \FormulaRefAuto{IntegerOrderTranslationAxiom}[ax]{%
      \FormulaRefAuto{IntegerZeroCarrierAxiom}[ax],4,15,14}}
  \proofstepwidestar[2]{\IntZero+(-\IntEmbed(m))=-\IntEmbed(m)}{%
    \rAEb{\FormulaRefAuto{IntegerAddZeroAxiom}[ax]{15}}}
  \proofstepwidestar[2]{\IntEmbed(m)+(-\IntEmbed(m))=\IntZero}{%
    \rAEa{\FormulaRefAuto{IntegerAddInverseAxiom}[ax]{4}}}
  \proofstepwidestar[2]{-\IntEmbed(m)\leq\IntZero}{\rIE{16,17,18}}
  \proofstepwidestar[2]{\Succ(m)\in\mathbb N}{%
    \FormulaRefAuto{PeanoSuccClosure}[ax]{2}}
  \proofstepwidestar[2]{\IntEmbed(\Succ(m))\in\mathbb Z}{%
    \FormulaRefAuto{NaturalIntegerEmbeddingFunction}[thm]{20}}
  \proofstepwidestar[2]{\Succ(m)\neq0}{%
    \FormulaRefAuto{PeanoSuccNonzero}[ax]{2}}
  \proofstepwidestar[2]{\IntZero<\IntEmbed(\Succ(m))}{%
    \FormulaRefAuto{NaturalIntegerEmbeddingPositiveModel}[thm]{20,22}}
  \proofstepwidestar[2]{-\IntEmbed(m)<\IntEmbed(\Succ(m))}{%
    \FormulaRefAuto{IntegerLeqStrictTransitiveFromAxioms}[thm]{%
      15,\FormulaRefAuto{IntegerZeroCarrierAxiom}[ax],21,19,23}}
  \proofstepwidestar[2,10]{a<\IntEmbed(\Succ(m))}{\rIE{10,24}}
  \proofstepwidestar[1,2,3]{a<\IntEmbed(\Succ(m))}{%
    \rOE{3,8,9,10,25}}
\end{tabproofwide}

\FormulaThmDeltaK[Ganzzahlschranke für positive ganzzahlige Faktoren]{%
  a\in\mathbb Z\dsep b\in\PositiveIntegers\vdash
  \exists n\in\mathbb N\,
  a<\IntEmbed(n)\cdot b%
}{RationalIntegerFractionBoundModel}{%
  \DeltaRow{Mengen}{a\dsep b\dsep n}%
  \DeltaRow{Funktion}{\IntEmbed}%
}
\begin{tabproofwide}
  \proofstepwidestar[1]{a\in\mathbb Z}{\rA}
  \proofstepwidestar[2]{b\in\PositiveIntegers}{\rA}
  \proofstepwidestar[1]{\exists m\in\mathbb N\,
    (a=\IntEmbed(m)\lor a=-\IntEmbed(m))}{%
    \FormulaRefAuto{IntegerSignedNormalForm}[thm]{1}}
  \proofstepwidestar[4]{m\in\mathbb N\land
    (a=\IntEmbed(m)\lor a=-\IntEmbed(m))}{\rA}
  \proofstepwidestar[4]{\Succ(m)\in\mathbb N}{%
    \FormulaRefAuto{PeanoSuccClosure}[ax]{\rAEa{4}}}
  \proofstepwidestar[1,4]{a<\IntEmbed(\Succ(m))}{%
    \FormulaRefAuto{IntegerNormalFormSuccessorBound}[thm]{1,4}}
  \proofstepwidestar[4]{\Succ(m)\neq0}{%
    \FormulaRefAuto{PeanoSuccNonzero}[ax]{\rAEa{4}}}
  \proofstepwidestar[2,4]{\IntEmbed(\Succ(m))
    \leq\IntEmbed(\Succ(m))\cdot b}{%
    \FormulaRefAuto{PositiveIntegerEmbeddingFactorLowerBound}[thm]{5,7,2}}
  \proofstepwidestar[4]{\IntEmbed(\Succ(m))\in\mathbb Z}{%
    \FormulaRefAuto{NaturalIntegerEmbeddingFunction}[thm]{5}}
  \proofstepwidestar[2]{b\in\mathbb Z\land\IntZero<b}{%
    \FormulaRefAuto{PositiveIntegersDef}[def]{2}}
  \proofstepwidestar[2,4]{\IntEmbed(\Succ(m))\cdot b\in\mathbb Z}{%
    \FormulaRefAuto{IntegerMulClosureAxiom}[ax]{9,\rAEa{10}}}
  \proofstepwidestar[1,2,4]{a<\IntEmbed(\Succ(m))\cdot b}{%
    \FormulaRefAuto{IntegerStrictLeqTransitiveFromAxioms}[thm]{%
      1,9,11,6,8}}
  \proofstepwidestar[1,2,4]{\exists n\in\mathbb N\,
    a<\IntEmbed(n)\cdot b}{\rEI{\rAI{5,12}}}
  \proofstepwidestar[1,2]{\exists n\in\mathbb N\,
    a<\IntEmbed(n)\cdot b}{\rEE{3,4,13}}
\end{tabproofwide}

\chapter{Abschluss der Modellkonstruktion}

\section{Herleitung der zweiseitigen Induktion im Modell}

\FormulaThmDeltaK[Zweiseitige Induktion im Quotientenmodell]{%
  \begin{aligned}[t]
    &P(\IntZero)\dsep
    \forall z\in\mathbb Z\,
      (P(z)\rightarrow P(\IntSucc(z)))\dsep{}\\
    &\forall z\in\mathbb Z\,
      (P(z)\rightarrow P(\IntPred(z)))
    \vdash \forall z\in\mathbb Z\,P(z)
  \end{aligned}%
}{IntegerTwoSidedInductionModel}{%
  \DeltaRow{Einstelliges Prädikat}{P}%
  \DeltaRow{Mengen}{z\dsep n}%
}
\begin{tabproofsplitwide}
  \proofpartwideindR[Induktion auf dem nichtnegativen Strahl]{%
    \begin{aligned}[t]
      &P(\IntZero)\dsep
      \forall z\in\mathbb Z\,
      (P(z)\rightarrow P(\IntSucc(z)))\vdash{}\\
      &\forall n\in\mathbb N\,P(\IntEmbed(n))
    \end{aligned}
  }{%
    P(\IntZero)\dsep
    \forall z\in\mathbb Z\,
    (P(z)\rightarrow P(\IntSucc(z)))
    \vdash\forall n\in\mathbb N\,P(\IntEmbed(n))
  }[IntegerTwoSidedInductionPositiveRay]
    \proofstepwidestarbreakkeep[1]{P(\IntZero)}{\rA}
    \proofstepwidestarbreakkeep[2]{\forall z\in\mathbb Z\,
      (P(z)\rightarrow P(\IntSucc(z)))}{\rA}
    \proofstepwidestarbreakkeep[]{\IntZero=\IntEmbed(0)}{%
      \FormulaRefAuto{IntegerZeroDef}}
    \proofstepwidestarbreakkeep[1]{P(\IntEmbed(0))}{\rIE{3,1}}

    \proofstepwidestarbreakkeep[5]{n\in\mathbb N}{\rA}
    \proofstepwidestarbreakkeep[6]{P(\IntEmbed(n))}{\rA}
    \proofstepwidestarbreakkeep[5]{\IntEmbed(n)\in\mathbb Z}{%
      \FormulaRefAuto{F\colon A\to B\dsep x\in A\vdash F(x)\in B}{%
        \FormulaRefAuto{NaturalIntegerEmbeddingFunction},5}}
    \proofstepwidestarbreakkeep[2,5]{%
      P(\IntEmbed(n))\rightarrow
      P(\IntSucc(\IntEmbed(n)))}{\rRE{\rUE{2},7}}
    \proofstepwidestarbreakkeep[2,5,6]{P(\IntSucc(\IntEmbed(n)))}{\rRE{8,6}}
    \proofstepwidestarbreakkeep[5]{%
      \IntSucc(\IntEmbed(n))=\IntEmbed(\Succ(n))}{%
      \FormulaRefAuto{IntegerSuccessorOnEmbedding}{5}}
    \proofstepwidestarbreakkeep[2,5,6]{P(\IntEmbed(\Succ(n)))}{\rIE{10,9}}
    \proofstepwidestarbreakkeep[2,5]{%
      P(\IntEmbed(n))\rightarrow P(\IntEmbed(\Succ(n)))}{\rRI{6,11}}
    \proofstepwidestarbreakkeep[2]{\forall n\in\mathbb N\,
      (P(\IntEmbed(n))\rightarrow P(\IntEmbed(\Succ(n))))}{%
      \rUI{\rRI{5,12}}}
    \proofstepwidestarbreakkeep[1,2]{\forall n\in\mathbb N\,
      P(\IntEmbed(n))}{%
      \FormulaRefAuto{PeanoInductionAxiom}{4,13}}
  \closeproofpartwideindR

  \proofpartwideindR[Induktion auf dem nichtpositiven Strahl]{%
    \begin{aligned}[t]
      &P(\IntZero)\dsep
      \forall z\in\mathbb Z\,
      (P(z)\rightarrow P(\IntPred(z)))\vdash{}\\
      &\forall n\in\mathbb N\,P(-\IntEmbed(n))
    \end{aligned}
  }{%
    P(\IntZero)\dsep
    \forall z\in\mathbb Z\,
    (P(z)\rightarrow P(\IntPred(z)))
    \vdash\forall n\in\mathbb N\,P(-\IntEmbed(n))
  }[IntegerTwoSidedInductionNegativeRay]
    \proofstepwidestarbreakkeep[1]{P(\IntZero)}{\rA}
    \proofstepwidestarbreakkeep[2]{\forall z\in\mathbb Z\,
      (P(z)\rightarrow P(\IntPred(z)))}{\rA}
    \proofstepwidestarbreakkeep[]{\IntEmbed(0)=\IntZero}{%
      \FormulaRefAuto{a=b\vdash b=a}{\FormulaRefAuto{IntegerZeroDef}}}
    \proofstepwidestarbreakkeep[]{-\IntZero=\IntZero}{%
      \FormulaRefAuto{IntegerNegativeZero}[thm]}
    \proofstepwidestarbreakkeep[1]{P(-\IntEmbed(0))}{\rIE{3,4,1}}

    \proofstepwidestarbreakkeep[6]{n\in\mathbb N}{\rA}
    \proofstepwidestarbreakkeep[7]{P(-\IntEmbed(n))}{\rA}
    \proofstepwidestarbreakkeep[6]{-\IntEmbed(n)\in\mathbb Z}{%
      \FormulaRefAuto{IntegerNegClosureAxiom}[ax]{%
        \FormulaRefAuto{F\colon A\to B\dsep x\in A\vdash F(x)\in B}{%
          \FormulaRefAuto{NaturalIntegerEmbeddingFunction},6}}}
    \proofstepwidestarbreakkeep[2,6]{%
      P(-\IntEmbed(n))\rightarrow
      P(\IntPred(-\IntEmbed(n)))}{\rRE{\rUE{2},8}}
    \proofstepwidestarbreakkeep[2,6,7]{P(\IntPred(-\IntEmbed(n)))}{\rRE{9,7}}
    \proofstepwidestarbreakkeep[6]{%
      \IntPred(-\IntEmbed(n))=-\IntEmbed(\Succ(n))}{%
      \FormulaRefAuto{IntegerPredecessorOnNegativeEmbedding}{6}}
    \proofstepwidestarbreakkeep[2,6,7]{P(-\IntEmbed(\Succ(n)))}{\rIE{11,10}}
    \proofstepwidestarbreakkeep[2,6]{%
      P(-\IntEmbed(n))\rightarrow P(-\IntEmbed(\Succ(n)))}{\rRI{7,12}}
    \proofstepwidestarbreakkeep[2]{\forall n\in\mathbb N\,
      (P(-\IntEmbed(n))\rightarrow P(-\IntEmbed(\Succ(n))))}{%
      \rUI{\rRI{6,13}}}
    \proofstepwidestarbreakkeep[1,2]{\forall n\in\mathbb N\,
      P(-\IntEmbed(n))}{%
      \FormulaRefAuto{PeanoInductionAxiom}{5,14}}
  \closeproofpartwideindR

  \proofpartwide{Schluss mittels Normalform}
    \proofstepwidestarbreakkeep[1]{P(\IntZero)}{\rA}
    \proofstepwidestarbreakkeep[2]{\forall z\in\mathbb Z\,
      (P(z)\rightarrow P(\IntSucc(z)))}{\rA}
    \proofstepwidestarbreakkeep[3]{\forall z\in\mathbb Z\,
      (P(z)\rightarrow P(\IntPred(z)))}{\rA}
    \proofstepwidestarbreakkeep[1,2]{\forall n\in\mathbb N\,
      P(\IntEmbed(n))}{%
      \FormulaRefAuto{IntegerTwoSidedInductionPositiveRay}{1,2}}
    \proofstepwidestarbreakkeep[1,3]{\forall n\in\mathbb N\,
      P(-\IntEmbed(n))}{%
      \FormulaRefAuto{IntegerTwoSidedInductionNegativeRay}{1,3}}
    \proofstepwidestarbreakkeep[6]{z\in\mathbb Z}{\rA}
    \proofstepwidestarbreakkeep[6]{\exists n\in\mathbb N\,
      (z=\IntEmbed(n)\lor z=-\IntEmbed(n))}{%
      \FormulaRefAuto{IntegerSignedNormalForm}{6}}
    \proofstepwidestarbreakkeep[8]{n\in\mathbb N\land
      (z=\IntEmbed(n)\lor z=-\IntEmbed(n))}{\rA}
    \proofstepwidestarbreakkeep[4,8]{P(\IntEmbed(n))}{\rRE{\rUE{4},\rAEa{8}}}
    \proofstepwidestarbreakkeep[5,8]{P(-\IntEmbed(n))}{\rRE{\rUE{5},\rAEa{8}}}
    \proofstepwidestarbreakkeep[11]{z=\IntEmbed(n)}{\rA}
    \proofstepwidestarbreakkeep[4,8,11]{P(z)}{\rIE{11,9}}
    \proofstepwidestarbreakkeep[13]{z=-\IntEmbed(n)}{\rA}
    \proofstepwidestarbreakkeep[5,8,13]{P(z)}{\rIE{13,10}}
    \proofstepwidestarbreakkeep[4,5,8]{P(z)}{%
      \rOE{\rAEb{8},11,12,13,14}}
    \proofstepwidestarbreakkeep[4,5,6]{P(z)}{\rEE{7,8,15}}
    \proofstepwidestarbreakkeep[1,2,3]{\forall z\in\mathbb Z\,P(z)}{%
      \rUI{\rRI{6,16}}}
  \closeproofpartwide
\end{tabproofsplitwide}

\section{Abschluss des Modellnachweises}

Die vorangegangenen Herleitungen erfassen nun Ringoperationen, Ordnung,
Schrittstruktur und zweiseitige Induktion. Der folgende Satz bündelt diese
Ergebnisse und beendet die konkrete Konstruktion. Er ist zugleich der letzte
Satz dieses Bandes, in dem der konstruierte Träger selbst auftritt.

\FormulaThmDeltaK[Das Quotientenmodell erfüllt die Ganzzahlaxiome]{%
  \begin{aligned}[t]
    &\mathbb Z=\QuotientSet{D_{\mathbb Z}}{\IntPairRel}\ \land{}\\
    &\IntegerPeano{%
      \mathbb Z,\IntZero,\IntOne,-,+,\cdot,
      \IntSucc,\IntPred,\leq_{\mathbb Z}}
  \end{aligned}%
}{IntQuotientModelsIntegerPeano}{%
  \DeltaRow{Mengen}{z\dsep D_{\mathbb Z}}%
  \DeltaRow{Funktionen}{\IntSucc\dsep\IntPred}%
  \DeltaRow{Relationsoperatoren}{\IntPairRel\dsep\leq_{\mathbb Z}}%
}
\begin{tabproofsplitwide}
  \proofpartwideindR[Träger und Abbildungen]{%
    \begin{aligned}[t]
      &\IntZero\in\mathbb Z\ \land\
      \IntSucc\colon\mathbb Z\to\mathbb Z\ \land{}\\
      &\IntPred\colon\mathbb Z\to\mathbb Z
    \end{aligned}
  }{%
    \IntZero\in\mathbb Z\land
    \IntSucc\colon\mathbb Z\to\mathbb Z\land
    \IntPred\colon\mathbb Z\to\mathbb Z
  }[IntegerPeanoModelCarrierMaps]
    \proofstepwidestar[]{\IntZero\in\mathbb Z}{%
      \FormulaRefAuto{IntZeroCarrier}}
    \proofstepwidestar[]{\IntSucc\colon\mathbb Z\to\mathbb Z}{%
      \FormulaRefAuto{IntegerSuccessorFunction}}
    \proofstepwidestar[]{\IntPred\colon\mathbb Z\to\mathbb Z}{%
      \FormulaRefAuto{IntegerPredecessorFunction}}
    \proofstepwidestar[]{\IntZero\in\mathbb Z\land
      \IntSucc\colon\mathbb Z\to\mathbb Z\land
      \IntPred\colon\mathbb Z\to\mathbb Z}{\rAI{1,\rAI{2,3}}}
  \closeproofpartwideindR

  \proofpartwideindR[Gegenseitige Inversen]{%
    \begin{aligned}[t]
      &\forall z\in\mathbb Z\,
        \IntPred(\IntSucc(z))=z\ \land{}\\
      &\forall z\in\mathbb Z\,
        \IntSucc(\IntPred(z))=z
    \end{aligned}
  }{%
    \forall z\in\mathbb Z\,\IntPred(\IntSucc(z))=z
    \land
    \forall z\in\mathbb Z\,\IntSucc(\IntPred(z))=z
  }[IntegerPeanoModelInverses]
    \proofstepwidestar[]{\forall z\in\mathbb Z\,
      \IntPred(\IntSucc(z))=z}{%
      \FormulaRefAuto{IntegerPredAfterSucc}}
    \proofstepwidestar[]{\forall z\in\mathbb Z\,
      \IntSucc(\IntPred(z))=z}{%
      \FormulaRefAuto{IntegerSuccAfterPred}}
    \proofstepwidestar[]{%
      \forall z\in\mathbb Z\,\IntPred(\IntSucc(z))=z
      \land
      \forall z\in\mathbb Z\,\IntSucc(\IntPred(z))=z}{\rAI{1,2}}
  \closeproofpartwideindR

  \proofpartwideindR[Geordnete Nachfolgerstruktur]{%
    \TotOrd{\mathbb Z,\leq_{\mathbb Z}}
    \land\forall z\in\mathbb Z\,(z<\IntSucc(z))%
  }{%
    \TotOrd{\mathbb Z,\leq_{\mathbb Z}}
    \land\forall z\in\mathbb Z\,(z<\IntSucc(z))%
  }[IntegerPeanoModelOrder]
    \proofstepwidestar[]{\TotOrd{\mathbb Z,\leq_{\mathbb Z}}}{%
      \FormulaRefAuto{IntegerOrderTotal}}
    \proofstepwidestar[]{\forall z\in\mathbb Z\,
      (z<\IntSucc(z))}{\FormulaRefAuto{IntegerSuccessorPositive}}
    \proofstepwidestar[]{\TotOrd{\mathbb Z,\leq_{\mathbb Z}}
      \land\forall z\in\mathbb Z\,(z<\IntSucc(z))}{\rAI{1,2}}
  \closeproofpartwideindR

  \proofpartwideindR[Induktionsschema]{%
    \begin{aligned}[t]
      &P(\IntZero)\dsep
      \forall z\in\mathbb Z\,
      (P(z)\to P(\IntSucc(z)))\dsep{}\\
      &\forall z\in\mathbb Z\,
      (P(z)\to P(\IntPred(z)))
      \vdash\forall z\in\mathbb Z\,P(z)
    \end{aligned}
  }{%
    P(\IntZero)\dsep
    \forall z\in\mathbb Z\,(P(z)\to P(\IntSucc(z)))\dsep
    \forall z\in\mathbb Z\,(P(z)\to P(\IntPred(z)))
    \vdash\forall z\in\mathbb Z\,P(z)
  }[IntegerPeanoModelInduction]
    \proofstepwidestar[1]{P(\IntZero)}{\rA}
    \proofstepwidestar[2]{\forall z\in\mathbb Z\,
      (P(z)\to P(\IntSucc(z)))}{\rA}
    \proofstepwidestar[3]{\forall z\in\mathbb Z\,
      (P(z)\to P(\IntPred(z)))}{\rA}
    \proofstepwidestar[1,2,3]{\forall z\in\mathbb Z\,P(z)}{%
      \FormulaRefAuto{IntegerTwoSidedInductionModel}{1,2,3}}
  \closeproofpartwideindR

  \proofpartwide{Schluss}
    \proofstepwidestarbreakkeep[]{%
      \mathbb Z=\QuotientSet{D_{\mathbb Z}}{\IntPairRel}}{%
      \FormulaRefAuto{IntegersAsQuotient}}
    \proofstepwidestarbreakkeep[]{%
      \IntegerArithmetic{%
        \mathbb Z,\IntZero,\IntOne,-,+,\cdot,\leq_{\mathbb Z}}}{%
      \FormulaRefAuto{IntegerOrderedArithmeticAxioms}}
    \proofstepwidestarbreakkeep[]{\IntZero\in\mathbb Z}{%
      \rAEa{\FormulaRefAuto{IntegerPeanoModelCarrierMaps}}}
    \proofstepwidestarbreakkeep[]{\IntSucc\colon\mathbb Z\to\mathbb Z}{%
      \rAEa{\rAEb{\FormulaRefAuto{IntegerPeanoModelCarrierMaps}}}}
    \proofstepwidestarbreakkeep[]{\IntPred\colon\mathbb Z\to\mathbb Z}{%
      \rAEb{\rAEb{\FormulaRefAuto{IntegerPeanoModelCarrierMaps}}}}
    \proofstepwidestarbreakkeep[6]{z\in\mathbb Z}{\rA}
    \proofstepwidestarbreakkeep[6]{\IntPred(\IntSucc(z))=z}{%
      \rRE{6,\rAEa{\FormulaRefAuto{IntegerPeanoModelInverses}}}}
    \proofstepwidestarbreakkeep[6]{\IntSucc(\IntPred(z))=z}{%
      \rRE{6,\rAEb{\FormulaRefAuto{IntegerPeanoModelInverses}}}}
    \proofstepwidestarbreakkeep[]{\TotOrd{\mathbb Z,\leq_{\mathbb Z}}}{%
      \rAEa{\FormulaRefAuto{IntegerPeanoModelOrder}}}
    \proofstepwidestarbreakkeep[6]{z<\IntSucc(z)}{%
      \rRE{6,\rAEb{\FormulaRefAuto{IntegerPeanoModelOrder}}}}
    \proofstepwidestarbreakkeep[11]{P(\IntZero)}{\rA}
    \proofstepwidestarbreakkeep[12]{%
      \forall z\in\mathbb Z(P(z)\to P(\IntSucc(z)))}{\rA}
    \proofstepwidestarbreakkeep[13]{%
      \forall z\in\mathbb Z(P(z)\to P(\IntPred(z)))}{\rA}
    \proofstepwidestarbreakkeep[11,12,13]{\forall z\in\mathbb Z\,P(z)}{%
      \FormulaRefAuto{IntegerPeanoModelInduction}[thm]{11,12,13}}
    \proofstepwidestarbreakkeep[]{%
      \IntegerPeano{%
        \mathbb Z,\IntZero,\IntOne,-,+,\cdot,
        \IntSucc,\IntPred,\leq_{\mathbb Z}}}{%
      \FormulaRefAuto{IntegerPeanoAxioms}{2,3,4,5,7,8,10,14}}
    \proofstepwidestarbreakkeep[]{%
      \mathbb Z=\QuotientSet{D_{\mathbb Z}}{\IntPairRel}
      \land
      \IntegerPeano{%
        \mathbb Z,\IntZero,\IntOne,-,+,\cdot,
        \IntSucc,\IntPred,\leq_{\mathbb Z}}}{%
      \rAI{1,15}}
  \closeproofpartwide
\end{tabproofsplitwide}

\chapter{Axiomatische Beschreibung der ganzen Zahlen}

\section{Abstraktionsschnitt}

Der Modellnachweis ist abgeschlossen. Ab hier bezeichnet \(\mathbb Z\) eine
Struktur mit den unten aufgeführten Konstanten, Operationen, Schrittoperatoren
und der Ordnung. Ihre Eigenschaften werden ausschließlich durch die
registrierten Axiome und die daraus abgeleiteten Sätze bestimmt. Die konkrete
Realisierung hat von diesem Punkt an nur noch die Rolle eines
Existenznachweises und ist kein Bestandteil der weiteren Sprache.

Die Definitionen
\(\IntSucc(z)=z+\IntOne\) und
\(\IntPred(z)=z+(-\IntOne)\) verbinden die Schrittstruktur mit der
Arithmetik. Zusammen mit der geordneten Ringstruktur und der zweiseitigen
Induktion bildet dies die vollständige Schnittstelle für alle späteren
Verwendungen der ganzen Zahlen.

\section{Axiome der Schrittstruktur}

\FormulaAxiomAutoR[Null-Axiom der ganzen Zahlen]{%
  \IntZero\in\mathbb Z%
}{IntegerPeanoZero}
\par\noindent\emph{Modellnachweis vor dem Abstraktionsschnitt:}
\FormulaRefAuto{IntZeroCarrier}[thm].\par

\FormulaAxiomAutoR[Nachfolgerfunktion der ganzen Zahlen]{%
  \IntSucc\colon\mathbb Z\to\mathbb Z%
}{IntegerPeanoSuccFunction}
\par\noindent\emph{Modellnachweis vor dem Abstraktionsschnitt:}
\FormulaRefAuto{IntegerSuccessorFunction}[thm].\par

\FormulaAxiomAutoR[Vorgängerfunktion der ganzen Zahlen]{%
  \IntPred\colon\mathbb Z\to\mathbb Z%
}{IntegerPeanoPredFunction}
\par\noindent\emph{Modellnachweis vor dem Abstraktionsschnitt:}
\FormulaRefAuto{IntegerPredecessorFunction}[thm].\par

\FormulaAxiomAutoR[Vorgänger nach Nachfolger]{%
  z\in\mathbb Z\vdash\IntPred(\IntSucc(z))=z%
}{IntegerPeanoPredSucc}
\par\noindent\emph{Modellnachweis vor dem Abstraktionsschnitt:}
\FormulaRefAuto{IntegerPredAfterSucc}[thm].\par

\FormulaAxiomAutoR[Nachfolger nach Vorgänger]{%
  z\in\mathbb Z\vdash\IntSucc(\IntPred(z))=z%
}{IntegerPeanoSuccPred}
\par\noindent\emph{Modellnachweis vor dem Abstraktionsschnitt:}
\FormulaRefAuto{IntegerSuccAfterPred}[thm].\par

\FormulaAxiomAutoR[Totale Ordnung der ganzen Zahlen]{%
  \TotOrd{\mathbb Z,\leq_{\mathbb Z}}%
}{IntegerPeanoOrder}
\par\noindent\emph{Modellnachweis vor dem Abstraktionsschnitt:}
\FormulaRefAuto{IntegerOrderTotal}[thm].\par

\FormulaAxiomAutoR[Strikte Nachfolgerorientierung]{%
  z\in\mathbb Z\vdash z<\IntSucc(z)%
}{IntegerPeanoSuccPositive}
\par\noindent\emph{Modellnachweis vor dem Abstraktionsschnitt:}
\FormulaRefAuto{IntegerSuccessorPositive}[thm].\par

\FormulaAxiomAutoR[Zweiseitige Induktion der ganzen Zahlen]{%
  \begin{aligned}[t]
    &P(\IntZero)\dsep
    \forall z\in\mathbb Z\,
      (P(z)\to P(\IntSucc(z)))\dsep{}\\
    &\forall z\in\mathbb Z\,
      (P(z)\to P(\IntPred(z)))
    \vdash\forall z\in\mathbb Z\,P(z)
  \end{aligned}%
}{IntegerTwoSidedInductionAxiom}
\par\noindent\emph{Modellnachweis vor dem Abstraktionsschnitt:}
\FormulaRefAuto{IntegerTwoSidedInductionModel}[thm].\par

\FormulaAxiomAutoR[Zweiseitiges Induktionsaxiom]{%
  \begin{aligned}[t]
    &P(\IntZero)\dsep
    \forall z\in\mathbb Z\,
      (P(z)\to P(\IntSucc(z)))\dsep{}\\
    &\forall z\in\mathbb Z\,
      (P(z)\to P(\IntPred(z)))
    \vdash\forall z\in\mathbb Z\,P(z)
  \end{aligned}%
}{IntegerPeanoInduction}
\par\noindent\emph{Registrierte Kurzform:}
\FormulaRefAuto{IntegerTwoSidedInductionAxiom}[ax].\par

\section{Vollständige Axiomatik}

Die vollständige Schnittstelle verbindet die bereits registrierte geordnete
Ringarithmetik mit der Schrittstruktur. Anders als auf \(\mathbb N\) ist der
Nachfolger auf \(\mathbb Z\) bijektiv. Die totale Ordnung und die strikte
Nachfolgerorientierung verhindern zyklische Nachfolgermodelle; die
zweiseitige Induktion verbindet den positiven und den negativen Strahl.

Metalogisch ist dabei zwischen zwei Lesarten zu unterscheiden: Als
erststufiges Axiomenschema ist die Beschreibung nicht kategorisch und
besitzt nichtstandardmäßige Modelle. Unter voller zweitstufiger Lesart der
Induktion beschreibt sie dagegen die intendierte zweiseitige, von
\(\IntZero\) durch Nachfolger und Vorgänger erzeugte Kette.

\FormulaDefDeltaK[Erweiterte Peano-Axiome der ganzen Zahlen]{%
  \IntegerPeano{%
    \mathbb Z,\IntZero,\IntOne,-,+,\cdot,
    \IntSucc,\IntPred,\leq_{\mathbb Z}}%
}{IntegerPeanoAxioms}{%
  \DeltaRow{\textbf{Ring und Ordnung}}{}%
  \DeltaRow{Arithmetik}{%
    \IntegerArithmetic{%
      \mathbb Z,\IntZero,\IntOne,-,+,\cdot,\leq_{\mathbb Z}}}
    [\FormulaRefAuto{IntegerOrderedArithmeticAxioms}]%
  \DeltaRow{\textbf{Schrittstruktur}}{}%
  \DeltaRow{Null}{\IntZero\in\mathbb Z}
    [\FormulaRefAuto{IntegerPeanoZero}]%
  \DeltaRow{Nachfolger}{\IntSucc\colon\mathbb Z\to\mathbb Z}
    [\FormulaRefAuto{IntegerPeanoSuccFunction}]%
  \DeltaRow{Vorgänger}{\IntPred\colon\mathbb Z\to\mathbb Z}
    [\FormulaRefAuto{IntegerPeanoPredFunction}]%
  \DeltaRow{Links invers}
    {z\in\mathbb Z\vdash\IntPred(\IntSucc(z))=z}
    [\FormulaRefAuto{IntegerPeanoPredSucc}]%
  \DeltaRow{Rechts invers}
    {z\in\mathbb Z\vdash\IntSucc(\IntPred(z))=z}
    [\FormulaRefAuto{IntegerPeanoSuccPred}]%
  \DeltaRow{Orientierung}
    {z\in\mathbb Z\vdash z<\IntSucc(z)}
    [\FormulaRefAuto{IntegerSuccessorPositiveAxiom}]%
  \DeltaRow{\textbf{Induktionsschema}}{}%
  \DeltaRow{Induktion}{%
    \begin{aligned}[t]
      &P(\IntZero)\dsep{}\\[-2pt]
      &\forall z\in\mathbb Z\,(P(z)\to P(\IntSucc(z)))\dsep{}\\[-2pt]
      &\forall z\in\mathbb Z\,(P(z)\to P(\IntPred(z)))\dsep{}\\[-2pt]
      &\vdash\forall z\in\mathbb Z\,P(z)
    \end{aligned}}
    [\FormulaRefAuto{IntegerTwoSidedInductionAxiom}]%
  \DeltaRow{\textbf{Signatur}}{}%
  \DeltaRow{Ganze Zahlen}{\mathbb Z}%
  \DeltaRow{Konstanten}{\IntZero\dsep\IntOne}%
  \DeltaRow{Operationen}{-\dsep+\dsep\cdot}%
  \DeltaRow{Schrittoperatoren}{\IntSucc\dsep\IntPred}%
  \DeltaRow{Ordnung}{\leq_{\mathbb Z}}%
}

\section*{Axiomatische Verwendung}

Von nun an ist \FormulaRefAuto{IntegerPeanoAxioms}[def] die alleinige
Grundlage für Aussagen über ganze Zahlen. Verwendet werden nur diese Axiome
und ihre bereits hergeleiteten Folgerungen. Die konkrete Realisierung ist für
solche Argumente weder Voraussetzung noch zulässiger Beweisschritt.

\end{document}
